\chapter{\texorpdfstring{Heterotic Strings}{11 Heterotic Strings}} \label{cha:11}

\section{World-sheet Supersymmetry} \label{sec:11.1}

In the previous chapter, we saw the method of expanding the world-sheet constraint algebra by adding supercurrents $T_{F}(z)$ and $\tilde{T}_{F}(\bar{z})$. Now we look at how far this idea can go—specifically, seeking a set of holomorphic or antiholomorphic currents such that their Laurent coefficients form a closed algebra.

We first specify the difference between global symmetries and constraints. Global symmetries on the world-sheet are like global symmetries in spacetime; they imply relationships between masses and amplitudes. However, we usually also pick out a portion of the symmetries as \emph{constraints} to be attached to the system, meaning that physical states must be annihilated by these constraints. In the bosonic string, spacetime Poincaré symmetry is a global symmetry of the world-sheet theory, while conformal symmetry is a constraint. Our current interest is in constraint algebras. In fact, the possible sets we find are very small, but later we will encounter some additive algebras that appear as global symmetries.

First, it should be noted that the collection of sets capable of acting as world-sheet symmetry algebras is very large. For example, in the bosonic string, any product of factors $\partial^{n}X^{\mu}$ is a holomorphic current. In most cases, the OPE of such currents will generate an infinite number of new currents. Even if restricted to a finite number of currents, infinite possibilities still remain.

We first concentrate on holomorphic currents. In a unitary CFT, a primary operator can be holomorphic only if its weight is $(h,0)$ with $h>0$ (holomorphicity requires $\tilde{h}=0$, and unitarity gives $h>0$). Although the complete world-sheet theory lacks a positive-definite norm due to ghost fields and timelike oscillators, the spatial part is indeed unitary, so it resides in a unitary representation of the symmetry. Since $\tilde{h}=0$, the spin of the current is also equal to $h$. (Note that $h+\tilde{h}$ determines the scaling behavior of the current, while $h-\tilde{h}$ determines the rotation behavior, because $z\bar{z}=r^{2}$ and $z/\bar{z}=\me^{2\mi\varphi}$, where $z=r\me^{\mi\varphi}$). We assume the currents are Hermitian and now examine the following possibilities:

\emph{Spin} $h\geq 2$: For algebras of currents with spin $>2$, we usually call them $W$-\emph{algebras}. Most are known, including several classes of infinite algebras, but there is no complete classification yet. There have been attempts to use some of these algebras as constraint algebras. One difficulty is that the commutators of the generators are generally non-linear functions of the generators, which makes the construction of the BRST operator very non-trivial. A few constructed examples were found after gauge fixing to be just special cases of the bosonic string. Another difficulty is that the geometric interpretation is unclear. So we will focus our attention on algebras with $h\leq 2$. Additionally, a CFT can possess multiple $(2,0)$ currents as global symmetries. The bosonic string has at least 27, namely the energy-momentum tensor of the ghost fields and the energy-momentum tensors of the 26 $X^{\mu}$ fields. However, only their sum has a geometric interpretation as conformal invariance; therefore, we will assume there is only one $(2,0)$ constraint current, namely the total energy-momentum tensor.

\emph{Spin $h$ is not an integer multiple of $\frac{1}{2}$}: For a current $j$ with spin $h$,
\begin{equation}
	j(z)j(0)\sim z^{-2h} \label{11.1.1}
\end{equation}
If $2h$ is not an integer, this will be multi-valued. Although many CFTs are already known to possess such currents, difficulties arise from non-locality when we try to use such currents as constraints. Attempts to construct such \emph{fractional strings} have yielded partial results, but whether such theories exist remains unclear. So we restrict $h$ to be an integer multiple of $\frac{1}{2}$.

With these constraints, the possible algebras are greatly restricted; the spin can only be $0, \frac{1}{2}, 1, \frac{3}{2}, 2$. The algebras allowed by the Jacobi identity are listed in Table \ref{tab:11.1}. The first two are obviously the conformal algebra and the $N=1$ superconformal algebra we have studied. The three $N=4$ algebras are related to each other. Among them, the second algebra is a special case of the first algebra, where the $U(1)$ current of the first algebra is the gradient of some scalar. The third is a subalgebra of the second.

\begin{table}[ht]
	\caption{World-sheet superconformal algebras. The table lists the number of currents of each spin, the total ghost central charge, the global symmetries generated by spin-1 currents, and the transformations of supercharges under these symmetries.}
	\label{tab:11.1}%
	\centering
	\begin{tabular}[c]{cccccll}
		\hline\hline
		$\vphantom{\Big[}n_{3/2}=N$ & $n_{1}$ & $n_{1/2}$ & $n_{0}$ & $c^{\text{g}}$ & symmetry & $T_{F}\:$rep. \\
		\hline
		0 & 0 & 0 & 0 & $-26$ & &  \\
		1 & 0 & 0 & 0 & $-15$ & &  \\
		2 & 1 & 0 & 0 & $-6$ & $U(1)$ & $\pm1$  \\
		3 & 3 & 1 & 0 & 0 & $SU(2)$ & $\mathbf{3}$  \\
		4 & 7 & 4 & 0 & 0 & $SU(2)\times SU(2)\times U(1)$ & $(\mathbf{2},\mathbf{2},0)$  \\
		4 & 6 & 4 & 1 & 0 & $SU(2)\times SU(2)$ & $(\mathbf{2},\mathbf{2})$  \\
		4 & 3 & 0 & 0 & 12 & $SU(2)$ & $\mathbf{2}$  \\
		\hline\hline
	\end{tabular}
\end{table}

The ghost central charge is determined by the number of currents of each spin; the ghost central charge associated with a current of spin $h$ is
\begin{subequations}
	\begin{align}
		c_{h}&=(-1)^{2h+1}[3(2h-1)^{2}-1] \: , \label{11.1.2a} \\
		c_{2}&=-26\:,\qquad c_{3/2}=+11\:,\qquad c_{1}=-2\:,\qquad c_{1/2}=-1\:, \qquad c_{0}=-2\:. \label{11.1.2b}
	\end{align} \label{11.1.2}
\end{subequations}
The sign $(-1)^{2h+1}$ is the result of taking the statistics of the ghost fields into account, anticommuting for integer spin and commuting for half-integer spin. Since the matter field central charge $c^{\text{m}}$ is $-c^{\text{g}}$, only one new algebra has a critical dimension, namely $N=2$.
\begin{tcolorbox}[breakable]
	\begin{remark}
		Gauge fixing a constraint current of weight $h$ introduces a ghost pair: an antighost $b$ of weight $\lambda=h$ (it has the same indices as the current) and a ghost $c$ of weight $1-h$. The pair has the opposite statistics to the current, so it is Grassmann odd for integer $h$ and Grassmann even for half-integer $h$. The central charges of the two kinds of pair are \eqref{2.5.12} and its commuting analogue,
		\begin{equation*}
			c_{bc}=-3(2\lambda-1)^{2}+1\:,\qquad c_{\beta\gamma}=+3(2\lambda-1)^{2}-1\:,
		\end{equation*}
		so with $\lambda=h$,
		\begin{align*}
			c_{h}=\begin{cases} -[3(2h-1)^{2}-1] & h\in\mathds{Z} \\ +[3(2h-1)^{2}-1] & h\in\mathds{Z}+\tfrac{1}{2}\end{cases}
			\;=(-1)^{2h+1}\bigl[3(2h-1)^{2}-1\bigr]\:.
		\end{align*}
		Evaluating,
		\begin{align*}
			c_{2}&=-(27-1)=-26\:,& c_{3/2}&=+(12-1)=11\:,& c_{1}&=-(3-1)=-2\:,\\
			c_{1/2}&=+(0-1)=-1\:,& c_{0}&=-(3-1)=-2\:.
		\end{align*}
		The ghost central charges in Table \ref{tab:11.1} follow from $c^{\text{g}}=c_{2}+n_{3/2}c_{3/2}+n_{1}c_{1}+n_{1/2}c_{1/2}+n_{0}c_{0}$:
		\begin{align*}
			N=2&: \quad -26+2(11)+1(-2)=-6\:, \\
			N=3&: \quad -26+3(11)+3(-2)+1(-1)=0\:, \\
			N=4\ (7,4,0)&: \quad -26+4(11)+7(-2)+4(-1)=0\:, \\
			N=4\ (6,4,1)&: \quad -26+4(11)+6(-2)+4(-1)+1(-2)=0\:, \\
			N=4\ (3,0,0)&: \quad -26+4(11)+3(-2)=12\:.
		\end{align*}
		A critical dimension requires $c^{\text{m}}=-c^{\text{g}}>0$, which among the new algebras happens only for $N=2$, with $c^{\text{m}}=6$.
	\end{remark}
\end{tcolorbox}

For $N=2$, it is convenient to combine two real supercurrents into one complex supercurrent
\begin{equation}
	T_{F}^{\pm} = 2^{-1/2}(T_{F1}\pm \mi T_{F2}) \:. \label{11.1.3}
\end{equation}
In this way, the $N=2$ algebra in operator product form is
\begin{subequations}
	\begin{align}
		T_{B}(z)T_{F}^{\pm}(0)&\sim \frac{3}{2z^{2}} T_{F}^{\pm}(0) + \frac{1}{z}\partial T_{F}^{\pm}(0)\:,\label{11.1.4a} \\
		T_{B}(z) j(0) &\sim \frac{1}{z^{2}} j(0) +\frac{1}{z}\partial j(0) \:, \label{11.1.4b} \\
		T_{F}^{+}(z)T_{F}^{-}(0)&\sim \frac{2c}{3z^{3}}+\frac{2}{z^{2}}j(0)+\frac{2}{z}T_{B}(0)+\frac{1}{z}\partial j(0) \:,\label{11.1.4c} \\
		T^{+}_{F}(z)T^{+}_{F}(0)& \sim T^{-}_{F}(z)T^{-}_{F}(0) \sim 0 \:, \label{11.1.4d} \\
		j(z)T_{F}^{\pm}(0) &\sim \pm \frac{1}{z} T_{F}^{\pm}(0) \:, \label{11.1.4e} \\
		j(z)j(0) &\sim \frac{c}{3z^{2}}\:. \label{11.1.4f}
	\end{align}
\end{subequations}
This implies that $T_{F}^{\pm}$ and $j$ are primary fields, and $T_{F}^{\pm}$ possesses charges $\pm1$ under the $U(1)$ generated by $j$. According to the Jacobi identity, the constant $c$ in $T_{F}^{+}T_{F}^{-}$ and $jj$ must be the central charge.

The minimal linear representation of the $N=2$ algebra possesses two real scalar fields and two fermions; we combine them into a complex scalar $Z$ and a complex fermion $\psi$. The action is
\begin{equation}
	S=\frac{1}{2\pi} \int \dif^{2}z\Bigl( \partial \overline{Z}\bar{\partial}Z 
	+\overline{\psi}\bar{\partial}\psi + \tilde{\overline{\psi}}\partial\tilde{\psi}\Bigr)\:. \label{11.1.5}
\end{equation}
The currents are
\begin{subequations}
	\begin{align}
		T_{B} &= -\partial \overline{Z}\partial Z -\frac{1}{2}(\overline{\psi}\partial\psi+\psi\partial\overline{\psi})\:,\qquad j=-\overline{\psi}\psi \:,\label{11.1.6a} \\
		T_{F}^{+} &= 2^{1/2} \mi\psi \partial\overline{Z} \:, \qquad 
		T_{F}^{-} = 2^{1/2} \mi\overline{\psi} \partial Z\:. \label{11.1.6b}
	\end{align} \label{11.1.6}
\end{subequations}
Simultaneously, antiholomorphic currents also exist, so this $Z\psi\tilde{\psi}$ CFT possesses $(2,2)$ superconformal symmetry.
\begin{tcolorbox}[breakable]
	\begin{remark}
		We check that the currents (\ref{11.1.6}) realize the algebra (\ref{11.1.4a})--(\ref{11.1.4f}) with $c=3$. First the propagators. Write $Z=2^{-1/2}(X^{1}+\mi X^{2})$ and $\psi=2^{-1/2}(\psi^{1}+\mi\psi^{2})$. Then
		\begin{align*}
			\partial\overline{Z}\bar{\partial}Z &\eqs{1} \tfrac{1}{2}\bigl(\partial X^{1}\bar{\partial}X^{1}+\partial X^{2}\bar{\partial}X^{2}\bigr)+\tfrac{\mi}{2}\bigl(\partial X^{1}\bar{\partial}X^{2}-\partial X^{2}\bar{\partial}X^{1}\bigr)\:, \\
			\overline{\psi}\bar{\partial}\psi &\eqs{2} \tfrac{1}{2}\bigl(\psi^{1}\bar{\partial}\psi^{1}+\psi^{2}\bar{\partial}\psi^{2}\bigr)+\tfrac{\mi}{2}\,\bar{\partial}\bigl(\psi^{1}\psi^{2}\bigr)\:.
		\end{align*}
		At (1) the imaginary part is $\tfrac{\mi}{2}[\partial(X^{1}\bar{\partial}X^{2})-\bar{\partial}(X^{1}\partial X^{2})]$, a total derivative. At (2) the cross terms $\psi^{1}\bar{\partial}\psi^{2}-\psi^{2}\bar{\partial}\psi^{1}=\psi^{1}\bar{\partial}\psi^{2}+(\bar{\partial}\psi^{1})\psi^{2}$ also form a total derivative. So (\ref{11.1.5}) is $\frac{1}{4\pi}\int\dif^{2}z\,(\partial X^{i}\bar{\partial}X^{i}+\psi^{i}\bar{\partial}\psi^{i}+\cdots)$, the usual free action with $\alpha'=2$, and $X^{i}X^{j}\sim-\delta^{ij}\ln\lvert z\rvert^{2}$, $\psi^{i}\psi^{j}\sim\delta^{ij}/z$. Recombining,
		\begin{equation*}
			\partial\overline{Z}(z)\,\partial Z(0)\sim-\frac{1}{z^{2}}\:,\quad \partial Z\,\partial Z\sim 0\:,\qquad
			\psi(z)\overline{\psi}(0)\sim\overline{\psi}(z)\psi(0)\sim\frac{1}{z}\:,\quad \psi\psi\sim\overline{\psi}\,\overline{\psi}\sim0\:.
		\end{equation*}
		\textbf{The $jj$ and $jT_F$ OPEs.} With $j=-\overline{\psi}\psi$,
		\begin{align*}
			j(z)j(0)&=\overline{\psi}\psi(z)\,\overline{\psi}\psi(0)
			\eqs{1}\frac{1}{z^{2}}+\frac{1}{z}\Bigl[:\!\overline{\psi}(z)\psi(0)\!:+:\!\psi(z)\overline{\psi}(0)\!:\Bigr]+\cdots
			\eqs{2}\frac{1}{z^{2}}+O(z^{0})\:,\\
			j(z)\psi(0)&=-\overline{\psi}\psi(z)\,\psi(0)\eqs{3}+\overline{\psi}(z)\psi(0)\,\psi(z)\sim\frac{1}{z}\psi(0)\:,\qquad
			j(z)\overline{\psi}(0)\sim-\frac{1}{z}\overline{\psi}(0)\:.
		\end{align*}
		At (1) the double contraction pairs $\overline{\psi}(z)$ with $\psi(0)$ and $\psi(z)$ with $\overline{\psi}(0)$; the pairs are nested, so there is no sign. At (2) the two single-contraction terms add to $:\!\overline{\psi}\psi+\psi\overline{\psi}\!:=0$ at $z\to0$, so there is no $1/z$ term, as expected for an abelian current. Thus $c/3=1$ and $c=3$. At (3) $\psi(0)$ was moved past $\psi(z)$, which costs a sign, to sit next to $\overline{\psi}(z)$. So $\psi$ has charge $+1$ and $\overline{\psi}$ charge $-1$. The $Z$ fields are neutral, so $T_{F}^{+}\propto\psi\partial\overline{Z}$ has charge $+1$ and $T_{F}^{-}$ has charge $-1$, which is (\ref{11.1.4e}).

		\textbf{The $T_F^+T_F^-$ OPE.} Bosons commute with fermions, so
		\begin{align*}
			T_{F}^{+}(z)T_{F}^{-}(0)&=-2\,\psi(z)\overline{\psi}(0)\;\partial\overline{Z}(z)\partial Z(0)\\
			&\eqs{1}-2\Bigl(\frac{1}{z}\Bigr)\Bigl(-\frac{1}{z^{2}}\Bigr)-\frac{2}{z}:\!\partial\overline{Z}\partial Z\!:(0)+\frac{2}{z^{2}}:\!\psi(z)\overline{\psi}(0)\!:\\
			&\eqs{2}\frac{2}{z^{3}}+\frac{2}{z^{2}}:\!\psi\overline{\psi}\!:(0)+\frac{1}{z}\Bigl[2:\!\partial\psi\,\overline{\psi}\!:-2:\!\partial\overline{Z}\partial Z\!:\Bigr](0)\\
			&\eqs{3}\frac{2c}{3z^{3}}+\frac{2}{z^{2}}j(0)+\frac{1}{z}\bigl[2T_{B}(0)+\partial j(0)\bigr]\:.
		\end{align*}
		At (1) the three terms are the double contraction, the fermion contraction alone, and the boson contraction alone. At (2) $\psi(z)$ was Taylor expanded about $0$. At (3) $:\!\psi\overline{\psi}\!:=-\overline{\psi}\psi=j$ and $c=3$ were used. The $1/z$ term needs the identity
		\begin{align*}
			2T_{B}+\partial j&=-2\partial\overline{Z}\partial Z-\overline{\psi}\partial\psi-\psi\partial\overline{\psi}-\partial\overline{\psi}\,\psi-\overline{\psi}\partial\psi\\
			&\eqs{4}-2\partial\overline{Z}\partial Z-2\overline{\psi}\partial\psi
			=-2\partial\overline{Z}\partial Z+2\partial\psi\,\overline{\psi}\:,
		\end{align*}
		where at (4) $\psi\partial\overline{\psi}=-\partial\overline{\psi}\,\psi$ cancels the third and fourth terms. Finally $T_{F}^{+}T_{F}^{+}\propto\psi\psi\,\partial\overline{Z}\partial\overline{Z}\sim0$ because neither factor has a contraction, which is (\ref{11.1.4d}).

		\textbf{Primarity.} $\psi$ has weight $\frac12$ and $\partial Z$ has weight $1$, so $T_{F}^{\pm}$ are weight-$\frac{3}{2}$ primaries. For $j$, the only possible anomalous term in $T_{B}j$ is the $z^{-3}$ double contraction with $T_{\psi}=-\frac{1}{2}(\overline{\psi}\partial\psi+\psi\partial\overline{\psi})$:
		\begin{align*}
			\overline{\psi}\partial\psi(z)\,\overline{\psi}\psi(0)\big|_{z^{-3}}&=\Bigl(\partial_{z}\frac{1}{z}\Bigr)\frac{1}{z}=-\frac{1}{z^{3}}\:,\qquad
			\psi\partial\overline{\psi}(z)\,\overline{\psi}\psi(0)\big|_{z^{-3}}\eqs{5}-\frac{1}{z}\Bigl(\partial_{z}\frac{1}{z}\Bigr)=+\frac{1}{z^{3}}\:.
		\end{align*}
		At (5) the contractions $\psi(z)\overline{\psi}(0)$ and $\partial\overline{\psi}(z)\psi(0)$ cross, which gives a minus sign. The two terms cancel, so $T_{B}j$ has no $z^{-3}$ term and $j$ is a weight-1 primary, which is (\ref{11.1.4b}).
	\end{remark}
\end{tcolorbox}

The central charge of the $Z\psi\tilde{\psi}$ CFT is 3, so two such CFTs will cancel the ghost central charge. Since there are two real scalar fields in each CFT, the critical dimension is 4. However, these dimensions combine into complex pairs, so spacetime can only be pure Euclidean space or $(2,2)$ spacetime, and cannot be Minkowski $(3,1)$ spacetime. Furthermore, this theory only has 4-dimensional translation invariance but lacks 4-dimensional Lorentz invariance. The symmetry here is $U(2)$ or $U(1,1)$. Finally, the spectrum here is very small. Constraints completely determine the two sets of $Z\psi\tilde{\psi}$ oscillators. Therefore, only the center-of-mass motion of the singlet remains. This has some mathematical interest, but whether it has physical applications remains unclear.

Thus we return to the original $N=0$ and $N=1$ algebras. However, there is another generalization, where the left-moving and right-moving sectors of the closed string have different algebras. Holomorphic and antiholomorphic algebras commute with each other, so there is no reason they must be equal. In open strings, boundary conditions relate the holomorphic and antiholomorphic parts together, so open strings do not have a similar construction.

This gives a new possibility, $(N,\tilde{N})=(0,1)$ \emph{heterotic strings}; $(N,\tilde{N})=(1,0)$ is isomorphic to it and will not be examined separately. Additionally, $(0,2)$ heterotic string theory and $(1,2)$ heterotic string theory have some mathematical interest.

It should be emphasized that the analysis in this section has many explicit and implicit assumptions, and these assumptions should be treated with caution until we find all string theories. In fact, there are indeed some theories that do not fall into this classification. One is the Green-Schwarz form of superstrings. They do not have simple covariant gauge fixing, but in light-cone gauge, it can be proved via bosonization that they are actually equivalent to RNS superstrings. Another exception is \emph{topological string theory}, which does not satisfy the spin-statistics we assumed under covariant constraints. This string theory has no physical degrees of freedom, but because its observables are topological, it has some interesting mathematical applications.

\section{$SO(32)$ and $E_{8}\times E_{8}$ Heterotic Strings} \label{sec:11.2}

The $(0,1)$ heterotic string combines the ghost fields and constraints of the left-moving bosonic string and the right-moving Type II string. In fact, we can go further and combine the left-moving bosonic string and the right-moving Type II string such that the left side has 26 flat dimensions and the right side has 10 flat dimensions. This is achievable, but its physical meaning is unclear, so we keep the same dimensions on both sides. Thus it is 10. We start from the following fields
\begin{equation}
	X^{\mu}(z,\bar{z})\:,\qquad \tilde{\psi}^{\mu}(\bar{z}) \:, \qquad \mu = 0,\cdots,9\:, \label{11.2.1}
\end{equation}
The total central charge is $(c,\tilde{c})=(10,15)$. The ghost central charge is $(c^{\text{g}},\tilde{c}^{\text{g}})=(-26,-15)$, so the remaining matter fields have $(c,\tilde{c})=(16,0)$. The simplest possibility is to take 32 left-moving spin-$\frac{1}{2}$ fields
\begin{equation}
	\lambda^{A}(z) \:, \qquad A=1,\cdots,32\:. \label{11.2.2}
\end{equation}
The total matter field action is
\begin{equation}
	S=\frac{1}{4\pi} \int \dif^{2}z\: \biggl( \frac{2}{\alpha^{\prime}}\partial X^{\mu}\bar{\partial}X_{\mu}
	+\lambda^{A}\bar{\partial}\lambda^{A}+\tilde{\psi}^{\mu}\partial\tilde{\psi}_{\mu}\biggr)\:.\label{11.2.3}
\end{equation}
The operator products are
\begin{subequations}
	\begin{align}
		X^{\mu}(z,\bar{z})X^{\nu}(0,0) &\sim -\eta^{\mu\nu} \frac{\alpha^{\prime}}{2} \ln\lvert z\rvert^{2} \:,\label{11.2.4a} \\
		\lambda^{A}(z)\lambda^{B}(z) &\sim \delta^{AB}\frac{1}{z} \:, \label{11.2.4b} \\
		\tilde{\psi}^{\mu}(\bar{z})\tilde{\psi}^{\nu}(0) &\sim \eta^{\mu\nu}\frac{1}{\bar{z}} \:.\label{11.2.4c}
	\end{align}
\end{subequations}
The energy-momentum tensor and supercurrent are
\begin{subequations}
	\begin{align}
		T_{B} &= -\frac{1}{\alpha^{\prime}}\partial X^{\mu}\partial X_{\mu} -\frac{1}{2}\lambda^{A}\partial\lambda^{A}\:,\label{11.2.5a} \\
		\tilde{T}_{B}&= -\frac{1}{\alpha^{\prime}}\bar{\partial}X^{\mu}\bar{\partial}X_{\mu} 
		-\frac{1}{2}\tilde{\psi}^{\mu}\bar{\partial}\tilde{\psi}_{\mu} \:, \label{11.2.5b} \\
		\tilde{T}_{F}&= \mi(2/\alpha^{\prime})^{1/2}\tilde{\psi}^{\mu}\bar{\partial}X_{\mu} \:.\label{11.2.5c}
	\end{align} \label{11.2.5}
\end{subequations}
\begin{tcolorbox}[breakable]
	\begin{remark}
		Each free boson contributes $1$ to the central charge and each real fermion $\frac{1}{2}$, so
		\begin{equation*}
			c=10+32\cdot\tfrac{1}{2}=26\:,\qquad \tilde{c}=10+10\cdot\tfrac{1}{2}=15\:,
		\end{equation*}
		which cancel the ghosts $(c^{\text{g}},\tilde{c}^{\text{g}})=(-26,-15)$. The right-moving side must also close into the $N=1$ algebra with $\tilde{T}_{F}\tilde{T}_{F}\sim 2\tilde{c}/3\bar{z}^{3}+2\tilde{T}_{B}/\bar{z}$. Using $\bar{\partial}X^{\mu}(\bar{z})\bar{\partial}X^{\nu}(0)\sim-\frac{\alpha'}{2}\eta^{\mu\nu}/\bar{z}^{2}$,
		\begin{align*}
			\tilde{T}_{F}(\bar{z})\tilde{T}_{F}(0)&=-\frac{2}{\alpha'}\,\tilde{\psi}^{\mu}(\bar{z})\tilde{\psi}^{\nu}(0)\;\bar{\partial}X_{\mu}(\bar{z})\bar{\partial}X_{\nu}(0)\\
			&\eqs{1}-\frac{2}{\alpha'}\Bigl[\frac{\eta^{\mu\nu}}{\bar{z}}\Bigl(-\frac{\alpha'}{2}\frac{\eta_{\mu\nu}}{\bar{z}^{2}}\Bigr)
			+\frac{\eta^{\mu\nu}}{\bar{z}}\bar{\partial}X_{\mu}\bar{\partial}X_{\nu}
			-\frac{\alpha'}{2}\frac{1}{\bar{z}^{2}}:\!\tilde{\psi}^{\mu}(\bar{z})\tilde{\psi}_{\mu}(0)\!:\Bigr]\\
			&\eqs{2}\frac{10}{\bar{z}^{3}}+\frac{2}{\bar{z}}\Bigl(-\frac{1}{\alpha'}\bar{\partial}X^{\mu}\bar{\partial}X_{\mu}\Bigr)+\frac{1}{\bar{z}}\bar{\partial}\tilde{\psi}^{\mu}\tilde{\psi}_{\mu}
			\eqs{3}\frac{2\cdot 15}{3\bar{z}^{3}}+\frac{2}{\bar{z}}\tilde{T}_{B}(0)\:.
		\end{align*}
		At (1) the terms are the double contraction and the two single contractions. At (2) $\eta^{\mu\nu}\eta_{\mu\nu}=10$ and $:\!\tilde{\psi}^{\mu}(\bar{z})\tilde{\psi}_{\mu}(0)\!:=\bar{z}\,\bar{\partial}\tilde{\psi}^{\mu}\tilde{\psi}_{\mu}+O(\bar{z}^{2})$, since $\tilde{\psi}^{\mu}\tilde{\psi}_{\mu}=0$. At (3) $\bar{\partial}\tilde{\psi}^{\mu}\tilde{\psi}_{\mu}=-\tilde{\psi}^{\mu}\bar{\partial}\tilde{\psi}_{\mu}$. The result has $\tilde{c}=15$, as required.
	\end{remark}
\end{tcolorbox}

The world-sheet theory possesses $SO(9,1)\times SO(32)$ symmetry. $SO(32)$ acts on $\lambda^{A}$ and is an internal symmetry. $\lambda^{A}$ cannot have timelike characteristics because there are no fermi constraints to eliminate negative-norm states. Therefore, although the action for $\lambda^{A}$ is the same as for $\psi$ in the RNS superstring, the two theories are very different due to the differing constraints.

The right-moving ghosts are the same as in the RNS superstring, and the left-moving ghosts are the same as in the bosonic string. The proofs of the nilpotent BRST charge and the no-ghost theorem are straightforward.

To complete the description of this theory, we need to specify the boundary conditions on the fields and the sectors in the spectrum. The current situation is more complex than in Type II strings; now Poincaré invariance and BRST invariance do not require all $\lambda^{A}$ to have the same boundary conditions. The periodicity of $T_{B}$ only requires that $\lambda^{A}$ be periodic in the sense of a difference by an $O(32)$ rotation,
\begin{equation}
	\lambda^{A}(w+2\pi) = O^{AB}\lambda^{B}(w) \:. \label{11.2.6} 
\end{equation}
We do not perform a systematic analysis of compatible theories but will give all known theories. Nine theories are known based on action (\ref{11.2.3}), six of which have tachyons. Of the remaining three tachyon-free theories, two have spacetime supersymmetry; these two theories will be our subject.

In Type IIA and Type IIB superstrings, GSO projection acts separately on the left and right moves. In any supersymmetric heterotic theory, this is also true. The world-sheet current with spacetime symmetry is $\mathscr{V}_{\mathbf{s}}$ in equation (\ref{10.4.25}), where $\mathbf{s}$ belongs to $\mathbf{16}$. For the corresponding charge to be well-defined, the OPE of this current with any vertex operator must be single-valued. For the right-moving spinor part of this vertex operator, the spin eigenvalue $\mathbf{s}^{\prime}$ must satisfy
\begin{equation}
	\mathbf{s}\cdot \mathbf{s}^{\prime} + \frac{l}{2} \in \mathds{Z} \:, \label{11.2.7}
\end{equation}
for any $\mathbf{s}\in\mathbf{16}$, where $l$ equals $-1$ in the NS sector and $-\frac{1}{2}$ in the R sector. Taking $\mathbf{s}=(\frac{1}{2},\frac{1}{2},\frac{1}{2},\frac{1}{2},\frac{1}{2})$, this condition is precisely the right-moving GSO projection
\begin{equation}
	\exp(\pi\mi\tilde{F})=1\: ; \label{11.2.8}
\end{equation}
any other $\mathbf{s}\in\mathbf{16}$ gives the same condition.
\begin{tcolorbox}[breakable]
	\begin{remark}
		The right-moving part of a vertex operator in bosonized form is $\me^{l\tilde{\phi}}\me^{\mi\mathbf{s}'\cdot\tilde{\mathbf{H}}}$ times polynomials in derivatives, which do not affect the monodromy. Dropping the tildes, and with $\phi\phi\sim-\ln z$ and $H^{a}H^{b}\sim-\delta^{ab}\ln z$,
		\begin{align*}
			\mathscr{V}_{\mathbf{s}}(z)\,\me^{l\phi}\me^{\mi\mathbf{s}'\cdot\mathbf{H}}(0)
			&=\me^{-\phi/2}(z)\me^{l\phi}(0)\;\me^{\mi\mathbf{s}\cdot\mathbf{H}}(z)\me^{\mi\mathbf{s}'\cdot\mathbf{H}}(0)
			\eqs{1} z^{-(-\frac{1}{2})l}\,z^{\mathbf{s}\cdot\mathbf{s}'}\bigl[\cdots\bigr]
			= z^{\mathbf{s}\cdot\mathbf{s}'+l/2}\bigl[\cdots\bigr]\:.
		\end{align*}
		At (1) $\me^{q_{1}\phi}(z)\me^{q_{2}\phi}(0)\sim z^{-q_{1}q_{2}}$ and $\me^{\mi\mathbf{a}\cdot\mathbf{H}}(z)\me^{\mi\mathbf{b}\cdot\mathbf{H}}(0)\sim z^{\mathbf{a}\cdot\mathbf{b}}$ were used. The OPE is single-valued iff the exponent is an integer, which is (\ref{11.2.7}). For $\mathbf{s}=(\frac{1}{2})^{5}$,
		\begin{equation*}
			\mathbf{s}\cdot\mathbf{s}'+\frac{l}{2}=\frac{1}{2}\Bigl(\sum_{a}s'_{a}+l\Bigr)\in\mathds{Z}
			\quad\Longleftrightarrow\quad \exp\Bigl[\pi\mi\Bigl(\sum_{a}s'_{a}+l\Bigr)\Bigr]=1\:,
		\end{equation*}
		and $\sum_{a}s'_{a}+l$ is precisely $\tilde{F}$, since the ghost $\me^{l\phi}$ carries $F=l$ (see the discussion after (\ref{10.4.25})). For example, $\me^{-\phi}\tilde{\psi}^{\mu}$ has $\sum s'+l=\pm1-1\in 2\mathds{Z}$ and survives, while the tachyon $\me^{-\phi}$ has $0-1$ and is removed. Any other $\mathbf{s}\in\mathbf{16}$ is obtained by flipping an even number of signs, $\mathbf{s}=(\tfrac{1}{2})^{5}-\boldsymbol{\delta}$ with $\boldsymbol{\delta}$ having $1$ in the flipped places, and
		\begin{equation*}
			\mathbf{s}\cdot\mathbf{s}'=(\tfrac{1}{2})^{5}\cdot\mathbf{s}'-\sum_{a\in\text{flipped}}s'_{a}\:.
		\end{equation*}
		The last sum is an integer in the NS sector ($s'_{a}\in\mathds{Z}$) and in the R sector (an even number of half-integers), so the condition does not change.
	\end{remark}
\end{tcolorbox}

Now we also attach GSO projection to the left-moving spinors. That is, we take the same periodicity for all 32 components
\begin{equation}
	\lambda^{A}(w+2\pi)=\pm\lambda^{A}(w)  \:, \label{11.2.9}
\end{equation}
and attach to the left-moving fermion number
\begin{equation}
	\exp(\pi\mi F)=1\:. \label{11.2.10}
\end{equation}
It is easy to prove that the OPE is local and closed via bosonization, the method being the same as for Type IIA and Type IIB strings. Combining 32 real fermions into 16 complex fermions gives
\begin{equation}
	\lambda^{K\pm}=2^{-1/2}(\lambda^{2K-1}\pm\mi\lambda^{2K})\:, \qquad K=1,\cdots,16 \:. \label{11.2.11}
\end{equation}
These can be bosonized and represented with 16 left-moving scalars $H^{K}(z)$. By analogy with the definition of $F$ in Type II strings, define
\begin{equation}
	F=\sum_{K=1}^{16}q_{K} \:, \label{11.2.12}
\end{equation}
where $\lambda^{K\pm}$ have $q_{K}=\pm1$. $F$ is additive, so the OPE is closed, and the projection (\ref{11.2.10}) ensures there is no branch cut with R-sector vertex operators. In the description of bosonization, we have 26 left-moving bosons and 10 right-moving bosons, so (\ref{11.2.3}) is actually a fusion of the bosonic string and Type II string.

Modular invariance is straightforward. The partition function of $\lambda$ is
\begin{equation}
	Z_{16}(\tau)= \frac{1}{2}\Bigl[
	Z^{0}{}_{0}(\tau)^{16}+ Z^{0}{}_{1}(\tau)^{16}+Z^{1}{}_{0}(\tau)^{16}+Z^{1}{}_{1}(\tau)^{16}
	\Bigr] \:. \label{11.2.13}
\end{equation}
\begin{tcolorbox}[breakable]
	\begin{remark}
		Here each $Z^{\alpha}{}_{\beta}$ is the partition function of one complex (two real) fermion, so $32$ real fermions give the $16$th power. Reading off (\ref{10.7.14}) for the four spin structures,
		\begin{align*}
			Z^{0}{}_{0}(\tau+1)&=\me^{-\pi\mi/12}Z^{0}{}_{1}(\tau)\:,& Z^{0}{}_{1}(\tau+1)&=\me^{-\pi\mi/12}Z^{0}{}_{0}(\tau)\:,\\
			Z^{1}{}_{0}(\tau+1)&=\me^{\pi\mi/6}Z^{1}{}_{0}(\tau)\:,& Z^{1}{}_{1}(\tau+1)&=\me^{\pi\mi/6}Z^{1}{}_{1}(\tau)\:,\\
			Z^{0}{}_{0}(-1/\tau)&=Z^{0}{}_{0}(\tau)\:,& Z^{0}{}_{1}(-1/\tau)&=Z^{1}{}_{0}(\tau)\:,\qquad Z^{1}{}_{0}(-1/\tau)=Z^{0}{}_{1}(\tau)\:.
		\end{align*}
		The phases are $\me^{2\pi\mi E_{0}}$ with $E_{0}=-\frac{1}{24}$ (NS) and $+\frac{1}{12}$ (R), the zero-point energies of a complex fermion. $Z^{1}{}_{1}$ vanishes identically because of the R zero modes, so its phase under $S$ does not matter. Raising to the 16th power,
		\begin{align*}
			Z_{16}(\tau+1)&=\frac{1}{2}\Bigl[\me^{-16\pi\mi/12}\bigl(Z^{0}{}_{1}{}^{16}+Z^{0}{}_{0}{}^{16}\bigr)+\me^{16\pi\mi/6}\bigl(Z^{1}{}_{0}{}^{16}+Z^{1}{}_{1}{}^{16}\bigr)\Bigr]
			\eqs{1}\me^{2\pi\mi/3}Z_{16}(\tau)\:,\\
			Z_{16}(-1/\tau)&=\frac{1}{2}\Bigl[Z^{0}{}_{0}{}^{16}+Z^{1}{}_{0}{}^{16}+Z^{0}{}_{1}{}^{16}+Z^{1}{}_{1}{}^{16}\Bigr]=Z_{16}(\tau)\:.
		\end{align*}
		At (1) $\me^{-4\pi\mi/3}=\me^{8\pi\mi/3}=\me^{2\pi\mi/3}$. The NS and R phases coincide only because the power is a multiple of 8. On the right, the same rules give $Z_{\psi}^{+}(\tau+1)=\me^{2\pi\mi/3}Z_{\psi}^{+}(\tau)$, so $Z_{\psi}^{+}(\tau)^{\ast}$ picks up $\me^{-2\pi\mi/3}$. The transverse bosons enter as $Z_{X}^{8}$, which is modular invariant by itself. Therefore
		\begin{equation*}
			Z_{X}(\tau)^{8}\,Z_{16}(\tau)\,Z_{\psi}^{+}(\tau)^{\ast}\;\xrightarrow{\;\tau\to\tau+1\;}\;\me^{2\pi\mi/3}\me^{-2\pi\mi/3}\,Z_{X}^{8}Z_{16}Z_{\psi}^{+\ast}=Z_{X}^{8}Z_{16}Z_{\psi}^{+\ast}\:,
		\end{equation*}
		and the one-loop amplitude is modular invariant.
	\end{remark}
\end{tcolorbox}
The effect of modular transformation is only the permutation of these four terms; there is no phase under $\tau\to -1/\tau$, and a phase of $\exp(2\pi\mi/3)$ is generated under $\tau\to\tau+1$. This phase will cancel with the phase of the partition function $Z_{\psi}^{+}(\tau)^{\ast}$ of $\tilde{\psi}$. The form of (\ref{11.2.13}) is similar to $Z_{\psi}^{+}(\tau)$ of Type II strings, but now the signs of all terms are positive. From several perspectives, this sign is inevitable. Now there are 32 fermions instead of 8, so the sign in the modular transformation must be raised to the fourth power, and thus the first three signs must be the same. The $Z^{1}{}_{1}$ term transforms to itself, so as usual, its sign depends on the chirality of the R sector. The other three theories are defined by flipping the left-side or right-side R sector; they are physically equivalent. Additionally, the first and second terms in $Z_{\psi}^{+}(\tau)$ have opposite signs due to $F$ of the superconformal ghosts, but for heterotic strings, there are no such ghosts on the left side. The relative sign between the first and third terms comes from spacetime statistics, but $\lambda$ is a spacetime scalar, so their R-sector states are too. Therefore, modular invariance, OPE conservation of $F$, and spacetime spin-statistics are all consistent with the partition function (\ref{11.2.13}).

We now look for the lightest states. The right-moving part is the same as in Type II strings, with no tachyon, and the lowest level is at $\mathbf{8}_{v}+\mathbf{8}$. On the left-moving side, the normal ordering constant in the left-moving transverse Hamiltonian $H^{\bot}=\alpha^{\prime}m^{2}/4$ is
\begin{equation}
	\text{NS}:\quad -\frac{8}{24}-\frac{32}{48} =-1\:, \qquad \text{R}:\quad -\frac{8}{24}+\frac{32}{24}=+1\:.\label{11.2.14} 
\end{equation}
\begin{tcolorbox}[breakable]
	\begin{remark}
		The zero-point energy of a field is $+\frac{1}{2}$ (boson) or $-\frac{1}{2}$ (fermion) times the sum of its positive mode frequencies, regulated with the Hurwitz zeta function $\zeta(-1,a)=-\frac{1}{2}B_{2}(a)=-\frac{1}{2}(a^{2}-a+\frac{1}{6})$:
		\begin{align*}
			\text{periodic boson:}&\quad \tfrac{1}{2}\sum_{n=1}^{\infty}n=\tfrac{1}{2}\zeta(-1)=-\tfrac{1}{24}\:,\\
			\text{NS fermion:}&\quad -\tfrac{1}{2}\sum_{n=1}^{\infty}\bigl(n-\tfrac{1}{2}\bigr)=-\tfrac{1}{2}\zeta\bigl(-1,\tfrac{1}{2}\bigr)\eqs{1}-\tfrac{1}{48}\:,\\
			\text{R fermion:}&\quad -\tfrac{1}{2}\sum_{n=1}^{\infty}n=-\tfrac{1}{2}\zeta(-1)=+\tfrac{1}{24}\:.
		\end{align*}
		At (1) $B_{2}(\frac{1}{2})=\frac{1}{4}-\frac{1}{2}+\frac{1}{6}=-\frac{1}{12}$, so $\zeta(-1,\frac{1}{2})=\frac{1}{24}$. In light-cone gauge the left movers are 8 transverse bosons and the 32 $\lambda^{A}$, so
		\begin{equation*}
			\text{NS}:\ 8\bigl(-\tfrac{1}{24}\bigr)+32\bigl(-\tfrac{1}{48}\bigr)=-\tfrac{1}{3}-\tfrac{2}{3}=-1\:,\qquad
			\text{R}:\ 8\bigl(-\tfrac{1}{24}\bigr)+32\bigl(+\tfrac{1}{24}\bigr)=-\tfrac{1}{3}+\tfrac{4}{3}=+1\:,
		\end{equation*}
		which is (\ref{11.2.14}).
	\end{remark}
	Note that the normal ordering constant of the R sector is already $+1$, which means that even the lightest state in the R sector is massive.
\end{tcolorbox}
\noindent So the left-moving NS ground state is a tachyon. The first excited state is
\begin{equation}
	\lambda^{A}_{-1/2}\lvert 0\rangle_{\text{NS}}  \label{11.2.15}
\end{equation}
which has $H^{\bot}=-\frac{1}{2}$, but it is removed by projection (\ref{11.2.10}): since there is now no ghost contribution, the NS ground state now has $\exp(\pi\mi F)=+1$.
States with $H^{\bot}=0$ can now be obtained in two ways:
\begin{equation}
	\alpha_{-1}^{i}\lvert 0\rangle_{\text{NS}}\:,\qquad \lambda_{-1/2}^{A}\lambda_{-1/2}^{B}\lvert 0\rangle_{\text{NS}} \:. \label{11.2.16}    
\end{equation}
$\lambda^{A}$ transforms under internal $SO(32)$ symmetry. Under the full symmetry $SO(8)_{\text{spin}}\times SO(32)$, the NS ground state is invariant $(\mathbf{1},\mathbf{1})$. The second state in (\ref{11.2.16}) is antisymmetric under $A\leftrightarrow B$, so the transformation of the massless state (\ref{11.2.16}) is $(\mathbf{8}_{v},\mathbf{1})+(\mathbf{1},[2])$. The antisymmetric tensor representation is the adjoint representation of $SO(32)$, with a dimension of $32\times 31/2=\mathbf{496}$.

\begin{table}[h]
	\caption{Low-energy states of heterotic strings}
	\label{tab:11.2}%
	\centering
	\begin{tabular}[c]{ccccc}
		\hline\hline
		$\vphantom{\Bigl[}m^{2}$ & NS & R & $\widetilde{\text{NS}}$ & $\widetilde{\text{R}}$  \\
		\hline
		$-4/\alpha^{\prime}$ & $(\mathbf{1},\mathbf{1})$ & - & - &  -  \\
		0 & $(\mathbf{8}_{v},\mathbf{1})+(\mathbf{1},\mathbf{496})$ & - & $\mathbf{8}_{v}$ & $\mathbf{8}$   \\
		\hline\hline
	\end{tabular}
\end{table}

Table \ref{tab:11.2} summarizes the tachyon states and massless states on both sides. The left-moving part is represented by its $SO(8)\times SO(32)$ quantum numbers, while the right-moving part is represented by its $SO(8)$ quantum numbers. Closed strings must combine left-moving and right-moving states with the same mass. The left-moving part originally had a tachyon as in the bosonic string, but since there is no tachyon state on the right to pair with it, this theory is tachyon-free. At the massless level, the product
\begin{equation}
	(\mathbf{8}_{v},\mathbf{1})\times (\mathbf{8}_{v}+\mathbf{8}) =
	(\mathbf{1},\mathbf{1})+(\mathbf{28},\mathbf{1})+(\mathbf{35},\mathbf{1})+(\mathbf{56},\mathbf{1})+(\mathbf{8}^{\prime},\mathbf{1})\label{11.2.17}
\end{equation}
is the supergravity multiplet of Type I strings. The product
\begin{equation}
	(\mathbf{1},\mathbf{496})\times (\mathbf{8}_{v}+\mathbf{8}) = (\mathbf{8}_{v},\mathbf{496}) + (\mathbf{8},\mathbf{496})\label{11.2.18}
\end{equation}
is the $N=1$ gauge multiplet in the $SO(32)$ adjoint representation. Therefore, the latter is the gauge symmetry in spacetime.
\begin{tcolorbox}[breakable]
	\begin{remark}
		The $SO(8)$ products follow by splitting tensors into irreducible pieces. For two vectors,
		\begin{align*}
			\mathbf{8}_{v}\times\mathbf{8}_{v}:\quad e^{i}\tilde{e}^{j}
			&=\underbrace{\tfrac{1}{8}\delta^{ij}e^{k}\tilde{e}^{k}}_{\mathbf{1}}
			+\underbrace{\tfrac{1}{2}\bigl(e^{i}\tilde{e}^{j}-e^{j}\tilde{e}^{i}\bigr)}_{\mathbf{28}}
			+\underbrace{\Bigl[\tfrac{1}{2}\bigl(e^{i}\tilde{e}^{j}+e^{j}\tilde{e}^{i}\bigr)-\tfrac{1}{8}\delta^{ij}e^{k}\tilde{e}^{k}\Bigr]}_{\mathbf{35}}\:,
		\end{align*}
		with $1+\frac{8\cdot7}{2}+\bigl(\frac{8\cdot9}{2}-1\bigr)=64$. These are the dilaton, the antisymmetric tensor and the graviton. For a vector times a spinor $u_{\alpha}$ in the $\mathbf{8}$,
		\begin{align*}
			\mathbf{8}_{v}\times\mathbf{8}:\quad e^{i}u_{\alpha}
			&=\underbrace{\tfrac{1}{8}\bigl(\Gamma^{i}\Gamma^{j}e^{j}u\bigr)_{\alpha}}_{\mathbf{8}'}
			+\underbrace{\Bigl[e^{i}u_{\alpha}-\tfrac{1}{8}\bigl(\Gamma^{i}\Gamma^{j}e^{j}u\bigr)_{\alpha}\Bigr]}_{\mathbf{56}}\:.
		\end{align*}
		The spinor $\chi=\Gamma^{j}e^{j}u$ has the opposite chirality because $\Gamma^{j}$ flips chirality, so it is an $\mathbf{8}'$. The remainder $\rho^{i}$ is gamma-traceless, $\Gamma^{i}\rho^{i}=\Gamma^{j}e^{j}u-\frac{1}{8}\Gamma^{i}\Gamma^{i}\chi=\chi-\chi=0$, using $\Gamma^{i}\Gamma^{i}=8$. That removes 8 of the 64 components and leaves $56$, the gravitino. Its $\mathbf{8}'$ partner is the dilatino. In (\ref{11.2.18}) the $SO(32)$ label is a spectator and $(\mathbf{1},\mathbf{496})\times(\mathbf{8}_{v}+\mathbf{8})$ just multiplies the $\mathbf{496}$ by $\mathbf{8}_{v}+\mathbf{8}$.
	\end{remark}
\end{tcolorbox}

This is exactly the same as the massless spectrum of the Type I open plus closed $SO(32)$ theory. However, the massive spectra of these two theories are different. In open strings, gauge quantum numbers are carried by the endpoints of the $SO(32)$ vector, so even at massive energy levels, no higher representations exist other than the second-rank tensor representation of that gauge group. But in heterotic strings, gauge quantum numbers are carried by fields, which propagate across the entire world-sheet. At massive energy levels, there can be any number of such excitations, which would give arbitrarily large gauge group representations. However, we will see in Chapter 14 that Type I theory and heterotic $SO(32)$ theory are actually identical.

The second type of heterotic string theory divides $\lambda^{A}$ into two groups of 16 each and attaches independent boundary conditions to each group,
\begin{equation}
	\lambda^{A}(w+2\pi) = \left\{
	\begin{array}{l}
		\eta \lambda^{A}(w) \: ,\\
		\eta^{\prime} \lambda^{A}(w) \:, 
	\end{array} \quad 
	\begin{array}{l}
		A=1,\cdots,16\:,  \\
		A=17,\cdots,32 \:, 
	\end{array}
	\right. \label{11.2.19}
\end{equation}
where $\eta$ and $\eta^{\prime}$ are each $\pm1$. Correspondingly, there exist operators
\begin{equation}
	\exp(\pi\mi F_{1}) \:, \qquad \exp(\pi\mi F_{1}^{\prime}) \:, \label{11.2.20} 
\end{equation}
which anticommute with the first set and second set of $\lambda^{A}$, respectively. We take separate GSO projections for the right-moving and the two left-moving parts. That is, summing over $2^{3}=8$ different periodicity combinations under the projection
\begin{equation}
	\exp(\pi\mi F_{1}) =\exp (\pi\mi F_{1}^{\prime}) =\exp(\pi\mi \tilde{F}) = 1  \label{11.2.21}
\end{equation}
OPE closure and locality, as well as modular invariance, are easily proved. In particular, the partition function
\begin{equation}
	Z_{8}(\tau)^{2} = \frac{1}{4} \Bigl[Z^{0}{}_{0}(\tau)^{8} + Z^{0}{}_{1}(\tau)^{8}
	+Z^{1}{}_{0}(\tau)^{8} + Z^{1}{}_{1}(\tau)^{8} \Bigr]^{2} \label{11.2.22}
\end{equation}
has the same transformations as $Z_{\psi}^{\pm}$ and $Z_{16}$. What is important here is that the fermions are in groups of 16, which results in the signs in $Z_{\psi}^{\pm}$ being squared.
\begin{tcolorbox}[breakable]
	\begin{remark}
		Each group of 16 real fermions is 8 complex fermions. With the rules of (\ref{10.7.14}) listed after (\ref{11.2.13}),
		\begin{align*}
			Z_{8}(\tau+1)&=\frac{1}{2}\Bigl[\me^{-8\pi\mi/12}Z^{0}{}_{1}{}^{8}+\me^{-8\pi\mi/12}Z^{0}{}_{0}{}^{8}+\me^{8\pi\mi/6}Z^{1}{}_{0}{}^{8}+\me^{8\pi\mi/6}Z^{1}{}_{1}{}^{8}\Bigr]
			\eqs{1}\me^{-2\pi\mi/3}Z_{8}(\tau)\:,\\
			Z_{8}(\tau+1)^{2}&=\me^{-4\pi\mi/3}Z_{8}(\tau)^{2}=\me^{2\pi\mi/3}Z_{8}(\tau)^{2}\:,\qquad Z_{8}(-1/\tau)^{2}=Z_{8}(\tau)^{2}\:.
		\end{align*}
		At (1) $\me^{-2\pi\mi/3}=\me^{4\pi\mi/3}$, so the NS and R terms again pick up a common phase. The result is the same phase $\me^{2\pi\mi/3}$ as $Z_{16}$, which again cancels against $Z_{\psi}^{+}(\tau)^{\ast}$. By contrast, for a group of 8 real fermions (4 complex) the NS and R phases are $\me^{-\pi\mi/3}$ and $\me^{2\pi\mi/3}=-\me^{-\pi\mi/3}$. They differ by a sign, which is why $Z_{\psi}^{\pm}$ needs the relative minus signs.
	\end{remark}
\end{tcolorbox}

As before, the lightest states on the right side are massless $\mathbf{8}_{v}+\mathbf{8}$. On the left, the NS-NS$^{\prime}$ sector still has a normal ordering constant of $-1$, so the ground state is a tachyon, but there is no state on the right to match it. The first excited states are at $m^{2}=0$:
\begin{align}
	& \alpha_{-1}^{i}\lvert0\rangle_{\text{NS,NS}^{\prime}} \:, \nonumber \\
	& \lambda^{A}_{-1/2}\lambda^{B}_{-1/2} \lvert0\rangle_{\text{NS,NS}^{\prime}} \:, \quad
	1\leq A,B\leq 16 \text{ or } 17\leq A,B \leq 32 \:. \label{11.2.23}
\end{align}
This has one difference from the $SO(32)$ case: because there are separate GSO projections on the two sets of $\lambda$, $A$ and $B$ must come from the same set. Since $SO(32)$ is partially broken by the boundary conditions, we classify the states according to the remaining $SO(16)\times SO(16)$. The states (\ref{11.2.23}) contain one antisymmetric tensor adjoint representation for each $SO(16)$, with a dimension of $16\times 15/2=\mathbf{120}$.

In the left-moving R-NS$^{\prime}$ sector, the normal ordering constant is
\begin{equation}
	-\frac{8}{24} + \frac{16}{24} - \frac{16}{48} =0 \:, \label{11.2.24}
\end{equation}
so the ground state is massless. Using 8 raising operators and 8 lowering operators to build out the 16 $\lambda^{A}$ zero modes gives the first $SO(16)$ 256-dimensional spinor representation. GSO projection divides it into two irreducible representations, $\mathbf{128}+\mathbf{128}^{\prime}$, the former being in the spectrum. The NS-R$^{\prime}$ sector gives another $SO(16)$ $\mathbf{128}$, and the R-R$^{\prime}$ sector still has no massless states.
\begin{tcolorbox}[breakable]
	\begin{remark}
		With the zero-point energies of the box after (\ref{11.2.14}), $8$ bosons, $16$ R fermions and $16$ NS fermions give
		\begin{equation*}
			8\bigl(-\tfrac{1}{24}\bigr)+16\bigl(+\tfrac{1}{24}\bigr)+16\bigl(-\tfrac{1}{48}\bigr)=-\tfrac{1}{3}+\tfrac{2}{3}-\tfrac{1}{3}=0\:,
		\end{equation*}
		and the R-R$'$ sector has $8(-\frac{1}{24})+32(\frac{1}{24})=+1$, which is massive. The R ground states come from the zero modes $\lambda^{A}_{0}$, $A=1,\dots,16$, which obey $\{\lambda_{0}^{A},\lambda_{0}^{B}\}=\delta^{AB}$. Pair them as in (\ref{11.2.11}):
		\begin{gather*}
			\lambda_{0}^{K\pm}=2^{-1/2}\bigl(\lambda_{0}^{2K-1}\pm\mi\lambda_{0}^{2K}\bigr)\:,\qquad K=1,\dots,8\:,\\
			\{\lambda_{0}^{K+},\lambda_{0}^{L-}\}=\delta^{KL}\:,\qquad \{\lambda_{0}^{K\pm},\lambda_{0}^{L\pm}\}=0\:.
		\end{gather*}
		These are 8 fermionic raising and lowering operators. Starting from the state annihilated by all $\lambda_{0}^{K-}$ they generate $2^{8}=256$ states, labeled by $q_{K}=\pm\frac{1}{2}$, the $SO(16)$ spinor weights. $\exp(\pi\mi F_{1})$ acts on them as the $SO(16)$ chirality $\prod_{K}(2q_{K})$ up to an overall sign convention. The projection keeps the $2^{7}=128$ weights with an even number of $-\frac{1}{2}$, the $\mathbf{128}$, and removes the $\mathbf{128}'$.
	\end{remark}
\end{tcolorbox}

To summarize, the massless $SO(8)_\text{spin }\times SO(16)\times SO(16)$ representations on the left side are
\begin{equation}
	(\mathbf{8}_{v},\mathbf{1},\mathbf{1})+(\mathbf{1},\mathbf{120},\mathbf{1})
	+(\mathbf{1},\mathbf{1},\mathbf{120})+(\mathbf{1},\mathbf{128},\mathbf{1})
	+(\mathbf{1},\mathbf{1},\mathbf{128}) \:. \label{11.2.25}
\end{equation}
Combining with the right-moving $\mathbf{8}_{v}$ gives massless vector bosons for each $SO(16)$, their transformations being $\mathbf{120}+\mathbf{128}$. Consistency of the spacetime theory requires massless vectors to transform according to the \emph{adjoint} representation of the gauge group. The exceptional group $E_{8}$ satisfies this requirement; $SO(16)$ is one of its subgroups, and under this subgroup $E_{8}$ transforms according to $\mathbf{120}+\mathbf{128}$. Clearly, the second heterotic string theory possesses the gauge group $E_{8}\times E_{8}$. Even though only $SO(16)\times SO(16)$ symmetry is explicit in the world-sheet theory under the current description, the world-sheet theory possesses full $E_{8}\times E_{8}$ symmetry. Additional currents are given by bosonization
\begin{equation}
	\exp \left[ \mi\sum_{K=1}^{16}q_{K}H^{K}(z) \right] \:. \label{11.2.26}
\end{equation}
This is a spin field, identical to the spin field in the fermion vertex operator (\ref{10.4.25}). For the first $E_{8}$, the charges are
\begin{equation}
	q_{K} = \left\{
	\begin{array}{l}
		\pm\tfrac{1}{2} \:,  \\
		0  \:,
	\end{array} \quad
	\begin{array}{l}
		K=1,\cdots,8  \\
		K=9,\cdots, 16 
	\end{array} \:,
	\qquad 
	\sum_{K=1}^{16} q_{K}\in 2\mathds{Z} \:, \label{11.2.27}
	\right.
\end{equation}
and similarly for the second. They are indeed $(1,0)$ operators. The massless spectrum consists of the $d=10, N=1$ supergravity multiplet plus the $E_{8}\times E_{8}$ gauge multiplet. The $SO(8)_{\text{spin}}\times E_{8}\times E_{8}$ quantum numbers of the massless fields are
\begin{align}
	&\quad (\mathbf{1},\mathbf{1},\mathbf{1})+(\mathbf{28},\mathbf{1},\mathbf{1})+(\mathbf{35},\mathbf{1},\mathbf{1})
	+(\mathbf{56},\mathbf{1},\mathbf{1})+(\mathbf{8}^{\prime},\mathbf{1},\mathbf{1})  \nonumber \\
	&+(\mathbf{8}_{v},\mathbf{248},\mathbf{1})+(\mathbf{8},\mathbf{248},\mathbf{1})
	+(\mathbf{8}_{v},\mathbf{1},\mathbf{248})+(\mathbf{8},\mathbf{1},\mathbf{248}) \:. \label{11.2.28}
\end{align}
\begin{tcolorbox}[breakable]
	\begin{remark}
		\textbf{Weight.} With $H^{K}(z)H^{L}(0)\sim-\delta^{KL}\ln z$, the exponential $\me^{\mi q\cdot H}$ has weight $\frac{1}{2}q^{2}$. For (\ref{11.2.27}), $q^{2}=8\cdot\frac{1}{4}=2$, so $h=1$: these are $(1,0)$ currents.

		\textbf{Counting.} There are $2^{8}$ sign choices. Writing $n_{-}$ for the number of $-\frac{1}{2}$ entries,
		\begin{equation*}
			\sum_{K=1}^{8}q_{K}=\tfrac{1}{2}\bigl[(8-n_{-})-n_{-}\bigr]=4-n_{-}\in2\mathds{Z}\quad\Longleftrightarrow\quad n_{-}\ \text{even}\:,
		\end{equation*}
		which leaves $2^{7}=128$ currents, the $\mathbf{128}$ of $SO(16)$ found in the R-NS$'$ sector.

		\textbf{Closure.} Two such currents have
		\begin{align*}
			\me^{\mi q\cdot H}(z)\,\me^{\mi q'\cdot H}(0)&\eqs{1}z^{q\cdot q'}\,\me^{\mi(q+q')\cdot H}(0)+\cdots\:,\qquad
			q\cdot q'\eqs{2}\tfrac{1}{4}\bigl[(8-d)-d\bigr]=2-\tfrac{d}{2}\:,
		\end{align*}
		where $d$ is the number of places where $q$ and $q'$ differ. At (1) the exponentials were contracted as in the box after (\ref{11.2.8}). At (2) agreeing entries contribute $+\frac{1}{4}$ and differing ones $-\frac{1}{4}$. Both have $n_{-}$ even, so $d$ is even and $q\cdot q'\in\mathds{Z}$. The OPE is therefore local. The singular cases are:
		\begin{align*}
			d=8\ (q'=-q):&\quad z^{-2}\bigl[1+\mi z\,q\cdot\partial H+\cdots\bigr]\:,\ \text{central term and Cartan currents }\mi\partial H^{K}\:,\\
			d=6:&\quad z^{-1}\,\me^{\mi(q+q')\cdot H}\:,\qquad q+q'=(\pm1,\pm1,0^{6})\ \text{up to order}\:.
		\end{align*}
		In the second case $q+q'$ is a root (\ref{11.4.16}) of $SO(16)$, i.e.\ the current is one of the fermion bilinears $\lambda^{K\pm}\lambda^{L\pm}$. The 128 spin currents together with the 120 $SO(16)$ currents therefore close under the OPE into a current algebra of dimension
		\begin{equation*}
			120+128=248=\operatorname{dim}(E_{8})\:.
		\end{equation*}
		The OPE of an $SO(16)$ root current $\me^{\mi\alpha\cdot H}$ with a spin current has $\alpha\cdot q\in\{0,\pm1\}$, so it again gives either nothing singular or another spin current; the $\mathbf{128}$ is a representation of $SO(16)$, as it must be.
	\end{remark}
\end{tcolorbox}

Consistency requires fermions to be in groups of 16. We can give a modular invariant theory using groups of 8 fermions, thus the left-moving partition function is $(Z_{\psi}^{\pm})^{4}$. However, we have seen that modular invariance requires the relative signs in $Z_{\psi}^{\pm}$. These signs would endow the left-moving R sector with a negative weight and correspond to the projection $\exp(\pi\mi F)$ on the NS sector. The former is inconsistent with spin-statistics because these states are spacetime scalars, and the latter is inconsistent with OPE closure, which makes interactions inconsistent. Therefore $SO(32)$ and $E_{8}\times E_{8}$ theories are the only supersymmetric heterotic strings in 10 dimensions.

\section{Other Ten-Dimensional Heterotic Strings} \label{sec:11.3}

Via the twist construction introduced in section 8.5, other heterotic string theories can be constructed from a single theory. The "minimal twist theory," namely the theory with the fewest path integral sectors, corresponds to the diagonal modular invariant
\begin{align}
	&\frac{1}{2}\Bigl[
	Z^{0}{}_{0}(\tau)^{16}Z^{0}{}_{0}(\tau)^{\ast 4}- Z^{0}{}_{1}(\tau)^{16}Z^{0}{}_{1}(\tau)^{\ast 4} \nonumber \\
	&\qquad \qquad - Z^{1}{}_{0}(\tau)^{16}Z^{1}{}_{0}(\tau)^{\ast 4}-Z^{1}{}_{1}(\tau)^{16}Z^{1}{}_{1}(\tau)^{\ast 4}
	\Bigr]  \:. \label{11.3.1}
\end{align}
\begin{tcolorbox}[breakable]
	\begin{remark}
		Each term pairs the same spin structure on 32 left-moving and 8 right-moving real fermions, i.e.\ 16 and 4 complex fermions. Under $\tau\to\tau+1$ the rules listed after (\ref{11.2.13}) give the phases
		\begin{align*}
			\text{NS}:&\quad \me^{-16\pi\mi/12}\,\bigl(\me^{-4\pi\mi/12}\bigr)^{\ast}=\me^{-\pi\mi}=-1\:,\qquad
			\text{R}:\quad \me^{16\pi\mi/6}\,\bigl(\me^{4\pi\mi/6}\bigr)^{\ast}=\me^{2\pi\mi}=+1\:.
		\end{align*}
		Writing $A_{\beta}^{\alpha}=Z^{\alpha}{}_{\beta}{}^{16}Z^{\alpha}{}_{\beta}{}^{\ast4}$, the bracket in (\ref{11.3.1}) transforms as
		\begin{align*}
			A^{0}_{0}-A^{0}_{1}-A^{1}_{0}-A^{1}_{1}
			\;&\xrightarrow{\tau\to\tau+1}\;(-A^{0}_{1})-(-A^{0}_{0})-A^{1}_{0}-A^{1}_{1}=A^{0}_{0}-A^{0}_{1}-A^{1}_{0}-A^{1}_{1}\:,\\
			\;&\xrightarrow{\tau\to-1/\tau}\;A^{0}_{0}-A^{1}_{0}-A^{0}_{1}-A^{1}_{1}=A^{0}_{0}-A^{0}_{1}-A^{1}_{0}-A^{1}_{1}\:.
		\end{align*}
		So the diagonal combination is invariant without any cancellation between left and right. That is the general property of left-right symmetric boundary conditions. Conversely, the relative minus sign between $A^{0}_{0}$ and $A^{0}_{1}$ is forced by $T$, and the equality of the coefficients of $A^{0}_{1}$ and $A^{1}_{0}$ by $S$.
	\end{remark}
\end{tcolorbox}
Corresponding to this modular invariant, for all fermions, $\lambda^{A}$ and $\tilde{\psi}^{\mu}$, all are taken simultaneously to be periodic or antiperiodic for each cycle on the torus. Written in terms of the spectrum, world-sheet fermions are either all R or all NS, and the diagonal projection is
\begin{equation}
	\exp[\pi\mi(F+\tilde{F})] =1 \:. \label{11.3.2}
\end{equation}
Except for a tachyon, this theory is self-consistent; the state
\begin{equation}
	\lambda_{-1/2}^{A}\lvert 0\rangle_{\text{NS,NS}}\:, \qquad m^{2}=-\frac{2}{\alpha^{\prime}}\:,\qquad
	\exp(\pi\mi F)=\exp(\pi\mi\tilde{F})=-1\:,  \label{11.3.3}
\end{equation}
transforms according to a vector under $SO(32)$. At the massless energy level, the states are
\begin{equation}
	\alpha^{i}_{-1}\tilde{\psi}^{j}_{-1/2} \lvert0\rangle_{\text{NS,NS}} \:, \qquad
	\lambda^{A}_{-1/2}\lambda_{-1/2}^{B}\tilde{\psi}^{j}_{-1/2}\lvert 0\rangle_{\text{NS,NS}} \:, \label{11.3.4}
\end{equation}
which are the graviton, dilaton, antisymmetric tensor, and $SO(32)$ gauge boson. There are fermions in the spectrum, but the lightest state resides at $m^{2}=4/\alpha^{\prime}$.
\begin{tcolorbox}[breakable]
	\begin{remark}
		A closed-string state has $\frac{\alpha'}{4}m^{2}=L_{0}=\tilde{L}_{0}$ (in light-cone gauge, including zero-point energies). From (\ref{11.2.14}) and its right-moving analog,
		\begin{equation*}
			L_{0}=N-1\ (\text{NS}),\quad N+1\ (\text{R})\:;\qquad \tilde{L}_{0}=\tilde{N}-\tfrac{1}{2}\ (\widetilde{\text{NS}}),\quad \tilde{N}\ (\widetilde{\text{R}})\:.
		\end{equation*}
		On the right the NS ground state has $\exp(\pi\mi\tilde{F})=-1$ (ghost contribution), while on the left there are no superconformal ghosts and the NS ground state has $\exp(\pi\mi F)=+1$. Each $\lambda_{-1/2}$ or $\tilde{\psi}_{-1/2}$ flips the sign.
		\begin{itemize}
			\item Tachyon: $\tilde{L}_{0}=-\frac{1}{2}$ requires the right NS ground state, with $\exp(\pi\mi\tilde{F})=-1$. Then (\ref{11.3.2}) requires $\exp(\pi\mi F)=-1$ with $L_{0}=-\frac{1}{2}$, which is $\lambda^{A}_{-1/2}\lvert0\rangle_{\text{NS}}$. So $m^{2}=\frac{4}{\alpha'}(-\frac{1}{2})=-\frac{2}{\alpha'}$ in the $\mathbf{32}$ of $SO(32)$. The left ground state ($L_{0}=-1$) has no partner.
			\item Massless: $\tilde{\psi}^{j}_{-1/2}\lvert0\rangle$ has $\exp(\pi\mi\tilde{F})=+1$, so it pairs with left states with $L_{0}=0$ and $\exp(\pi\mi F)=+1$, namely $\alpha^{i}_{-1}\lvert0\rangle$ and $\lambda^{A}_{-1/2}\lambda^{B}_{-1/2}\lvert0\rangle$. The first gives $\mathbf{8}_{v}\times\mathbf{8}_{v}=\mathbf{1}+\mathbf{28}+\mathbf{35}$, the second $(\mathbf{8}_{v},\mathbf{496})$.
			\item Fermions: spacetime fermions need the right side in the R sector, and the diagonal boundary conditions then put all $\lambda^{A}$ in the R sector too, with $L_{0}\geq+1$. Level matching gives $\tilde{L}_{0}=L_{0}\geq1$, so $m^{2}\geq\frac{4}{\alpha'}$.
		\end{itemize}
	\end{remark}
\end{tcolorbox}

Now we examine various twists. First, examine the $\mathds{Z}_{2}$ generated by $\exp(\pi\mi \tilde{F})$. Combining with the diagonal projection (\ref{11.3.2}), this gives the total projection
\begin{equation}
	\frac{1+\exp[\pi\mi(F+\tilde{F})]}{2}\cdot \frac{1+\exp(\pi\mi\tilde{F})}{2}
	= \frac{1+\exp(\pi\mi F)}{2}\cdot \frac{1+\exp(\pi\mi\tilde{F})}{2} \:. \label{11.3.5}
\end{equation}
\begin{tcolorbox}[breakable]
	\begin{remark}
		Multiplying out the left side,
		\begin{align*}
			\bigl(1+\me^{\pi\mi(F+\tilde{F})}\bigr)\bigl(1+\me^{\pi\mi\tilde{F}}\bigr)
			&=1+\me^{\pi\mi\tilde{F}}+\me^{\pi\mi(F+\tilde{F})}+\me^{\pi\mi F}\bigl(\me^{\pi\mi\tilde{F}}\bigr)^{2}\\
			&\eqs{1}1+\me^{\pi\mi\tilde{F}}+\me^{\pi\mi F}\me^{\pi\mi\tilde{F}}+\me^{\pi\mi F}\\
			&=\bigl(1+\me^{\pi\mi F}\bigr)\bigl(1+\me^{\pi\mi\tilde{F}}\bigr)\:.
		\end{align*}
		At (1) $\me^{\pi\mi\tilde{F}}$ is a $\mathds{Z}_{2}$ operator, $(\me^{\pi\mi\tilde{F}})^{2}=1$, and $\me^{\pi\mi F}$, $\me^{\pi\mi\tilde{F}}$ commute because they act on different fields. The group generated by the diagonal element and $\me^{\pi\mi\tilde{F}}$ is the same $\mathds{Z}_{2}\times\mathds{Z}_{2}$ as the one generated by $\me^{\pi\mi F}$ and $\me^{\pi\mi\tilde{F}}$, and the projector onto invariant states depends only on the group.
	\end{remark}
\end{tcolorbox}
This is equivalent to the projection (\ref{11.2.8}) plus (\ref{11.2.10}), which define the supersymmetric $SO(32)$ heterotic string. Furthermore, in the sector added by the spatial twist of $\exp(\pi\mi\tilde{F})$, $\lambda^{A}$ and $\tilde{\psi}^{\mu}$ have opposite periodicity. Therefore, this twisted theory is the $SO(32)$ heterotic string. Twisting with $\exp(\pi\mi\tilde{F})$ has the same effect.

Now examine the diagonal theory twisted by $\exp(\pi\mi F_{1})$, where $\exp(\pi\mi F_{1})$ inverts the signs of the first 16 $\lambda^{A}$ and was used to construct the $E_{8}\times E_{8}$ heterotic string. The corresponding theory is not supersymmetric—just as in equation (\ref{11.2.8}), a theory will be supersymmetric if and only if the projection includes $\exp(\pi\mi \tilde{F})$. It has gauge group $E_{8}\times SO(16)$ and has a tachyon in $(\mathbf{1},\mathbf{16})$ (see the box below). Twisting once more with $\exp(\pi\mi\tilde{F})$ will yield the supersymmetric $E_{8}\times E_{8}$ heterotic string.
\begin{tcolorbox}[breakable]
	\begin{proof}
		Label the sectors by the boundary conditions of $(\tilde{\psi}^{\mu};\lambda^{1,\dots,16};\lambda^{17,\dots,32})$. The diagonal theory has $(\text{NS};\text{NS};\text{NS})$ and $(\text{R};\text{R};\text{R})$. The spatial twist by $\exp(\pi\mi F_{1})$ flips the periodicity of the first group, adding $(\text{NS};\text{R};\text{NS})$ and $(\text{R};\text{NS};\text{R})$. The projection is onto states invariant under both generators. Write $F=F_{1}+F_{1}'$:
		\begin{equation*}
			\exp(\pi\mi F_{1})=1\:,\qquad \exp[\pi\mi(F+\tilde{F})]=1\quad\Longleftrightarrow\quad \exp[\pi\mi(F_{1}'+\tilde{F})]=1\:.
		\end{equation*}
		Use the zero-point energies of the box after (\ref{11.2.14}), the rule that $\exp(\pi\mi\tilde{F})=-1$ on the right NS ground state and $\exp(\pi\mi F_{1})=\exp(\pi\mi F_{1}')=+1$ on the left NS ground state, and that each $\lambda_{-1/2}$ flips the corresponding sign.

		\textbf{Gauge bosons} are $(\text{left }L_{0}=0)\otimes\tilde{\psi}^{\mu}_{-1/2}\lvert0\rangle$, and the right factor has $\exp(\pi\mi\tilde{F})=+1$. Hence they need $\exp(\pi\mi F_{1})=\exp(\pi\mi F_{1}')=+1$ on the left.
		\begin{align*}
			(\text{NS};\text{NS};\text{NS}),\ L_{0}=-1+N:&\quad \lambda^{a}_{-1/2}\lambda^{b}_{-1/2}\ (a,b\leq16)\ \checkmark,\quad
			\lambda^{a'}_{-1/2}\lambda^{b'}_{-1/2}\ (a',b'\geq17)\ \checkmark,\\
			&\quad \lambda^{a}_{-1/2}\lambda^{b'}_{-1/2}\ \times\:,\\
			(\text{NS};\text{R};\text{NS}),\ L_{0}\eqs{1}0:&\quad \text{ground states }\mathbf{256}\ \xrightarrow{\ \exp(\pi\mi F_{1})=1\ }\ \mathbf{128}\ \text{of }SO(16)_{1}\:.
		\end{align*}
		At (1) $8(-\frac{1}{24})+16(\frac{1}{24})+16(-\frac{1}{48})=0$ as in (\ref{11.2.24}). The mixed states $\lambda^{a}\lambda^{b'}$ have $\exp(\pi\mi F_{1})=-1$ and are removed. The ground states in the second line have $\exp(\pi\mi F_{1}')=+1$ because the second group is in its NS ground state. The massless vectors are therefore
		\begin{equation*}
			(\mathbf{120}+\mathbf{128},\mathbf{1})+(\mathbf{1},\mathbf{120})=(\mathbf{248},\mathbf{1})+(\mathbf{1},\mathbf{120})\:,
		\end{equation*}
		the adjoint of $E_{8}\times SO(16)$, by the same $\mathbf{120}+\mathbf{128}=\mathbf{248}$ as in the $E_{8}\times E_{8}$ theory. The sector $(\text{R};\text{NS};\text{R})$ has a right-moving R state, so it contributes fermions, not vectors.

		\textbf{Tachyons} need $\tilde{L}_{0}=-\frac{1}{2}$, i.e.\ the right NS ground state with $\exp(\pi\mi\tilde{F})=-1$. They therefore need $L_{0}=-\frac{1}{2}$ with $\exp(\pi\mi F_{1})=+1$ and $\exp(\pi\mi F_{1}')=-1$. The $(\text{NS};\text{R};\text{NS})$ sector has $L_{0}\geq0$, so only $(\text{NS};\text{NS};\text{NS})$ with one $\lambda_{-1/2}$ can contribute:
		\begin{align*}
			\lambda^{a}_{-1/2}\lvert0\rangle\ (a\leq16)&:\quad \exp(\pi\mi F_{1})=-1\ \times\:,\\
			\lambda^{a'}_{-1/2}\lvert0\rangle\ (a'\geq17)&:\quad \exp(\pi\mi F_{1})=+1,\ \exp(\pi\mi F_{1}')=-1\ \checkmark\:.
		\end{align*}
		The surviving tachyon is the vector of the second $SO(16)$ and a singlet of $E_{8}$, i.e.\ it lies in $(\mathbf{1},\mathbf{16})$. It has $m^{2}=-2/\alpha'$.
	\end{proof}
\end{tcolorbox}

This step can be continued by dividing $\lambda^{A}$ into groups of 8, 4, 2, and 1 each. Write the $SO(32)$ index $A$ in binary, $A=1+d_{1}d_{2}d_{3}d_{4}d_{5}$, where each digit $d_{i}$ is 0 or 1. For $i=1,\ldots,5$, define the operator $\exp(\pi\mi F_{i})$ such that it anticommutes with those $\lambda_{A}$ for which $d_{i}=0$ and commutes with those $\lambda_{A}$ for which $d_{i}=1$. Clearly, there are five possible twist groups, with 2, 4, 8, 16, and 32 elements respectively, generated by choosing 1, 2, 3, 4, or 5 elements from $\exp(\pi\mi F_{i})$ and constructing various products. The first gives the $E_{8}\times SO(16)$ theory described earlier; further twisting produces gauge groups $SO(24)\times SO(8)$, $E_{7}\times E_{7}\times SO(4)$, $SU(16)\times SO(2)$, and $E_{8}$. None of these theories are supersymmetric, and all have tachyons. Further twisting with $\exp(\pi\mi\tilde{F})$ gives supersymmetric theories, which are either $SO(32)$ or $E_{8}\times E_{8}$ theories. In this construction, the gauge group is non-trivial, with more currents coming from the R sector.

Let's review the logic of the twist construction. For a sector twisted by group element $h$, the corresponding vertex operator will produce a branch cut in the fields, but the projection onto $h$-invariant states means these branch cuts will not appear in the products of the vertex operators. Since $h$ is a symmetry, the projection is protected by interaction. On the torus, the sum of timelike and spacelike twists is modularly invariant, and this can be generalized to any genus. However, we learned in section 10.7 that due to the anomalous phase in modular transformation, this assumed modular invariance of summing over path integral boundaries is not enough. Only for right-left symmetric path integrals will the phase automatically cancel.

At one loop, anomalous phases only appear in the transformation $\tau\to\tau+1$, and they are interpreted as the level matching condition $L_{0}-\tilde{L}_{0}\in\mathds{Z}$ not being satisfied. There is further a theorem that for an Abelian twist group (like the product of $\mathds{Z}_{2}$ examined here), if each twisted sector has infinitely many level-matched states before attaching the projection, then the one-loop amplitude, and indeed all amplitudes, are precisely modularly invariant. In a heterotic string, take such a sector: $k$ $\lambda^{A}$ satisfy R boundary conditions, and $32-k$ $\lambda^{A}$ satisfy NS boundary conditions; the zero-point energy (note: the right-moving part is zero; only left-moving is considered here) is
\begin{equation}
	{-}\frac{8}{24}+\frac{k}{24}-\frac{(32-k)}{48}=-1+\frac{k}{16} \:. \label{11.3.6}
\end{equation}
\begin{tcolorbox}[breakable]
	\begin{remark}
		Each R fermion contributes $+\frac{1}{24}$ and each NS fermion $-\frac{1}{48}$ (box after (\ref{11.2.14})), so
		\begin{equation*}
			-\frac{8}{24}+\frac{k}{24}-\frac{32-k}{48}\eqs{1}-\frac{16}{48}-\frac{32}{48}+\frac{2k+k}{48}=-1+\frac{k}{16}\:.
		\end{equation*}
		At (1) everything was put over $48$. Each unit of R-sector fermion changes the vacuum energy by $\frac{1}{24}+\frac{1}{48}=\frac{1}{16}$, which is the familiar $\frac{1}{16}$ weight of a single-fermion twist field. NS oscillators raise $L_{0}$ in steps of $\frac{1}{2}$ and R oscillators in integer steps, so $L_{0}\in\frac{k}{16}+\frac{1}{2}\mathds{Z}$. With $\tilde{L}_{0}\in\frac{1}{2}\mathds{Z}$ on the right, level matching $L_{0}-\tilde{L}_{0}\in\mathds{Z}$ can be satisfied by infinitely many states (before projection) iff $\frac{k}{16}\in\frac{1}{2}\mathds{Z}$, i.e.\ $k\in8\mathds{Z}$.
	\end{remark}
\end{tcolorbox}
The oscillator will raise the energy by an integer multiple of $\frac{1}{2}$, so the energy of the left-moving part is $\frac{1}{16}k\operatorname{mod}\frac{1}{2}$. On the right-moving side, for Lorentz invariance, we still take common boundary conditions for all fermions, so the energy is an integer multiple of $\frac{1}{2}$. Therefore, if $k$ is an integer multiple of 8, the level matching condition is precisely satisfied. And as we have seen, OPE closure and spacetime spin-statistics actually require $k$ to be an integer multiple of 16. Defining twists $\exp(\pi\mi F_{i})$ such that any product of them happens to anticommute with 16 $\lambda^{A}$ satisfies this condition.


When the level matching condition is satisfied, multiple modularly invariant and compatible theories can actually exist. Consider a twisted theory with a partition function
\begin{equation}
	Z=\frac{1}{\operatorname{order}(H)}\sum_{h_{1},h_{2}\in H}Z_{h_{1},h_{2}} \:, \label{11.3.7}
\end{equation}
where $h_{1}$ and $h_{2}$ are the timelike and spacelike periods on the torus. Then, the partition function
\begin{equation}
	Z=\frac{1}{\operatorname{order}(H)}\sum_{h_{1},h_{2}\in H}\epsilon(h_{1},h_{2}) Z_{h_{1},h_{2}} \label{11.3.8}
\end{equation}
is also compatible (modularly invariant with closed and local OPE), where the phase $\epsilon(h_{1},h_{2})$ satisfies
\begin{subequations}
	\begin{align}
		\epsilon(h_{1},h_{2}) &= \epsilon(h_{2},h_{1})^{-1} \:, \label{11.3.9a} \\
		\epsilon(h_{1},h_{2})\epsilon(h_{1},h_{3}) &= \epsilon(h_{1},h_{2}h_{3}) \:,\label{11.3.9b} \\
		\epsilon(h,h) &= 1 \:. \label{11.3.9c}
	\end{align} \label{11.3.9}
\end{subequations}
Written as $\hat{h}_{2}$ in the original twisted theory, the new theory no longer projects onto $H$-invariant states, but rather onto states satisfying
\begin{equation}
	\hat{h}_{2}\lvert\psi\rangle_{h_{1}} = \epsilon(h_{1},h_{2})^{-1}\lvert \psi\rangle_{h_{1}} \label{11.3.10}    
\end{equation}
In other words, the state is now an eigenvector of $\hat{h}$, and the phase is sector-related; equivalently, we have a sector-related redefinition
\begin{equation}
	\hat{h}\to \epsilon(h_{1},h)\hat{h} \:. \label{11.3.11}
\end{equation}
The phase factor $\epsilon(h_{1},h_{2})$ is called \emph{discrete torsion}.

In the previous theory, namely the group generated by $\exp(\pi\mi F_{1})$ and $\exp{\pi\mi \tilde{F}}$, which generated the $E_{8}\times E_{8}$ string from the diagonal theory, there is an interesting possibility for discrete torsion. For
\begin{equation}
	(h_{1},h_{2})= \Bigl(\exp[\pi\mi(k_{1}F_{1}+l_{1}\tilde{F})],
	\exp[\pi\mi(k_{2}F_{1}+l_{2}\tilde{F})]\Bigr) \label{11.3.12} 
\end{equation}
the phase
\begin{equation}
	\epsilon(h_{1},h_{2})= (-1)^{k_{1}l_{2}+k_{2}l_{1}} \label{11.3.13}
\end{equation}
satisfies condition (\ref{11.3.9}). It modifies the projection that produces the supersymmetric $E_{8}\times E_{8}$ string
\begin{equation}
	\exp(\pi\mi F_{1})=\exp(\pi\mi F_{1}^{\prime}) = \exp(\pi\mi\tilde{F})=1\:, \label{11.3.14}
\end{equation}
to
\begin{equation}
	\exp[\pi\mi (F_{1}+\alpha_{1}^{\prime}+\tilde{\alpha})]
	=\exp[\pi\mi(F_{1}^{\prime}+\alpha_{1}+\tilde{\alpha})] 
	= \exp[\pi\mi(\tilde{F}+\alpha_{1}+\alpha_{1}^{\prime})]=1\:, \label{11.3.15}
\end{equation}
the sign conventions here being similar to those in equation (\ref{10.7.11}): under the transformation $w\to w+2\pi$, $\tilde{\psi}^{\mu}$, the first 16 $\lambda^{A}$, and the second 16 $\lambda^{A}$ pick up phases $-\exp(-\pi\mi\tilde{\alpha})$, $-\exp(-\pi\mi\alpha_{1})$, and $-\exp(-\pi\mi\alpha_{1}^{\prime})$, respectively. These $\alpha$ are conserved under OPE, so this projection is self-consistent. In other words, the spectrum is composed of the following sectors:
\begin{align*}
	&(\mathrm{NS}+,\mathrm{NS}+,\mathrm{NS}+) \:, \\
	&(\mathrm{NS}-,\mathrm{NS}-,\mathrm{R}+)\:,\quad (\mathrm{NS}-,\mathrm{R}+,\mathrm{NS}-)\:,\quad
	(\mathrm{NS}+,\mathrm{R}-,\mathrm{R}-)\:, \\
	&(\mathrm{R}+,\mathrm{NS}-,\mathrm{NS}-)\:,\quad (\mathrm{R}-,\mathrm{R}-,\mathrm{NS}+)\:, \quad
	(\mathrm{R}-,\mathrm{NS}+,\mathrm{R}-) \:, \\
	&(\mathrm{R}+,\mathrm{R}+,\mathrm{R}+) 
\end{align*}
where the three signs refer to $\tilde{\psi}^{\mu}$, the first 16 $\lambda^{A}$, and the second 16 $\lambda^{A}$, respectively.
\begin{tcolorbox}[breakable]
	\begin{remark}
		\textbf{The phase (\ref{11.3.13}) is a discrete torsion.} The group is $\mathds{Z}_{2}\times\mathds{Z}_{2}$ with elements labeled by $(k,l)\in\mathds{Z}_{2}^{2}$, and group multiplication adds labels mod 2. Then
		\begin{align*}
			\epsilon(h_{2},h_{1})^{-1}&\eqs{1}(-1)^{k_{2}l_{1}+k_{1}l_{2}}=\epsilon(h_{1},h_{2})\:,\\
			\epsilon(h_{1},h_{2})\epsilon(h_{1},h_{3})&=(-1)^{k_{1}(l_{2}+l_{3})+(k_{2}+k_{3})l_{1}}\eqs{2}\epsilon(h_{1},h_{2}h_{3})\:,\\
			\epsilon(h,h)&=(-1)^{2kl}=1\:.
		\end{align*}
		At (1) $\epsilon=\pm1$ is its own inverse and the exponent is symmetric under $1\leftrightarrow2$. At (2) $h_{2}h_{3}$ has labels $(k_{2}+k_{3},l_{2}+l_{3})$, and the sign depends only on the labels mod 2. So (\ref{11.3.9a})--(\ref{11.3.9c}) hold.

		\textbf{The sector list.} The spatial twists produce all $2^{3}$ combinations of periodicities, labeled by $(\tilde{\alpha},\alpha_{1},\alpha_{1}')\in\{0,1\}^{3}$ with $0=$ NS and $1=$ R. For each, (\ref{11.3.15}) fixes the three eigenvalues:
		\begin{equation*}
			\me^{\pi\mi\tilde{F}}=(-1)^{\alpha_{1}+\alpha_{1}'}\:,\qquad \me^{\pi\mi F_{1}}=(-1)^{\alpha_{1}'+\tilde{\alpha}}\:,\qquad \me^{\pi\mi F_{1}'}=(-1)^{\alpha_{1}+\tilde{\alpha}}\:.
		\end{equation*}
		For example, $(\tilde{\alpha},\alpha_{1},\alpha_{1}')=(0,0,1)$ gives $(-,-,+)$, i.e.\ $(\text{NS}-,\text{NS}-,\text{R}+)$, and $(1,0,0)$ gives $(+,-,-)$, i.e.\ $(\text{R}+,\text{NS}-,\text{NS}-)$. Going through the eight cases reproduces the list above. In particular $(\text{NS},\text{NS},\text{NS})$ is projected onto $(+,+,+)$, which excludes the tachyon (it needs $\me^{\pi\mi\tilde{F}}=-1$), and $(\text{R},\text{NS},\text{NS})$ onto $(+,-,-)$. Without torsion the latter would have been $(+,+,+)$, which contains the gravitino.
	\end{remark}
\end{tcolorbox}

The gravitino is in the $(\mathrm{R}\pm,\mathrm{NS}+,\mathrm{NS}+)$ sectors, so it is not in this spectrum. The tachyon is in the $(\mathrm{NS}-,\mathrm{NS}-,\mathrm{NS}+)$ and $(\mathrm{NS}-,\mathrm{NS}+,\mathrm{NS}-)$ sectors, so it is also not in this spectrum. This twist leaves $SO(16)\times SO(16)$ gauge symmetry. Classifying the states according to their $SO(8)_{\text{spin}}\times SO(16)\times SO(16)$ quantum numbers, we find the following massless spectrum:
\begin{align*}
	(\mathrm{NS}+,\mathrm{NS}+,\mathrm{NS}+)&\::\quad (\mathbf{1},\mathbf{1},\mathbf{1}) 
	+ (\mathbf{28},\mathbf{1},\mathbf{1}) + (\mathbf{35},\mathbf{1},\mathbf{1})
	+ (\mathbf{8}_{v},\mathbf{120},\mathbf{1}) + (\mathbf{8}_{v},\mathbf{1},\mathbf{120}) \:, \\
	(\mathrm{R}+,\mathrm{NS}-,\mathrm{NS}-)&\::\quad (\mathbf{8},\mathbf{16},\mathbf{16})\: , \\
	(\mathrm{R}-,\mathrm{R}-,\mathrm{NS}+)&\::\quad  (\mathbf{8}^{\prime},\mathbf{128}^{\prime},\mathbf{1})\:, \\
	(\mathrm{R}-,\mathrm{NS}+,\mathrm{R}-)&\::\quad (\mathbf{8}^{\prime},\mathbf{1},\mathbf{128}^{\prime})\:.
\end{align*}
This shows that theories without supersymmetry can possibly contain no tachyons.
\begin{tcolorbox}[breakable]
	\begin{remark}
		Take the left zero-point energies $-1$ (NS,NS), $0$ (one group R) and $+1$ (both groups R), and on the right $-\frac{1}{2}$ (NS) and $0$ (R). Massless states need $L_{0}=\tilde{L}_{0}=0$. On the right that means $\tilde{\psi}_{-1/2}\lvert0\rangle$ ($\mathbf{8}_{v}$, $\me^{\pi\mi\tilde{F}}=+$) in NS, or the R ground states, with $\mathbf{8}$ for $\me^{\pi\mi\tilde{F}}=+$ and $\mathbf{8}'$ for $-$. Going through the sectors:
		\begin{align*}
			(\text{NS}+,\text{NS}+,\text{NS}+):&\quad \bigl[\alpha^{i}_{-1},\ \lambda^{a}_{-1/2}\lambda^{b}_{-1/2},\ \lambda^{a'}_{-1/2}\lambda^{b'}_{-1/2}\bigr]\otimes\mathbf{8}_{v}\\
			&\quad=(\mathbf{8}_{v}\times\mathbf{8}_{v},\mathbf{1},\mathbf{1})+(\mathbf{8}_{v},\mathbf{120},\mathbf{1})+(\mathbf{8}_{v},\mathbf{1},\mathbf{120})\:,\\
			(\text{R}+,\text{NS}-,\text{NS}-):&\quad \lambda^{a}_{-1/2}\lambda^{b'}_{-1/2}\lvert0\rangle\otimes\mathbf{8}=(\mathbf{8},\mathbf{16},\mathbf{16})\:,\\
			(\text{R}-,\text{R}-,\text{NS}+):&\quad (\text{R ground states of group 1 with }\me^{\pi\mi F_{1}}=-)\otimes\mathbf{8}'\\
			&\quad=(\mathbf{8}',\mathbf{128}',\mathbf{1})\:,\\
			(\text{R}-,\text{NS}+,\text{R}-):&\quad (\mathbf{8}',\mathbf{1},\mathbf{128}')\ \text{likewise}\:.
		\end{align*}
		In the first line $\mathbf{8}_{v}\times\mathbf{8}_{v}=\mathbf{1}+\mathbf{28}+\mathbf{35}$. The mixed bilinear $\lambda^{a}\lambda^{b'}$ has $(\me^{\pi\mi F_{1}},\me^{\pi\mi F_{1}'})=(-,-)$, so it is excluded there but is exactly what the second sector needs. In the third line the first group's zero-point energy is $0$ and the second group is in its NS ground state, which has $\me^{\pi\mi F_{1}'}=+$. The remaining four sectors give nothing massless:
		\begin{itemize}
			\item $(\text{NS}-,\text{NS}-,\text{R}+)$ and $(\text{NS}-,\text{R}+,\text{NS}-)$: the left side has $L_{0}\geq0$. On the right, $\me^{\pi\mi\tilde{F}}=-$ allows only even numbers of $\tilde{\psi}_{-1/2}$ on the NS ground state, so $\tilde{L}_{0}\in\{-\frac{1}{2},\frac{1}{2},\dots\}$ and never equals $0$.
			\item $(\text{NS}+,\text{R}-,\text{R}-)$ and $(\text{R}+,\text{R}+,\text{R}+)$: both groups are R, so $L_{0}\geq1$.
		\end{itemize}
		The tachyon would need $\tilde{L}_{0}=-\frac{1}{2}$, i.e.\ the right NS ground state with $\me^{\pi\mi\tilde{F}}=-$, together with $L_{0}=-\frac{1}{2}$, i.e.\ both $\lambda$ groups NS. Neither of the two sectors with that right side, $(\text{NS}-,\text{NS}-,\text{R}+)$ and $(\text{NS}-,\text{R}+,\text{NS}-)$, has both groups NS. The theory is therefore tachyon free, and the only massless spin-$\frac{3}{2}$ candidate $(\text{R}\pm,\text{NS}+,\text{NS}+)$ is absent, so there is no gravitino.
	\end{remark}
\end{tcolorbox}


\section{A Little Lie Algebra} \label{sec:11.4}

\subsection*{Basic definitions}

A Lie algebra is a vector space equipped with an anticommuting product $[T,T^{\prime}]$. In the form of basis $T^{a}$, this product is
\begin{equation}
	[T^{a},T^{b}] = \mi f\indices{^a^b_c}T^{c} \:, \label{11.4.1}
\end{equation}
where $f\indices{^a^b_c}$ are the \emph{structure constants}. This product is required to satisfy the Jacobi identity
\begin{equation}
	[T^{a},[T^{b},T^{c}]] + [T^{b},[T^{c},T^{a}]] + [T^{c},[T^{a},T^{b}]] =0 \:.\label{11.4.2}
\end{equation}
The corresponding Lie group is generated by the exponent
\begin{equation}
	\exp(\mi\theta_{a}T^{a}) \:, \label{11.4.3}
\end{equation}
where $\theta_{a}$ are real numbers. For a compact Lie group, the corresponding compact Lie algebra possesses a positive-definite inner product
\begin{equation}
	(T^{a},T^{b}) = d^{ab} \:, \label{11.4.4}
\end{equation}
which is invariant,
\begin{equation}
	([T,T^{\prime}],T^{\prime\prime}) + (T^{\prime},[T,T^{\prime\prime}]) =0\:. \label{11.4.5}
\end{equation}
The equivalent statement of invariance is that $f^{abc}$ is fully antisymmetric, and indices are raised or lowered with $d^{ab}$.
\begin{tcolorbox}[breakable]
	\begin{proof}
		Take $T=T^{b}$, $T'=T^{a}$, $T''=T^{c}$ in (\ref{11.4.5}) and define $f^{abc}=f\indices{^a^b_d}d^{dc}$. Then
		\begin{align*}
			0=([T^{b},T^{a}],T^{c})+(T^{a},[T^{b},T^{c}])
			\eqs{1}\mi f\indices{^b^a_d}d^{dc}+\mi f\indices{^b^c_d}d^{ad}
			\eqs{2}\mi\bigl(f^{bac}+f^{bca}\bigr)\:.
		\end{align*}
		At (1) (\ref{11.4.1}) and (\ref{11.4.4}) were used. At (2) $d^{ad}=d^{da}$. So $f^{abc}$ is antisymmetric in its last two indices. It is antisymmetric in the first two by (\ref{11.4.1}), and the two transpositions generate all permutations, so it is totally antisymmetric. Reading the steps backwards, total antisymmetry implies (\ref{11.4.5}) for basis elements, and by linearity for all elements.
	\end{proof}
\end{tcolorbox}

We are interested in simple Lie algebras, which have no non-trivial invariant subalgebras (ideals). General compact algebras are the direct sum of simple algebras and $U(1)$. For a simple algebra, the inner product is unique up to a normalization, and there exists a set of bases for the generators such that it is $\delta^{ab}$. For any representation $r$ of a Lie algebra (any set of matrices $t^{a}_{r,ij}$ satisfying (\ref{11.4.1}) as given by $f\indices{^a^b_c}$), its basis is invariant, so for a simple Lie algebra, it is proportional to $d^{ab}$,
\begin{equation}
	\operatorname{Tr}(t_{r}^{a}t_{r}^{b}) = T_{r}d^{ab} \label{11.4.6}
\end{equation}
where $T_{r}$ is some constant. Additionally, $t_{r}^{a}t_{r}^{b}d_{ab}$ commutes with all $t_{r}^{c}$, so it is proportional to the identity operator,
\begin{equation}
	t_{r}^{a}t_{r}^{b}d_{ab} = Q_{r}  \label{11.4.7}
\end{equation}
$Q_{r}$ is the \emph{Casimir invariant} of representation $r$.
\begin{tcolorbox}[breakable]
	First prove that $t_{r}^{a}t_{r}^{b}d_{ab}$ commutes with all $t_{r}^{c}$:
	\begin{align*}
		[d_{ab}t_{r}^{a}t_{r}^{b},t_{r}^{c}] &= d_{ab}[t_{r}^{a},t_{r}^{c}]t_{r}^{b} 
		+ d_{ab}t_{r}^{a}[t_{r}^{b},t_{r}^{c}] 
		=\mi \, d_{ab} (f\indices{^a^c_d}t_{r}^{d}t_{r}^{b}+ f\indices{^b^c_d}t_{r}^{a}t_{r}^{d}) \\
		&=\mi \, (f\indices{_b^c_d}t_{r}^{d}t_{r}^{b}+ f\indices{_a^c_d}t_{r}^{a}t_{r}^{d}) 
		=\mi \, (f\indices{^b^c^d}t_{r\,d}t_{r\,b}+ f\indices{^a^c^d}t_{r\,a}t_{r\,d}) \\
		&=\mi \, (f\indices{^b^c^d}t_{r\,d}t_{r\,b}+ f\indices{^d^c^b}t_{r\,d}t_{r\,b}) =0
	\end{align*}
	Then, according to Schur's lemma, (\ref{11.4.7}) is obtained.
\end{tcolorbox}

Typical Lie algebras include:
\begin{itemize}
	\item $SU(n)$: Generators are traceless Hermitian $n\times n$ matrices. The corresponding group consists of unitary matrices with determinant 1. This algebra is also denoted as $A_{n-1}$.
	\item $SO(n)$: Generators are antisymmetric Hermitian $n\times n$ matrices. The corresponding group $SO(n,\mathds{R})$ consists of real orthogonal matrices with determinant 1. When $n=2k$, this algebra is also denoted as $D_{k}$. When $n=2k+1$, it is denoted as $B_{k}$.
	\item $Sp(k)$: Generators are $2k\times 2k$ Hermitian matrices, but possess extra structure
	\begin{equation}
		M T M^{-1} = - T^{\mathrm{T}} \:. \label{11.4.8}
	\end{equation}
	where superscript $\mathrm{T}$ denotes transpose, and
	\begin{equation}
		M = \mi \begin{bmatrix}
			0 & I_{k} \\ -I_{k} & 0 
		\end{bmatrix} \label{11.4.9}
	\end{equation}
	where $I_{k}$ is the $k\times k$ identity matrix. The corresponding group also consists of unitary matrices $U$, but satisfies extra properties
	\begin{equation}
		M U M^{-1} = (U^{\mathrm{T}})^{-1} \:. \label{11.4.10}
	\end{equation}
	Sometimes this group is also denoted as $Sp(2k)$. It is simultaneously denoted as $C_{k}$.
	\begin{tcolorbox}[breakable]
		\begin{remark}
			For $U=\exp(\mi\theta_{a}T^{a})$ with each $T^{a}$ obeying (\ref{11.4.8}),
			\begin{equation*}
				MUM^{-1}\eqs{1}\exp\bigl(\mi\theta_{a}MT^{a}M^{-1}\bigr)\eqs{2}\exp\bigl(-\mi\theta_{a}T^{a\mathrm{T}}\bigr)\eqs{3}\bigl(U^{\mathrm{T}}\bigr)^{-1}\:.
			\end{equation*}
			At (1) conjugation acts term by term on the power series. At (2) (\ref{11.4.8}) was used. At (3) $\exp(A^{\mathrm{T}})=(\exp A)^{\mathrm{T}}$ and $\exp(-A)=(\exp A)^{-1}$. Equivalently, $U^{\mathrm{T}}MU=M$, so $U$ preserves the antisymmetric form $M$. For the dimension, note $M^{-1}=M$ and write $T$ in $k\times k$ blocks:
			\begin{gather*}
				T=\begin{bmatrix}A&B\\C&D\end{bmatrix}\:,\qquad
				MTM^{-1}=\begin{bmatrix}D&-C\\-B&A\end{bmatrix}\eqs{4}-\begin{bmatrix}A^{\mathrm{T}}&C^{\mathrm{T}}\\B^{\mathrm{T}}&D^{\mathrm{T}}\end{bmatrix}\\
				\Longrightarrow\quad D=-A^{\mathrm{T}},\ B=B^{\mathrm{T}},\ C=C^{\mathrm{T}}\:.
			\end{gather*}
			At (4) the condition (\ref{11.4.8}) was imposed. So $A$ is free ($k^{2}$ parameters) and $B,C$ are symmetric ($\frac{1}{2}k(k+1)$ each). Hermiticity then relates $C=B^{\dagger}$ and makes $A$ Hermitian, leaving $\operatorname{dim}Sp(k)=k^{2}+2\cdot\frac{1}{2}k(k+1)=2k^{2}+k$ real parameters.
		\end{remark}
	\end{tcolorbox}
\end{itemize}

For each compact group, by multiplying some of its generators by $\mi$, we can obtain various non-compact groups. For example, traceless \emph{imaginary} matrices generate the $SL(n,\mathds{R})$ group, which is the group of real matrices with determinant 1. The group $SO(m,n,\mathds{R})$ that preserves Lorentz-type inner products can be similarly obtained from $SO(m+n)$. Generators of another non-compact group also satisfy the symplectic condition (\ref{11.4.8}), but the generators are imaginary matrices rather than Hermitian matrices; this group consists of real matrices satisfying (\ref{11.4.10}). This non-compact group is also denoted as $Sp(k)$ or $Sp(2k)$; sometimes $USp(2k)$ is used to distinguish the compact unitary case.

These non-compact groups do not appear in Yang-Mills theories (the corresponding results would not be unitary), but they have other applications. In compactified string theories, some $SL(n,\mathds{R})$ and $SO(m,n,\mathds{R})$ appear as low-energy symmetries. And in classical mechanics, the real form of $Sp(k)$ appears as an invariant of the Poisson bracket.

\subsection*{Roots and weights}

For any Lie algebra $h$, choose the largest set of mutually commuting generators $H^{i}$, $i=1,\cdots,\operatorname{rank}(g)$. The remaining generators $E^{\alpha}$ can be chosen to have definite charges under $H^{i}$,
\begin{equation}
	[H^{i},E^{\alpha}]=\alpha^{i}E^{\alpha} \:. \label{11.4.11}
\end{equation}
The $\operatorname{rank}(g)$-dimensional vector $\alpha^{i}$ is called a \emph{root}. There is a theorem that for a given root there is only one generator, so the notation $E^{\alpha}$ is unambiguous. The Jacobi identity determines the remaining algebra as
\begin{equation}
	[E^{\alpha},E^{\beta}] = 
	\begin{cases}
		\epsilon(\alpha,\beta) E^{\alpha+\beta} &\qquad \text{if }\alpha+\beta \text{ is a root} \:, \\
		2\alpha\cdot H/\alpha^{2} &\qquad \text{if }\alpha+\beta=0\:, \\
		0 &\qquad \text{otherwise} \:.
	\end{cases} \label{11.4.12}
\end{equation}
The products $\alpha\cdot H$ and $\alpha^{2}$ are defined using contraction with $d_{ij}$, which is the inverse of the inner product (\ref{11.4.4}) restricted to the commutator algebra. Taking the trace with $H^{i}$, the second equation determines the normalization $(E^{\alpha},E^{-\alpha})=2/\alpha^{2}$. The constant $\epsilon(\alpha,\beta)$ can only take the value $\pm1$.
\begin{tcolorbox}[breakable]
	Now prove (\ref{11.4.12}) using the Jacobi identity:
	\begin{align*}
		[H^{i},[E^{\alpha},E^{\beta}]] &= [E^{\alpha},[H^{i},E^{\beta}]] + [E^{\beta},[E^{\alpha},H^{i}]]
		= \beta^{i}[E^{\alpha},E^{\beta}] - \alpha^{i} [E^{\beta},E^{\alpha}] \\
		&= (\alpha^{i}+\beta^{i})[E^{\alpha},E^{\beta}] \:,
	\end{align*}
	This shows that $[E^{\alpha},E^{\beta}]$ is either an eigenvector of $H^{i}$ or zero, which is case 3. If $\alpha+\beta$ is not zero, then it is a root, and $[E^{\alpha},E^{\beta}]$ is the corresponding $E^{\alpha+\beta}$, where a normalization coefficient remains to be determined. If $\alpha+\beta$ is zero, it indicates that $[E^{\alpha},E^{\beta}]$ commutes with $H^{i}$, then $[E^{\alpha},E^{\beta}]$ can only be a linear combination of $H^{i}$. Thus, let $[E^{\alpha},E^{-\alpha}]= a_{j}H^{j}$; then according to (\ref{11.4.5}),
	\begin{align*}
		(H^{i},[E^{\alpha},E^{-\alpha}])&= -([E^{\alpha},H^{i}],E^{-\alpha}) \\
		a^{i} & = \alpha^{i}(E^{\alpha},E^{-\alpha}) 
	\end{align*}
	Normalizing $(E^{\alpha},E^{-\alpha})$ to $2/\alpha^{2}$ yields the desired result. The result of the first term is arbitrary; here the coefficient is taken as $\pm1$.
\end{tcolorbox}

The matrices $t_{r}^{i}$ representing $H^{i}$ can be taken to be diagonal. Their common eigenvalues $w^{i}$ combine into a vector
\begin{equation}
	(w^{1},\cdots, w^{rank(g)}) \:, \label{11.4.13}
\end{equation}
This vector is called a weight, and its number is equal to the dimension of the representation. In the adjoint representation, roots and weights are the same.
\begin{tcolorbox}[breakable]
	Since $H^{i}$ mutually commute, they can be simultaneously diagonalized, possessing a common eigenvector $\lvert\lambda\rangle$; $\lvert\lambda\rangle$ spans the representation space, and the number of eigenvectors is the dimension of the representation. For the adjoint representation, $\lvert\lambda\rangle$ is $E^{\alpha}$.
\end{tcolorbox}

\noindent Examples:
\begin{itemize}
	\item For $SO(2k)=D_{k}$, examine $k$ $2\times2$ blocks along the diagonal, let $H^{i}$ be the $i$-th block for
	\begin{equation}
		\begin{bmatrix}
			0 & \mi \\ -\mi & 0
		\end{bmatrix}  \label{11.4.14}
	\end{equation}
	with zeros elsewhere. This is the maximal mutually commuting set. Under these $k$ $H^{i}$, the $2k$-vector $(1,\pm\mi,0,\cdots,0)$ possesses eigenvalues
	\begin{equation}
		(\pm 1, 0^{k-1}) \:; \label{11.4.15}
	\end{equation}
	they are weights of this vector representation. Other weights just shift $\pm1$ to other positions.
	
	The adjoint representation is an antisymmetric tensor, which is contained in the product of two vector representations. Weights are additive, so by adding any two different (due to antisymmetry) vector weights together, the root is obtained. This gives
	\begin{equation}
		(+1,+1,0^{k-2})\:,\quad (+1,-1,0^{k-2})\:,\quad (-1,-1,0^{k-2})\:, \label{11.4.16}
	\end{equation}
	and all their permutations. By adding any weight and its opposite weight together, we obtain $k$ zero roots, which are $H^{i}$.
	\begin{tcolorbox}[breakable]
		Examine the antisymmetric product $\lvert\lambda_{1}\rangle\wedge \lvert\lambda_{2}\rangle$ of 2 $2k$-vectors $\lvert \lambda_{1,2}\rangle$, and assume $\lvert\lambda_{1,2}\rangle$ only have eigenvalues $+1$ under $H^{1,2}$ respectively, then
		\begin{align*}
			&H^{1}\lvert \lambda_{1} \rangle \wedge \lvert \lambda_{2} \rangle = 1\\
			&H^{2} \lvert \lambda_{1} \rangle \wedge \lvert \lambda_{2} \rangle =1
		\end{align*}
		From this, (\ref{11.4.16}) is obtained
	\end{tcolorbox}
	In the spinor representation, weight $w^{i}$ is given by all $k$-vectors with all components being $\frac{1}{2}$, where $\mathbf{2^{k-1}}$ has an even number of $-\frac{1}{2}$, while $\mathbf{2^{k-1}}^{\prime}$ has an odd number.
	\item For $SO(2k+1)=B_{k}$, the same diagonal generators can be taken, only the last column and last row are zero. The weights in the vector representation are basically the same as above, but with one more $(0^{k})$ from the last row. The additional generators have roots
	\begin{equation}
		(\pm1,0^{k-1}) \label{11.4.17}
	\end{equation}
	and all permutations.
	\item The adjoint of $Sp(k)=C_{k}$ is a \emph{symmetric} tensor, so the same roots as $SO(2k)$ will be obtained, the difference being that vector weights need not be distinct. The corresponding roots are the roots of $SO(2k)$ and
	\begin{equation}
		(\pm2,0^{k-1}) \label{11.4.18}
	\end{equation}
	along with permutations. Usually this root is normalized so that the length of the longest root is $\sqrt{2}$, so all roots are divided by $2^{1/2}$.
	\item For $SU(n)=A_{n-1}$, it will be more convenient to discuss $U(n)$ first, although its algebra is not a simple algebra. $n$ mutually commuting generators $H^{i}$ can be taken to be 1 at position $ii$ and zero elsewhere. Then the charged generator with 1 at position $ij$ has an eigenvalue of $+1$ for $H^{i}$ and an eigenvalue of $-1$ for $H^{j}$: the roots are
	\begin{equation}
		(+1,-1,0^{n-2})  \label{11.4.19}
	\end{equation}
	of all permutations. Note that all roots reside on the hyperplane $\sum_{i} a^{i}=0$; this is because all eigenvalues of the $U(1)$ generator are zero. The roots of $SU(n)$ are those in $U(n)$ considered to reside on this hyperplane.
	\item We said the $E_{8}$ generator decomposes into the adjoint representation and spinor representation of $SO(16)$. The commuting generators of $SO(16)$ can also be taken as the commuting generators of $E_{8}$, so the roots of $E_{8}$ are the roots of $SO(16)$ plus the roots of the spinor, i.e., all permutations of roots (\ref{11.4.16}) plus
	\begin{equation}
		(+\tfrac{1}{2},+\tfrac{1}{2},+\tfrac{1}{2},+\tfrac{1}{2},+\tfrac{1}{2},
		+\tfrac{1}{2},+\tfrac{1}{2},+\tfrac{1}{2}) \label{11.4.20}
	\end{equation}
	and roots obtained by doing an even number of sign flips on this root. Equivalently, this set can be described as
	\begin{subequations}
		\begin{align}
			&\text{for all } i,\qquad \alpha^{i}\in \mathds{Z} \text{ or } \alpha^{i}\in\mathds{Z}+\tfrac{1}{2}\:,
			\label{11.4.21a} \\
			\noalign{\hbox{and}}
			&\sum_{i=1}^{8}\alpha^{i}\in 2\mathds{Z}\:,\qquad \sum_{i=1}^{8}(\alpha^{i})^{2}=2 \:. \label{11.4.21b}
		\end{align}
	\end{subequations} \label{11.4.21}
	\begin{tcolorbox}[breakable]
		\begin{remark}
			The two descriptions agree. The conditions (\ref{11.4.21}) split into two cases.
			\begin{itemize}
				\item All $\alpha^{i}\in\mathds{Z}$: $\sum(\alpha^{i})^{2}=2$ forces exactly two entries $\pm1$ and the rest zero. Their sum is $0$ or $\pm2$, always even. These are the $SO(16)$ roots (\ref{11.4.16}), and there are $\binom{8}{2}\cdot2^{2}=112$ of them.
				\item All $\alpha^{i}\in\mathds{Z}+\frac{1}{2}$: each $(\alpha^{i})^{2}\geq\frac{1}{4}$ and $8\cdot\frac{1}{4}=2$, so every entry is $\pm\frac{1}{2}$. With $n_{-}$ minus signs, $\sum\alpha^{i}=4-n_{-}$, which is even iff $n_{-}$ is even. That leaves $2^{7}=128$ vectors, the weights of one chiral $SO(16)$ spinor, which is (\ref{11.4.20}) with an even number of sign flips.
			\end{itemize}
			In total there are $112+128=240$ roots, and with the $8$ Cartan generators $\operatorname{dim}(E_{8})=248$. Every root has $\alpha^{2}=2$, so $E_{8}$ is simply laced.
		\end{remark}
	\end{tcolorbox}
\end{itemize}

For $A_{n}, D_{k}$ and $E_{8}$ (as well as $E_{6}$ and $E_{7}$ we have not yet described), all roots are of equal length. They are called \emph{simply-laced} algebras. For $B_{k}$ and $C_{k}$ (as well as $F_{4}$ and $G_{2}$), there are two different lengths of roots, so they are called long roots and short roots.

Later, a very useful quantity is the \emph{dual Coxeter number} $h(g)$ of Lie algebra $g$, defined as
\begin{equation}
	-\sum_{c,d} f\indices{^a^c_d}f\indices{^b^d_c}=h(g)\psi^{2}d^{ab} \: \label{11.4.22}
\end{equation}
where $\psi$ is any long root. As a reference, the corresponding values for all simple Lie algebras are given in Table \ref{tab:11.3}. Since the inverse of $d^{ab}$ appears in $\psi^{2}=\psi^{i}\psi^{j}d_{ij}$, the definition (\ref{11.4.22}) makes $h(g)$ independent of any normalization of inner product $d^{ab}$.

\begin{table}[h]
	\caption{Dimension and Coxeter number of simple Lie algebras}
	\label{tab:11.3}%
	\centering
	\begin{tabular}[c]{rcccccccc}
		\hline\hline
		\quad\vphantom{\Big(} & $SU(n)$ &  $SO(n),\:n\geq 4$ & $Sp(k)$ & $E_{6}$ & $E_{7}$ & $E_{8}$ & $F_{4}$ & $G_{2}$  \\
		\hline
		$\overset{}{\operatorname{dim}(g)}$ & $n^{2}-1$ & $n(n-1)/2$ & $2k^{2}+k$ & 78 & 133 & 248 & 52 &14 \\
		$h(g)$ & $n$ & $n-2$ & $k+1$ & 12 & 18 & 30 & 9 & 4 \\
		\hline\hline
	\end{tabular}
\end{table}


\begin{tcolorbox}[breakable]
	\begin{remark}
		\textbf{Dimensions.} Traceless Hermitian $n\times n$ matrices have $n^{2}-1$ real parameters. Antisymmetric imaginary ones have $\frac{1}{2}n(n-1)$. $Sp(k)$ has $2k^{2}+k$ (box after (\ref{11.4.10})).

		\textbf{Dual Coxeter number as a sum over roots.} In the adjoint representation $(t^{a})_{cb}=\mi f\indices{^a^b_c}$, since $[T^{a},T^{b}]=\mi f\indices{^a^b_c}T^{c}$. Hence
		\begin{equation*}
			\operatorname{Tr}_{\text{adj}}(t^{a}t^{b})=\sum_{c,d}\mi f\indices{^a^d_c}\;\mi f\indices{^b^c_d}=-\sum_{c,d}f\indices{^a^c_d}f\indices{^b^d_c}=h(g)\psi^{2}d^{ab}\:,
		\end{equation*}
		so (\ref{11.4.22}) says $T_{\text{adj}}=h(g)\psi^{2}$. Now take $a,b$ in the Cartan subalgebra. The weights of the adjoint are the roots, each once, plus $\operatorname{rank}(g)$ zero weights, so
		\begin{equation*}
			h(g)\psi^{2}d^{ij}=\operatorname{Tr}_{\text{adj}}(H^{i}H^{j})=\sum_{\alpha}\alpha^{i}\alpha^{j}
			\quad\xrightarrow{\ \times d_{ij}\ }\quad
			h(g)\,\psi^{2}\operatorname{rank}(g)=\sum_{\alpha}\alpha^{2}\:.
		\end{equation*}
		With $\psi^{2}=2$ this gives the entries of Table \ref{tab:11.3}:
		\begin{align*}
			\text{simply laced }(\alpha^{2}=2):&\quad h=\frac{\operatorname{dim}(g)-\operatorname{rank}(g)}{\operatorname{rank}(g)}\:,\\
			&\quad SU(n): \tfrac{n^{2}-1}{n-1}-1=n\:,\qquad SO(2k): \tfrac{k(2k-1)}{k}-1=2k-2\:,\\
			&\quad E_{6,7,8}: \tfrac{78}{6}-1=12\:,\ \tfrac{133}{7}-1=18\:,\ \tfrac{248}{8}-1=30\:,\\
			SO(2k+1):&\quad \frac{2k(k-1)\cdot2+2k\cdot1}{2k}=2k-1=n-2\:,\\
			Sp(k):&\quad \frac{2k\cdot2+2k(k-1)\cdot1}{2k}=k+1\:.
		\end{align*}
		For $SO(2k+1)$ there are $2k(k-1)$ long roots (\ref{11.4.16}) and $2k$ short roots (\ref{11.4.17}) of length$^{2}$ $1$. For $Sp(k)$, after dividing by $2^{1/2}$, the $2k$ roots (\ref{11.4.18}) are long and the $2k(k-1)$ roots of $SO(2k)$ type are short with length$^{2}$ $1$. The same count gives $F_{4}$ ($24$ long, $24$ short, rank 4): $(48+24)/8=9$, and $G_{2}$ ($6$ long, $6$ short with $\alpha^{2}=\frac{2}{3}$, rank 2): $(12+4)/4=4$. Finally, the trace of (\ref{11.4.7}) over $r$ and the contraction of (\ref{11.4.6}) with $d_{ab}$ both compute $\operatorname{Tr}(t^{a}_{r}t^{b}_{r})d_{ab}$, so $Q_{r}\operatorname{dim}(r)=T_{r}\operatorname{dim}(g)$. For the adjoint this gives
		\begin{equation*}
			Q_{g}=T_{\text{adj}}=h(g)\psi^{2}\:.
		\end{equation*}
	\end{remark}
\end{tcolorbox}

\subsection*{Useful facts for grand unification}

The exceptional group $E_{8}$ is related to the groups appearing in grand unification through several subgroups. This will play a role in the compactification of heterotic strings. 

The first subgroup is
\begin{equation}
	E_{8} \to SU(3) \times E_{6} \:. \label{11.4.23}
\end{equation}
Further decomposition of the $E_{6}$ representation can be seen in Table \ref{tab:11.4}. In the simple compactification of $E_{8}\times E_{8}$ strings, the fermions of the Standard Model can be considered to all come from the 248-dimensional adjoint representation of one of the $E_{8}$. Therefore, tracking the fate of this representation under successive symmetry breaking is beneficial. In $E_{8}\to SU(3)\times E_{6}$,
\begin{equation}
	\mathbf{248} \to (\mathbf{8},\mathbf{1})+ (\mathbf{1},\mathbf{78})+ (\mathbf{3},\mathbf{27})
	+(\bar{\mathbf{3}},\bar{\mathbf{27}}) \:. \label{11.4.24}
\end{equation}
That is, the adjoint representation of $E_{8}$ contains the adjoint representation of the subgroup (the first two representations in (\ref{11.4.24})); among the remaining 162 generators, half transform according to the triplet of $SU(3)$ and the complex 27-dimensional representation of $E_{6}$, and the other half are their conjugates. Further subgroups see Table \ref{tab:11.4}. The first three subgroups represent the breaking of $E_{6}$ via several smaller grand unified subgroups down to the Standard Model group; the fourth is another breaking mode.

\begin{table}[h]
	\caption{Subgroups and representations of grand unified groups}
	\label{tab:11.4}%
	\centering
	\begin{tabular}[c]{rcl}
		\hline\hline
		$E_{6}$ & $\to$ &  $SO(10)\times U(1)$  \\
		$\mathbf{78}$ &  $\to$ & $\mathbf{45}_{0}+\mathbf{16}_{-3}+\overline{\mathbf{16}}_{3}+\mathbf{1}_{0}$ \\
		$\mathbf{27}$ &  $\to$ & $\mathbf{1}_{4}+\mathbf{10}_{-2}+\mathbf{16}_{1}$ \\
		\hline
		$SO(10)$ & $\to$ &  $SU(5)\times U(1)$  \\
		$\mathbf{45}$ & $\to$ & $\mathbf{24}_{0}+\mathbf{10}_{4}+\overline{\mathbf{10}}_{-4}+\mathbf{1}_{0}$ \\
		$\mathbf{16}$ & $\to$ & $\mathbf{10}_{-1} + \bar{\mathbf{5}}_{3}+\mathbf{1}_{-5}$ \\
		$\mathbf{10}$ & $\to$ & $\mathbf{5}_{2}+\bar{\mathbf{5}}_{-2}$ \\
		\hline
		$SU(5)$ & $\to$ &  $SU(3)\times SU(2)\times U(1)$  \\
		$\mathbf{10}$ & $\to$ & $(\mathbf{3},\mathbf{2})_{1}+(\bar{\mathbf{3}},\mathbf{1})_{-4}
		+(\mathbf{1},\mathbf{1})_{6}$ \\
		$\bar{\mathbf{5}}$ & $\to$ & $(\bar{\mathbf{3}},\mathbf{1})_{2} + (\mathbf{1},\mathbf{2})_{-3}$ \\
		\hline
		$E_{6}$ & $\to$ &  $SU(3)\times SU(3)\times SU(3)$  \\
		$\mathbf{78}$ & $\to$ & $(\mathbf{8},\mathbf{1},\mathbf{1})+ (\mathbf{1},\mathbf{8},\mathbf{1}) 
		+(\mathbf{1},\mathbf{1},\mathbf{8}) + (\mathbf{3},\mathbf{3},\mathbf{3})
		+(\bar{\mathbf{3}},\bar{\mathbf{3}},\bar{\mathbf{3}}) $ \\
		$\mathbf{27}$ & $\to$ & $(\mathbf{3},\bar{\mathbf{3}},\mathbf{1}) + (\mathbf{1},\mathbf{3},\bar{\mathbf{3}})
		+(\bar{\mathbf{3}},\mathbf{1},\mathbf{3})$ \\
		\hline\hline
	\end{tabular}
\end{table}

In grand unification, exactly one $SU(3)\times SU(2)\times U(1)$ with quarks and leptons falls into $\mathbf{10}$ plus $\bar{\mathbf{5}}$ of $SU(5)$. Successively tracing back, we see that this generation falls into the representation $\mathbf{16}$ of $SO(10)$, leaving one extra state $\mathbf{1}_{-5}$. This extra state is neutral under $SU(5)$, so it is also neutral under $SU(3)\times SU(2)\times U(1)$, and it can be recognized as a right-handed neutrino. Backtracking to $E_{6}$, $\mathbf{27}$ contains the 15 states of a single generation, besides which there are 12 more states. Compared to $SU(5)$ grand unification, $SO(10)$ and $E_{6}$ are somewhat more unified, meaning that this generation is contained in a single representation, but uneconomically, this representation also contains other unseen states. In fact, the latter might not be that difficult. To see why, examine the decomposition of $\mathbf{27}$ of $E_{6}$ under $SU(3)\times SU(2)\times U(1)\subset SU(5)\subset SO(10)\subset E_{6}$:
\begin{align}
	\mathbf{27} &\to (\mathbf{3},\mathbf{2})_{1} + (\bar{\mathbf{3}},\mathbf{1})_{-4}
	+(\mathbf{1},\mathbf{1})_{6} + (\bar{\mathbf{3}},\mathbf{1})_{2}+(\mathbf{1},\mathbf{2})_{-3} \nonumber \\
	&\quad + [\mathbf{1}_{0}] \nonumber \\
	&\quad +[(\bar{\mathbf{3}},\mathbf{1})_{2}+(\mathbf{3},\mathbf{1})_{-2}]
	+[(\mathbf{1},\mathbf{2})_{-3}+(\mathbf{1},\mathbf{2})_{3}] + [\mathbf{1}_{0}] \:. \label{11.4.25}
\end{align}
The first line lists one generation, the second line is the extra state appearing in $\mathbf{16}$ of $SO(10)$, and the third line is the extra states in $\mathbf{27}$ of $E_{6}$. Each subset in square brackets $[\:]$ is a real representation of $SU(3)\times SU(2)\times U(1)$. The point is that for a real representation $r$, the \textsf{CPT} conjugate is also in representation $r$, so for particles and their antiparticles, this combined gauge plus $SO(2)$ helicity representation is $(r,+\frac{1}{2})+(r,-\frac{1}{2})$. This is the same as the \emph{massive} spin-$\frac{1}{2}$ particle in representation $r$, so gauge and spacetime symmetry are self-consistent with these particles having mass. In the most general invariant action, all particles in $[\:]$ brackets will have very large (grand unification scale) masses. Notably, for $\mathbf{10}+\bar{\mathbf{5}}$ of $SU(5)$, $\mathbf{16}$ of $SO(10)$, or $\mathbf{27}$ of $E_{6}$, the natural $SU(3)\times SU(2)\times U(1)$ spectrum in any of these three is standard quarks and lepton generations.
\begin{tcolorbox}[breakable]
	\begin{remark}
		\textbf{Deriving (\ref{11.4.24}) and the $E_{6}$ lines of Table \ref{tab:11.4}.} Use the chain $E_{8}\supset SO(16)\supset SO(6)\times SO(10)$ with $SO(6)\cong SU(4)\supset SU(3)\times U(1)$. First,
		\begin{align*}
			\mathbf{120}&\eqs{1}(\mathbf{15},\mathbf{1})+(\mathbf{1},\mathbf{45})+(\mathbf{6},\mathbf{10})\:,\qquad
			\mathbf{128}\eqs{2}(\mathbf{4},\mathbf{16})+(\bar{\mathbf{4}},\overline{\mathbf{16}})\:.
		\end{align*}
		At (1) the antisymmetric tensor of $\mathbf{16}=(\mathbf{6},\mathbf{1})+(\mathbf{1},\mathbf{10})$ was split into its three blocks, $15+45+60=120$. At (2) the $SO(16)$ spinor weights $(\pm\frac{1}{2})^{8}$ with an even number of minus signs split into (even, even) and (odd, odd) numbers of minus signs among the first 3 and last 5 entries. That gives spinor $\otimes$ spinor, with the $SO(6)$ spinors $\mathbf{4},\bar{\mathbf{4}}$ being the $SU(4)$ fundamentals. Next, with $\mathbf{4}\to\mathbf{3}_{1}+\mathbf{1}_{-3}$ under $SU(3)\times U(1)$ (the charge is traceless: $3\cdot1-3=0$),
		\begin{align*}
			\mathbf{6}=[\mathbf{4}\times\mathbf{4}]_{\text{A}}&\eqs{3}\bar{\mathbf{3}}_{2}+\mathbf{3}_{-2}\:,\qquad
			\mathbf{15}=\mathbf{4}\times\bar{\mathbf{4}}-\mathbf{1}\eqs{4}\mathbf{8}_{0}+\mathbf{1}_{0}+\mathbf{3}_{4}+\bar{\mathbf{3}}_{-4}\:.
		\end{align*}
		At (3) $[\mathbf{3}_{1}\times\mathbf{3}_{1}]_{\text{A}}=\bar{\mathbf{3}}_{2}$ and $\mathbf{3}_{1}\times\mathbf{1}_{-3}=\mathbf{3}_{-2}$. At (4) $(\mathbf{3}_{1}+\mathbf{1}_{-3})(\bar{\mathbf{3}}_{-1}+\mathbf{1}_{3})=\mathbf{8}_{0}+\mathbf{1}_{0}+\mathbf{3}_{4}+\bar{\mathbf{3}}_{-4}+\mathbf{1}_{0}$, and one singlet was removed. Collecting the $248=120+128$ states by their $SU(3)$ representation:
		\begin{align*}
			\mathbf{8}:&\quad \mathbf{8}_{0}\ (\text{from }\mathbf{15})\:,\\
			\mathbf{3}:&\quad \mathbf{3}_{4}\otimes\mathbf{1}\ (\mathbf{15})+\mathbf{3}_{-2}\otimes\mathbf{10}\ (\mathbf{6},\mathbf{10})+\mathbf{3}_{1}\otimes\mathbf{16}\ (\mathbf{4},\mathbf{16})
			=\mathbf{3}\otimes\bigl[\mathbf{1}_{4}+\mathbf{10}_{-2}+\mathbf{16}_{1}\bigr]\:,\\
			\bar{\mathbf{3}}:&\quad \text{the conjugates}\:,\\
			\mathbf{1}:&\quad \mathbf{1}_{0}\ (\mathbf{15})+\mathbf{45}_{0}+\mathbf{16}_{-3}\ (\mathbf{1}_{-3}\subset\mathbf{4})+\overline{\mathbf{16}}_{3}\:.
		\end{align*}
		The $SU(3)$ singlets form the adjoint of the commutant, $\mathbf{78}=\mathbf{45}_{0}+\mathbf{16}_{-3}+\overline{\mathbf{16}}_{3}+\mathbf{1}_{0}$ with $45+32+1=78$. The triplet partners form $\mathbf{27}=\mathbf{1}_{4}+\mathbf{10}_{-2}+\mathbf{16}_{1}$ with $1+10+16=27$. This is (\ref{11.4.24}), $248=8+78+3\cdot27+3\cdot27$, together with the first two rows of Table \ref{tab:11.4}.

		\textbf{$SO(10)\to SU(5)\times U(1)$.} The $SO(10)$ spinor weights are $(\pm\frac{1}{2})^{5}$ with an even number $n_{-}$ of minus signs, and the $U(1)\subset U(5)$ is $Q=-2\sum_{i}s_{i}$:
		\begin{equation*}
			n_{-}=0:\ Q=-5\ (1\text{ state})\:,\quad n_{-}=2:\ Q=-1\ \bigl(\tbinom{5}{2}=10\bigr)\:,\quad n_{-}=4:\ Q=+3\ \bigl(\tbinom{5}{4}=5\bigr)\:,
		\end{equation*}
		i.e.\ $\mathbf{16}\to\mathbf{1}_{-5}+\mathbf{10}_{-1}+\bar{\mathbf{5}}_{3}$. The vector weights $\pm e_{i}$ give $Q=\mp2$, so $\mathbf{10}\to\bar{\mathbf{5}}_{-2}+\mathbf{5}_{2}$. The adjoint weights $\pm e_{i}\pm e_{j}$ and $0^{5}$ give $\mathbf{45}\to\mathbf{10}_{4}+\overline{\mathbf{10}}_{-4}+(\mathbf{24}+\mathbf{1})_{0}$.

		\textbf{$SU(5)\to SU(3)\times SU(2)\times U(1)$.} With hypercharge $Y=\operatorname{diag}(-2,-2,-2,3,3)$ on the $\mathbf{5}$, $\mathbf{5}\to(\mathbf{3},\mathbf{1})_{-2}+(\mathbf{1},\mathbf{2})_{3}$ and so $\bar{\mathbf{5}}\to(\bar{\mathbf{3}},\mathbf{1})_{2}+(\mathbf{1},\mathbf{2})_{-3}$. Then
		\begin{equation*}
			\mathbf{10}=[\mathbf{5}\times\mathbf{5}]_{\text{A}}\eqs{5}(\bar{\mathbf{3}},\mathbf{1})_{-2-2}+(\mathbf{3},\mathbf{2})_{-2+3}+(\mathbf{1},\mathbf{1})_{3+3}
			=(\bar{\mathbf{3}},\mathbf{1})_{-4}+(\mathbf{3},\mathbf{2})_{1}+(\mathbf{1},\mathbf{1})_{6}\:.
		\end{equation*}
		At (5) $[\mathbf{3}\times\mathbf{3}]_{\text{A}}=\bar{\mathbf{3}}$ and $[\mathbf{2}\times\mathbf{2}]_{\text{A}}=\mathbf{1}$, and the charges add.
		In every line of Table \ref{tab:11.4} the $U(1)$ charges sum to zero over each representation, e.g.\ $\mathbf{27}$: $4-20+16=0$; $\mathbf{16}$: $-10+15-5=0$. This is required because a $U(1)$ generator inside a simple group is traceless.

		\textbf{(\ref{11.4.25}).} Composing the chain, $\mathbf{27}\to\mathbf{16}_{1}+\mathbf{10}_{-2}+\mathbf{1}_{4}$. The $\mathbf{16}$ gives $\mathbf{10}+\bar{\mathbf{5}}+\mathbf{1}$ of $SU(5)$, i.e.\ the first line of (\ref{11.4.25}) plus the $[\mathbf{1}_{0}]$ of the second line. The $\mathbf{10}$ of $SO(10)$ gives $\mathbf{5}+\bar{\mathbf{5}}=[(\mathbf{3},\mathbf{1})_{-2}+(\bar{\mathbf{3}},\mathbf{1})_{2}]+[(\mathbf{1},\mathbf{2})_{3}+(\mathbf{1},\mathbf{2})_{-3}]$. The $\mathbf{1}_{4}$ gives the last $[\mathbf{1}_{0}]$. Only the $SU(3)\times SU(2)\times U(1)$ hypercharge is displayed there, so the extra $U(1)$ labels are dropped.
	\end{remark}
\end{tcolorbox}

\section{Current Algebra} \label{sec:11.5}

In heterotic strings, the form of gauge boson vertex operators is $j(z)\tilde{\psi}^{\mu}(\bar{z})\me^{\mi k\cdot X}$, where $j(z)$ is the fermion bilinear form $\lambda^{A}\lambda^{B}$ or the spin field (\ref{11.2.26}). Similarly, for the toroidal compactification of bosonic strings, the form of gauge boson vertex operators is $j(z)\bar{\partial}X^{\mu}(\bar{z})\me^{\mi k\cdot X}$, where $j(z)$ is $\partial X^{m}$ for Kaluza-Klein gauge bosons and the exponential $\exp[\pm 2\mi\alpha^{\prime}X^{m}_{L,R}]$ for expanded gauge symmetries. All these currents are holomorphic $(1,0)$ operators.

We examine the set of $(1,0)$ currents $j^{a}(z)$ in a general CFT. Conformal invariance requires their OPEs to take the following form
\begin{equation}
	j^{a}(z)j^{b}(0) \sim \frac{k^{ab}}{z^{2}} + \frac{\mi c\indices{^a^b_c}}{z}\,j^{c}(0) \label{11.5.1}
\end{equation}
where $k^{ab}$ and $c\indices{^a^b_c}$ are constants. From a dimensional point of view, the coefficient of the $z^{-2}$ term must be a constant, and the coefficient of the $z^{-1}$ term must be proportional to a current. The Laurent coefficients
\begin{equation}
	j^{a}(z) = \sum_{m=-\infty}^{\infty} \frac{j_{m}^{a}}{z^{m+1}} \label{11.5.2}
\end{equation}
thus satisfy a closed algebra
\begin{equation}
	[j_{m}^{a},j_{n}^{b}]=mk^{ab}\delta_{m,-n}+\mi c\indices{^a^b_c}j^{c}_{m+n} \:. \label{11.5.3}
\end{equation}
\begin{tcolorbox}[breakable]
	Now prove (\ref{11.5.3}), using
	\[
	j_{m}^{a} = \oint_{0}\dif z\: j^{a}(z)z^{m} \:,
	\]
	The left side of equation (\ref{11.5.3}) can be written as
	\begin{align*}
		[j_{m}^{a},j_{n}^{b}] &= \frac{1}{(2\pi\mi)^{2}}\oint_{0} \dif w\: w^{n} \oint_{w}\dif z\: z^{m}
		\biggl\{\frac{k^{ab}}{(z-w)^{2}}+\frac{\mi c\indices{^a^b_c}}{z-w}j^{c}(w) \biggr\} \\
		&=\frac{1}{2\pi\mi} \oint_{0} \dif w\: w^{n} \biggl\{mw^{m-1}k^{ab}+ \mi c\indices{^a^b_c}w^{m}j^{c}(w) \biggr\}\\
		&=mk^{ab}\delta_{m,-n}+\mi c\indices{^a^b_c}j^{c}_{m+n} \:.
	\end{align*}
\end{tcolorbox}
\noindent In particular,
\begin{equation}
	[j^{a}_{0},j^{b}_{0}] = \mi c\indices{^a^b_c}j_{0}^{c} \:. \label{11.5.4}
\end{equation}
That is, the $m=0$ modes form the Lie algebra $g$, and
\begin{equation}
	c\indices{^a^b_c}=f\indices{^a^b_c} \:. \label{11.5.5}
\end{equation}
We first restrict to the case where $g$ is a simple algebra. The Jacobi identity for $j_{1}^{a}j_{0}^{b}j_{-1}^{c}$ requires
\begin{equation}
	f\indices{^b^c_d}k^{ad}+f\indices{^b^a_d}k^{dc} = 0 \:. \label{11.5.6}
\end{equation}
\begin{tcolorbox}[breakable]
	\begin{align*}
		0 &= [j_{1}^{a},[j_{0}^{b},j_{-1}^{c}]]+[j_{0}^{b},[j_{-1}^{c},j_{1}^{a}]]
		+[j_{-1}^{c},[j_{1}^{a},j_{0}^{b}]] \\
		&= [j_{1}^{a},\mi c\indices{^b^c_d}j^{d}_{-1}]+ [j_{0}^{b},-k^{ca}+\mi c\indices{^c^a_d}j_{0}^{d}]
		+[j_{-1}^{c},\mi c\indices{^a^b_d}j_{1}^{d}]\\
		&=\mi c\indices{^b^c_d}(k^{ad}+\mi c\indices{^a^d_e}j^{e}_{0})+\mi c\indices{^c^a_d}\mi c\indices{^b^d_e}j^{e}_{0}
		+\mi c\indices{^a^b_d}(-k^{cd}+\mi c\indices{^c^d_e}j^{e}_{0})
	\end{align*}
	By definition, $k^{ab}$ is symmetric with respect to $a,b$, and $c\indices{^a^b_c}$ is antisymmetric with respect to $a,b$; using the normal Jacobi identity on the $c\,c$ terms proves (\ref{11.5.6}).
\end{tcolorbox}
\noindent This relationship is the same as equation (\ref{11.4.5}) defining the inner product $d^{ab}$ of a Lie algebra, and since we assumed $g$ is simple, it must be
\begin{equation}
	k^{ab}=\hat{k}d^{ab} \:, \label{11.5.7}
\end{equation}
where $\hat{k}$ is some constant. Algebra (\ref{11.5.3}) has various names: \emph{current algebra}, \emph{affine Lie algebra}, (\emph{affine}) \emph{Kac-Moody algebra}. Currents are $(1,0)$ tensors, so
\begin{equation}
	[L_{m},j_{n}^{a}] = -nj_{m+n}^{a} \:. \label{11.5.8}
\end{equation}
\begin{tcolorbox}[breakable]
	\begin{remark}
		With $L_{m}=\oint\frac{\dif z}{2\pi\mi}z^{m+1}T_{B}(z)$, $j^{a}_{n}=\oint\frac{\dif w}{2\pi\mi}w^{n}j^{a}(w)$ and the weight-one OPE $T_{B}(z)j^{a}(w)\sim\frac{j^{a}(w)}{(z-w)^{2}}+\frac{\partial j^{a}(w)}{z-w}$,
		\begin{align*}
			[L_{m},j^{a}_{n}]&\eqs{1}\oint_{0}\frac{\dif w}{2\pi\mi}\,w^{n}\oint_{w}\frac{\dif z}{2\pi\mi}\,z^{m+1}\Bigl[\frac{j^{a}(w)}{(z-w)^{2}}+\frac{\partial j^{a}(w)}{z-w}\Bigr]\\
			&=\oint_{0}\frac{\dif w}{2\pi\mi}\Bigl[(m+1)w^{m+n}j^{a}(w)+w^{m+n+1}\partial j^{a}(w)\Bigr]\\
			&\eqs{2}\bigl[(m+1)-(m+n+1)\bigr]\oint_{0}\frac{\dif w}{2\pi\mi}\,w^{m+n}j^{a}(w)=-n\,j^{a}_{m+n}\:.
		\end{align*}
		At (1) the commutator of two contour integrals was turned into a $z$ contour around $w$. At (2) the second term was integrated by parts in $w$.
	\end{remark}
\end{tcolorbox}

Physically speaking, $j_{n}^{a}$ generates position-dependent $g$-transformations. Due to the existence of local currents, this is possible in quantum field theory. In a unitary theory, the \emph{central extension} or \emph{Schwinger term} $\hat{k}$ must always be positive. To prove this, note that
\begin{equation}
	\hat{k}d^{aa}= \langle 1\rvert\,[j_{1}^{a},j_{-1}^{a}]\,\lvert 1 \rangle = \lVert\,j_{-1}^{a}\lvert1\rangle\,\rVert^{2}  \label{11.5.9}
\end{equation}
(without summation over $a$). For a compact Lie algebra, $d^{aa}$ is positive, so $\hat{k}$ must be non-negative. It is zero if and only if $j_{-1}^{a}\lvert1\rangle=0$, but the vertex operator of $j_{-1}^{a}\lvert1\rangle$ is current $j^{a}$: any matrix element of $j^{a}$ can be obtained by pasting $j_{-1}^{a}$ into the world-sheet. Therefore, $\hat{k}=0$ if and only if the current is identically zero.

The coefficient $\hat{k}$ is quantized. To prove this, examine any root $\alpha$. Define
\begin{equation}
	J^{3} =\frac{\alpha\cdot H}{\alpha^{2}}\:, \qquad J^{\pm}=E^{\pm\alpha} \:, \label{11.5.10}
\end{equation}
using the general form (\ref{11.4.12}), it can be verified that they satisfy the $SU(2)$ algebra
\begin{equation}
	[J^{3},J^{\pm}]=J^{\pm}\:, \qquad [J^{+},J^{-}]=2J^{3} \:. \label{11.5.11}
\end{equation}
It can be verified that
\begin{subequations}
	\begin{align}
		&\frac{\alpha\cdot H_{0}}{\alpha^{2}}\:, \qquad E_{0}^{\alpha}\:,\qquad E_{0}^{-\alpha} \:, \label{11.5.12a}\\
		&\frac{\alpha\cdot H_{0}+\hat{k}}{\alpha^{2}}\:, \qquad E_{1}^{\alpha}\:, \qquad E_{-1}^{-\alpha} \:\label{11.5.12b}
	\end{align} \label{15.5.12}
\end{subequations}
\begin{tcolorbox}[breakable]
	\begin{remark}
		\textbf{(\ref{11.5.11}).} From (\ref{11.4.11}) and (\ref{11.4.12}),
		\begin{equation*}
			[J^{3},J^{\pm}]=\frac{\alpha_{i}[H^{i},E^{\pm\alpha}]}{\alpha^{2}}=\pm\frac{\alpha\cdot\alpha}{\alpha^{2}}E^{\pm\alpha}=\pm J^{\pm}\:,\qquad
			[J^{+},J^{-}]=[E^{\alpha},E^{-\alpha}]=\frac{2\alpha\cdot H}{\alpha^{2}}=2J^{3}\:.
		\end{equation*}
		\textbf{(\ref{11.5.12a}).} By (\ref{11.5.4}) the zero modes obey exactly the Lie algebra, so the same computation applies.

		\textbf{(\ref{11.5.12b}).} Now the central term of (\ref{11.5.3}) enters. With $k^{ab}=\hat{k}d^{ab}$, the central term of $[E^{\alpha}_{m},E^{-\alpha}_{n}]$ is $m\hat{k}(E^{\alpha},E^{-\alpha})\delta_{m,-n}=m\hat{k}\frac{2}{\alpha^{2}}\delta_{m,-n}$, using the normalization below (\ref{11.4.12}). Then
		\begin{align*}
			[J^{+},J^{-}]&=[E^{\alpha}_{1},E^{-\alpha}_{-1}]\eqs{1}\frac{2\alpha\cdot H_{0}}{\alpha^{2}}+1\cdot\hat{k}\,\frac{2}{\alpha^{2}}=2\,\frac{\alpha\cdot H_{0}+\hat{k}}{\alpha^{2}}=2J^{3}\:,\\
			[J^{3},J^{+}]&=\frac{1}{\alpha^{2}}[\alpha\cdot H_{0},E^{\alpha}_{1}]\eqs{2}\frac{\alpha\cdot\alpha}{\alpha^{2}}E^{\alpha}_{1}=J^{+}\:,\qquad
			[J^{3},J^{-}]=\frac{1}{\alpha^{2}}[\alpha\cdot H_{0},E^{-\alpha}_{-1}]=-J^{-}\:.
		\end{align*}
		At (1) the non-central part is the Lie bracket with mode number $1+(-1)=0$. At (2) the central term of $[H^{i}_{0},E^{\alpha}_{1}]$ vanishes because $0\cdot\delta_{0,-1}=0$, and $\hat{k}$ is a c-number. So the pseudospin is an $SU(2)$ algebra whose $J^{3}$ is shifted by $\hat{k}/\alpha^{2}$. On any state the eigenvalues of $2J^{3}$ for both algebras are integers. Taking the difference, $2\hat{k}/\alpha^{2}\in\mathds{Z}$.
	\end{remark}
\end{tcolorbox}
also satisfy the $SU(2)$ algebra. The first is the usual center-of-mass Lie algebra, and the second is called \emph{pseudospin}. Since the eigenvalues of $2J_{3}$ in $SU(2)$ must be integers, $2\hat{k}/\alpha^{2}$ must also be an integer. If $\alpha$ is taken to be the long root (denoted as $\psi$) in the algebra, this condition is the strongest. Thus the \emph{level}
\begin{equation}
	k=\frac{2\hat{k}}{\psi^{2}} \label{11.5.13}
\end{equation}
is a non-negative integer, and for non-trivial currents, it is positive.

It is customary to normalize the inner product of the Lie algebra so that the length of the long root is $\sqrt{2}$, which makes $\hat{k}=k$ the coefficient of the leading term in the OPE. Incidentally, from this, it follows that under this normalization, the generators (\ref{11.4.14}) are normalized, so the $SO(n)$ inner product is half the trace of the vector representation. Similarly, the inner product of $SU(n)$ is equal to the trace in the fundamental representation. From then on, we take a set of bases for the generators such that $d^{ab}=\delta^{ab}$.
\begin{tcolorbox}[breakable]
	\begin{remark}
		The roots in (\ref{11.4.16}) were computed from $[H^{i},E^{\alpha}]=\alpha^{i}E^{\alpha}$ with the $H^{i}$ of (\ref{11.4.14}). Their length is $\alpha^{2}=\alpha^{i}\alpha^{j}d_{ij}$, and $(\pm1,\pm1,0^{k-2})$ has length $\sqrt{2}$ exactly when $d_{ij}=\delta_{ij}$, i.e.\ $(H^{i},H^{j})=\delta^{ij}$. In the vector representation
		\begin{equation*}
			\operatorname{Tr}_{v}(H^{i}H^{j})=\delta^{ij}\operatorname{Tr}\begin{bmatrix}0&\mi\\-\mi&0\end{bmatrix}^{2}=\delta^{ij}\operatorname{Tr}I_{2}=2\delta^{ij}
			\quad\Longrightarrow\quad (T,T')=\tfrac{1}{2}\operatorname{Tr}_{v}(TT')\:,
		\end{equation*}
		and an invariant inner product on a simple algebra is fixed once it is fixed on the Cartan generators. For $SU(n)$, $H^{i}=E_{ii}$ gives the roots $e_{i}-e_{j}$ of (\ref{11.4.19}), with length $\sqrt{2}$ in the metric $\operatorname{Tr}_{\text{f}}(H^{i}H^{j})=\delta^{ij}$. So $(T,T')=\operatorname{Tr}_{\text{f}}(TT')$, restricted to the traceless hyperplane. In the language of (\ref{11.4.6}) this means $T_{v}=2=\psi^{2}$ for $SO(n)$ and $T_{\text{f}}=1=\frac{1}{2}\psi^{2}$ for $SU(n)$.
	\end{remark}
\end{tcolorbox}


The level represents the relative size of the $z^{-2}$ and $z^{-1}$ terms in the OPE. For $U(1)$, the structure constant is zero, and only the $z^{-2}$ term appears. Therefore, there is no such thing as a level for $U(1)$. It will be more convenient to normalize all $U(1)$ currents to
\begin{equation}
	j^{a}(z)j^{b}(0) \sim \frac{\delta^{ab}}{z^{2}} \label{11.5.14}
\end{equation}
From this OPE and holomorphicity, it can be concluded that each $U(1)$ current algebra is isomorphic to a free bosonic CFT,
\begin{equation}
	j^{a}=\mi\partial H^{a} \:. \label{11.5.15}
\end{equation}

In heterotic strings, current algebra is composed of $n$ real fermions $\lambda^{A}(z)$. Currents
\begin{equation}
	\mi\lambda^{A}\lambda^{B} \:. \label{11.5.16}
\end{equation}
form the $SO(n)$ algebra. The maximal commuting set of currents is $\mi\lambda^{2K-1}\lambda^{2K}$, where $K=1,\cdots,[n/2]$. These correspond to generators (\ref{11.4.14}), and their normalization makes the length of root (\ref{11.4.16}) $\sqrt{2}$. The level is the coefficient of the leading term in the OPE; the leading term is $1/z^{2}$, so the level is $k=1$. The case $n=3$ is an exception: there is no long root, only short roots $\pm1$, so we must readjust the diagonal currents to $2^{1/2}\mi\lambda^{1}\lambda^{2}$, and the level is $k=2$.
\begin{tcolorbox}[breakable]
	\begin{remark}
		Let $J^{AB}=\mi\lambda^{A}\lambda^{B}$ and use $\lambda^{A}(z)\lambda^{B}(0)\sim\delta^{AB}/z$. Wick's theorem for $J^{AB}(z)J^{CD}(0)=-\lambda^{A}_{1}\lambda^{B}_{1}\lambda^{C}_{2}\lambda^{D}_{2}$ gives
		\begin{align*}
			J^{AB}(z)J^{CD}(0)&\eqs{1}-\frac{\delta^{AD}\delta^{BC}-\delta^{AC}\delta^{BD}}{z^{2}}\\
			&\qquad-\frac{1}{z}\Bigl[\delta^{BC}\lambda^{A}\lambda^{D}-\delta^{AC}\lambda^{B}\lambda^{D}-\delta^{BD}\lambda^{A}\lambda^{C}+\delta^{AD}\lambda^{B}\lambda^{C}\Bigr](0)\\
			&\eqs{2}\frac{\delta^{AC}\delta^{BD}-\delta^{AD}\delta^{BC}}{z^{2}}\\
			&\qquad+\frac{\mi}{z}\Bigl[\delta^{BC}J^{AD}-\delta^{AC}J^{BD}-\delta^{BD}J^{AC}+\delta^{AD}J^{BC}\Bigr](0)\:.
		\end{align*}
		At (1) the double contraction $(A,D)(B,C)$ is nested ($+$) and $(A,C)(B,D)$ is crossed ($-$). In each single contraction the sign counts the fermions that had to be moved to bring the contracted pair together. At (2) $\lambda^{A}\lambda^{D}=-\mi J^{AD}$. The $1/z$ term is $\mi$ times the $SO(n)$ commutator of the generators $(t^{AB})_{CD}\propto\delta^{AC}\delta^{BD}-\delta^{AD}\delta^{BC}$, so the currents form the $SO(n)$ current algebra. For the Cartan currents,
		\begin{equation*}
			J^{2K-1,2K}(z)\,J^{2L-1,2L}(0)\sim\frac{\delta^{KL}}{z^{2}}\:,
		\end{equation*}
		so $k^{ab}=\delta^{ab}$ in the normalization where these currents are orthonormal, which by the previous box is the one with $\psi^{2}=2$. Hence $\hat{k}=1$ and $k=2\hat{k}/\psi^{2}=1$. For $n=3$ the only roots are $\pm1$ (short), so the Cartan current must be rescaled to $2^{1/2}J^{12}$ to make the root length $\sqrt{2}$. Its OPE is $2/z^{2}$, so $k=2$.
	\end{remark}
\end{tcolorbox}

For any real representation $r$ of any Lie algebra, we can build currents with $\operatorname{dim}(r)$ real fermions
\begin{equation}
	\frac{1}{2}\lambda^{A}\lambda^{B}\,t_{r,AB}^{a} \:. \label{11.5.17}
\end{equation}
They satisfy a current algebra of level $k=T_{r}/\psi^{2}$, where $T_{r}$ is defined by equation (\ref{11.4.6}). The situation in the previous paragraph is the $n$-dimensional vector representation of $SO(n)$, for which $T_{R}=\psi^{2}$. Another example is $nk$ fermions transforming according to the $k$-fold vector representation, which gives level $k$.
\begin{tcolorbox}[breakable]
	\begin{remark}
		For a real representation the Hermitian matrices $t^{a}_{r}$ can be chosen imaginary and antisymmetric, $t^{a}_{r,BA}=-t^{a}_{r,AB}$. Using the double contractions from the previous box,
		\begin{align*}
			j^{a}(z)j^{b}(0)\Big|_{z^{-2}}&=\tfrac{1}{4}t^{a}_{r,AB}t^{b}_{r,CD}\,\frac{\delta^{AD}\delta^{BC}-\delta^{AC}\delta^{BD}}{z^{2}}
			=\frac{1}{4z^{2}}\Bigl[t^{a}_{r,AB}t^{b}_{r,BA}-t^{a}_{r,AB}t^{b}_{r,AB}\Bigr]\\
			&\eqs{1}\frac{1}{4z^{2}}\Bigl[\operatorname{Tr}(t^{a}_{r}t^{b}_{r})+\operatorname{Tr}(t^{a}_{r}t^{b}_{r})\Bigr]\eqs{2}\frac{T_{r}\,d^{ab}}{2z^{2}}\:.
		\end{align*}
		At (1) antisymmetry turned $-t^{b}_{AB}$ into $t^{b}_{BA}$. At (2) (\ref{11.4.6}) was used. So $\hat{k}=\frac{1}{2}T_{r}$ and $k=2\hat{k}/\psi^{2}=T_{r}/\psi^{2}$. The four single contractions give $\frac{1}{4z}\lambda\bigl(3t^{a}_{r}t^{b}_{r}-t^{b}_{r}t^{a}_{r}\bigr)\lambda$. Only the antisymmetric part of the matrix survives between $\lambda^{A}\lambda^{B}$, and $(t^{a}t^{b})^{\mathrm{T}}=t^{b}t^{a}$, so this is $\frac{1}{2z}\lambda[t^{a}_{r},t^{b}_{r}]\lambda=\frac{\mi f\indices{^a^b_c}}{z}j^{c}$, reproducing the Lie algebra. For the vector of $SO(n)$, $T_{v}=\psi^{2}$ (previous boxes), so $k=1$. For $k$ copies of the vector $T_{r}=kT_{v}$, so the level is $k$.
	\end{remark}
\end{tcolorbox}

As a final example, we examine the $SU(2)$ symmetry at the self-dual point of toroidal compactification. Its current is $\exp[2^{1/2}\mi H(z)]$. Then current $\mi\partial H$ is normalized to weight $2^{1/2}$ and length $\sqrt{2}$. The OPE of $\mi\partial H$ with itself starts from $1/z^{2}$, so the level is still $k=1$.

In some cases, there may be sectors where the current is not periodic, $j^{a}(w+2\pi)=R^{ab}\,j^{b}(w)$, where $R^{ab}$ is any automorphism of the algebra. In these cases, the modes of the current are fractional and satisfy \emph{twisted} affine Lie algebras.

\subsection*{The Sugawara construction}

In current algebras with conformal symmetry, there is a remarkable relationship between the energy-momentum tensor and the currents, which gives many interesting structures. Define the operator
\begin{equation}
	:jj(z_{1}):=\lim_{z_{2}\to z_{1}}\Biggl(j^{a}(z_{1})j^{a}(z_{2})-\frac{\hat{k}\operatorname{dim}(g)}{z_{12}^{2}}\Biggr) \:, \label{11.5.18}
\end{equation}
where summation over $a$ is implied. We first wish to prove that, up to a normalization, the OPE of $:jj: $ with $j^{a}$ is the same as the OPE of $T_{B}$ with $j^{a}$. 

The OPE of the product $:jj: $ is not the product of OPEs, because compared with the distance to the third current, the two currents in $:jj: $ are closer to each other; we must use holomorphicity to make a less direct discussion. Examine the following product:
\begin{align}
	j^{a}(z_{1})j^{a}(z_{2})j^{c}(z_{3})&=\frac{\hat{k}}{z_{31}^{2}}j^{c}(z_{2})+\frac{\mi f^{cad}}{z_{31}}j^{d}(z_{1})j^{a}(z_{2})+\frac{\hat{k}}{z_{32}^{2}}j^{c}(z_{1}) \nonumber \\
	&\quad +\frac{\mi f^{cad}}{z_{32}}j^{a}(z_{1})j^{d}(z_{2})+ \text{terms holomorphic in } z_{3} \:. \label{11.5.19}
\end{align}
We have used the current-current OPE to determine the singularity as $z_{3}$ approaches $z_{1}$ or $z_{2}$. Take $z_{2}\to z_{1}$ in this relationship and do a Laurent expansion in $z_{21}$, keeping the $z_{21}^{0}$ order terms, and obtain
\begin{align}
	:jj(z_{1}):j^{c}(z_{3}) &\sim \frac{2\hat{k}}{z_{13}^{2}}j^{c}(z_{1})+\frac{f^{cad}f^{ead}}{z_{13}^{2}}\,j^{e}(z_{1}) \nonumber\\
	&=\frac{2\hat{k}+h(g)\psi^{2}}{z_{13}^{2}}j^{c}(z_{1}) \nonumber \\
	&=(k+h(g))\psi^{2}\biggl[\frac{1}{z_{13}^{2}}j^{c}(z_{3})+\frac{1}{z_{13}}\partial j^{c}(z_{3})\biggr] \:. \label{11.5.20}
\end{align}
where $h(g)$ is the dual Coxeter number.
\begin{tcolorbox}[breakable]
	\begin{remark}
		Set $\epsilon=z_{12}$, so $z_{2}=z_{1}-\epsilon$ and $\frac{1}{z_{32}}=\frac{1}{z_{31}+\epsilon}=\frac{1}{z_{31}}-\frac{\epsilon}{z_{31}^{2}}+O(\epsilon^{2})$. By definition (\ref{11.5.18}), $:\!jj(z_{1})\!:j^{c}(z_{3})$ is the $\epsilon^{0}$ term of (\ref{11.5.19}); the subtracted c-number does not depend on $z_{3}$. The four terms of (\ref{11.5.19}) contribute as follows.

		\textbf{Double-pole terms.} $\frac{\hat{k}}{z_{31}^{2}}j^{c}(z_{2})+\frac{\hat{k}}{z_{32}^{2}}j^{c}(z_{1})\to\frac{2\hat{k}}{z_{31}^{2}}j^{c}(z_{1})$ at order $\epsilon^{0}$.

		\textbf{Simple-pole terms.} Expand the remaining current pairs with the OPE. The $\hat{k}\delta^{ad}$ pieces drop out against $f^{cad}$, which leaves
		\begin{align*}
			\frac{\mi f^{cad}}{z_{31}}j^{d}(z_{1})j^{a}(z_{2})&=\frac{\mi f^{cad}}{z_{31}}\Bigl[\frac{\mi f^{dae}}{\epsilon}j^{e}(z_{2})+(j^{d}j^{a})(z_{1})+O(\epsilon)\Bigr]\:,\\
			\frac{\mi f^{cad}}{z_{32}}j^{a}(z_{1})j^{d}(z_{2})&=\mi f^{cad}\Bigl(\frac{1}{z_{31}}-\frac{\epsilon}{z_{31}^{2}}\Bigr)\Bigl[\frac{\mi f^{ade}}{\epsilon}j^{e}(z_{2})+(j^{a}j^{d})(z_{1})+O(\epsilon)\Bigr]\:,
		\end{align*}
		where $(AB)$ denotes the non-singular part of the product. Now use $j^{e}(z_{2})=j^{e}(z_{1})-\epsilon\,\partial j^{e}(z_{1})+\cdots$ and $f^{ade}=-f^{dae}$. Order by order:
		\begin{align*}
			\epsilon^{-1}:&\quad -\frac{f^{cad}f^{dae}}{z_{31}}j^{e}-\frac{f^{cad}f^{ade}}{z_{31}}j^{e}=0\:,\\
			\epsilon^{0},\ \partial j:&\quad +\frac{f^{cad}f^{dae}}{z_{31}}\partial j^{e}+\frac{f^{cad}f^{ade}}{z_{31}}\partial j^{e}=0\:,\\
			\epsilon^{0},\ j:&\quad +\frac{f^{cad}f^{ade}}{z_{31}^{2}}j^{e}\eqs{1}\frac{f^{cad}f^{ead}}{z_{31}^{2}}j^{e}\eqs{2}\frac{h(g)\psi^{2}}{z_{31}^{2}}j^{c}\:,\\
			\epsilon^{0},\ (jj):&\quad \frac{\mi f^{cad}}{z_{31}}\Bigl[(j^{d}j^{a})+(j^{a}j^{d})\Bigr](z_{1})\eqs{3}0\:.
		\end{align*}
		At (1) $f^{ade}=f^{ead}$ by cyclicity. At (2) $f^{cad}f^{ead}=-\sum f^{cad}f^{eda}=h(g)\psi^{2}\delta^{ce}$ by (\ref{11.4.22}). At (3) the bracket goes into itself under $a\leftrightarrow d$, while $f^{cad}$ is antisymmetric, so the contraction vanishes. This is why only the singular parts of the $jj$ OPE contribute. Adding everything and using $z_{31}^{2}=z_{13}^{2}$,
		\begin{equation*}
			:\!jj(z_{1})\!:j^{c}(z_{3})\sim\frac{2\hat{k}+h(g)\psi^{2}}{z_{13}^{2}}j^{c}(z_{1})
			\eqs{4}(k+h(g))\psi^{2}\Bigl[\frac{j^{c}(z_{3})}{z_{13}^{2}}+\frac{\partial j^{c}(z_{3})}{z_{13}}\Bigr]\:.
		\end{equation*}
		At (4) $2\hat{k}=k\psi^{2}$ by (\ref{11.5.13}), and $j^{c}(z_{1})=j^{c}(z_{3})+z_{13}\partial j^{c}(z_{3})+\cdots$. This is (\ref{11.5.20}).
	\end{remark}
\end{tcolorbox}
\noindent Define
\begin{equation}
	T_{B}^{\mathrm{s}}(z) =\frac{1}{(k+h(g))\psi^{2}}\,:jj(z):\:. \label{11.5.21}
\end{equation}
The OPE of $T_{B}^{\mathrm{s}}$ with the currents is the same as the OPE of the energy-momentum tensor $T_{B}(z)$ with the currents,
\begin{equation}
	T_{B}^{\mathrm{s}}(z)j^{c}(0)\sim
	T_{B}(z)j^{c}(0) \:. \label{11.5.22}
\end{equation}

Now, replace $j^{c}(z_{3})$ with $T_{B}^{\mathrm{s}}$ and repeat the above discussion,
\begin{align}
	&j^{a}(z_{1})j^{a}(z_{2})T_{B}^{\mathrm{s}}(z_{3})=\frac{1}{z_{31}^{2}}j^{a}(z_{1})j^{a}(z_{2})
	+\frac{1}{z_{31}}\partial j^{a}(z_{1}) j^{a}(z_{2}) \nonumber \\
	&\quad +\frac{1}{z_{32}^{2}} j^{a}(z_{1})j^{a}(z_{2})+\frac{1}{z_{32}}j^{a}(z_{1})\partial j^{a}(z_{2})
	+\text{terms holomorphic in } z_{3} \:. \label{11.5.23}
\end{align}
Again expand in $z_{21}$ and keep the $z_{21}^{0}$ terms, which gives
\begin{equation}
	T_{B}^{\mathrm{s}}(z_{1})T_{B}^{\mathrm{s}}(z_{3}) \sim \frac{c^{g,k}}{2z_{13}^{4}}+\frac{2}{z_{13}^{2}}T_{B}^{\mathrm{s}}(z_{3})+\frac{1}{z_{13}}\partial T_{B}^{\mathrm{s}}(z_{3}) \label{11.5.24}
\end{equation}
where
\begin{equation}
	c^{g,k}=\frac{k\operatorname{dim}(g)}{k+h(g)} \:. \label{11.5.25}
\end{equation}
\begin{tcolorbox}[breakable]
	\begin{remark}
		Again set $\epsilon=z_{12}$ and write $j^{a}(z_{1})j^{a}(z_{2})=\frac{\hat{k}\operatorname{dim}(g)}{\epsilon^{2}}+A(z_{1},z_{2})$, where $A$ is regular with $A\to\,:\!jj(z_{1})\!:$ as $\epsilon\to0$. There is no $1/\epsilon$ term because $f^{aac}=0$. The four terms of (\ref{11.5.23}) are $\frac{1}{z_{31}^{2}}jj$, $\frac{1}{z_{31}}\partial_{1}(jj)$, $\frac{1}{z_{32}^{2}}jj$ and $\frac{1}{z_{32}}\partial_{2}(jj)$. Use
		\begin{equation*}
			\frac{1}{z_{32}}=\frac{1}{z_{31}}-\frac{\epsilon}{z_{31}^{2}}+\frac{\epsilon^{2}}{z_{31}^{3}}-\frac{\epsilon^{3}}{z_{31}^{4}}+\cdots\:,\qquad
			\frac{1}{z_{32}^{2}}=\frac{1}{z_{31}^{2}}-\frac{2\epsilon}{z_{31}^{3}}+\frac{3\epsilon^{2}}{z_{31}^{4}}+\cdots\:.
		\end{equation*}
		\textbf{c-number part.} With $\partial_{1}\epsilon^{-2}=-2\epsilon^{-3}$ and $\partial_{2}\epsilon^{-2}=+2\epsilon^{-3}$,
		\begin{align*}
			&\hat{k}\operatorname{dim}(g)\Bigl[\frac{1}{\epsilon^{2}z_{31}^{2}}-\frac{2}{\epsilon^{3}z_{31}}+\frac{1}{\epsilon^{2}z_{32}^{2}}+\frac{2}{\epsilon^{3}z_{32}}\Bigr]\\
			&\qquad\eqs{1}\hat{k}\operatorname{dim}(g)\Bigl[\frac{1+1-2}{\epsilon^{2}z_{31}^{2}}+\frac{-2+2}{\epsilon z_{31}^{3}}+\frac{3-2}{z_{31}^{4}}+O(\epsilon)\Bigr]
			=\frac{\hat{k}\operatorname{dim}(g)}{z_{13}^{4}}+O(\epsilon)\:.
		\end{align*}
		At (1) the expansions above were inserted and the powers of $\epsilon$ collected. The singular terms in $\epsilon$ cancel, as they must since they multiply nothing singular in $z_{3}$ after the subtraction.

		\textbf{Operator part.} At $\epsilon^{0}$, $A$ gives
		\begin{equation*}
			\frac{2:\!jj(z_{1})\!:}{z_{31}^{2}}+\frac{\partial\!:\!jj(z_{1})\!:}{z_{31}}
			\eqs{2}\frac{2:\!jj(z_{1})\!:}{z_{13}^{2}}-\frac{\partial\!:\!jj(z_{1})\!:}{z_{13}}
			\eqs{3}\frac{2:\!jj(z_{3})\!:}{z_{13}^{2}}+\frac{\partial\!:\!jj(z_{3})\!:}{z_{13}}\:.
		\end{equation*}
		At (2) $z_{31}=-z_{13}$. At (3) $:\!jj(z_{1})\!:$ was expanded about $z_{3}$, which turns $-1$ into $2-1=+1$ in the simple pole. Multiplying by $1/[(k+h(g))\psi^{2}]$ converts $:\!jj(z_{1})\!:$ into $T^{\mathrm{s}}_{B}(z_{1})$. Comparing with (\ref{11.5.24}),
		\begin{equation*}
			\frac{c^{g,k}}{2}=\frac{\hat{k}\operatorname{dim}(g)}{(k+h(g))\psi^{2}}\eqs{4}\frac{k\operatorname{dim}(g)}{2(k+h(g))}\:,
		\end{equation*}
		where at (4) $\hat{k}=\frac{1}{2}k\psi^{2}$. This is (\ref{11.5.25}). Equivalently, the $z^{-4}$ coefficient of $:\!jj\!:(z):\!jj\!:(0)$ is $(k+h(g))\psi^{2}\,\hat{k}\operatorname{dim}(g)$. A naive free-field double contraction would give only $2\hat{k}^{2}\operatorname{dim}(g)=k\psi^{2}\,\hat{k}\operatorname{dim}(g)$. The missing $h(g)\psi^{2}\,\hat{k}\operatorname{dim}(g)$ comes from the $ff$ terms, which is how $h(g)$ enters $c^{g,k}$.
	\end{remark}
\end{tcolorbox}
This is the standard form of an energy-momentum tensor with central charge $c^{g,k}$. Laurent coefficients are
\begin{subequations}
	\begin{align}
		L_{0}^{\mathrm{s}} &= \frac{1}{(k+h(g))\psi^{2}}\biggl(j_{0}^{a}j_{0}^{a}
		+2\sum_{n=1}^{\infty}j_{-n}^{a}j_{n}^{a}\biggr) \:, \label{11.5.26a} \\
		L_{m}^{\mathrm{s}} &= \frac{1}{(k+h(g))\psi^{2}}\sum_{n=-\infty}^{\infty}j_{n}^{a}j_{m-n}^{a}\:,\qquad m\neq0\:,
		\label{11.5.26b}
	\end{align} \label{11.5.26}
\end{subequations}
they also satisfy the Virasoro algebra with this central charge. Note that since holomorphicity requires $L_{0}^{\mathrm{s}}$ and $j_{n}^{a}$ with $n\geq0$ to annihilate state $\lvert 1\rangle$, the normal ordering constant in $L_{0}^{\mathrm{s}}$ is zero.
\begin{tcolorbox}[breakable]
	\begin{remark}
		The subtraction in (\ref{11.5.18}) removes exactly the singular part of $j^{a}(z)j^{a}(w)$, which is only a double pole. So
		\begin{equation*}
			:\!jj(w)\!:=\oint_{w}\frac{\dif z}{2\pi\mi}\,\frac{j^{a}(z)j^{a}(w)}{z-w}\:,
		\end{equation*}
		since $\frac{1}{(z-w)^{3}}$ has no residue. Deform the contour into $\lvert z\rvert>\lvert w\rvert$ minus $\lvert z\rvert<\lvert w\rvert$, where the radially ordered product is $j(z)j(w)$ and $j(w)j(z)$ respectively:
		\begin{align*}
			:\!jj(w)\!:&\eqs{1}\oint_{\lvert z\rvert>\lvert w\rvert}\frac{\dif z}{2\pi\mi}\sum_{l\geq0}\frac{w^{l}}{z^{l+1}}\,j^{a}(z)j^{a}(w)
			+\oint_{\lvert z\rvert<\lvert w\rvert}\frac{\dif z}{2\pi\mi}\sum_{l\geq0}\frac{z^{l}}{w^{l+1}}\,j^{a}(w)j^{a}(z)\\
			&\eqs{2}\sum_{n\leq-1}j^{a}_{n}w^{-n-1}\,j^{a}(w)+\sum_{n\geq0}j^{a}(w)\,j^{a}_{n}w^{-n-1}\:.
		\end{align*}
		At (1) $\frac{1}{z-w}$ was expanded in the appropriate region. At (2) $j^{a}(z)=\sum_{n}j^{a}_{n}z^{-n-1}$ was inserted. The residue picks $n=-l-1\leq-1$ in the first term and $n=l\geq0$ in the second. The coefficient of $w^{-m-2}$ is
		\begin{equation*}
			\bigl(:\!jj\!:\bigr)_{m}=\sum_{n\leq-1}j^{a}_{n}j^{a}_{m-n}+\sum_{n\geq0}j^{a}_{m-n}j^{a}_{n}\:.
		\end{equation*}
		For $m\neq0$ the order in the second sum does not matter: $[j^{a}_{m-n},j^{a}_{n}]=(m-n)\hat{k}\operatorname{dim}(g)\delta_{m,0}+\mi f^{aac}j^{c}_{m}=0$. So $(:\!jj\!:)_{m}=\sum_{n}j^{a}_{n}j^{a}_{m-n}$, which is (\ref{11.5.26b}). For $m=0$ the $n=0$ term is $j_{0}^{a}j_{0}^{a}$, and the terms $n\leq-1$ and $n\geq1$ combine into $2\sum_{n\geq1}j^{a}_{-n}j^{a}_{n}$, which is (\ref{11.5.26a}) with no additive constant.
	\end{remark}
\end{tcolorbox}

We have determined the $:jj::jj:$ OPE using the $jj$ OPE. we cannot do this directly because the $jj$ OPE holds if and only if the distance between the two operators is closer than the distance between any other operators; in this case, there are always two extra currents in the vicinity. Naively using the OPE will give wrong normalizations for $T^{\mathrm{s}}$ and $c^{g,k}$. The above discussion uses the OPE only when it holds, then takes advantage of holomorphicity. The operator $T_{B}^{\mathrm{s}}$ constructed from the product of two currents is called the \emph{Sugawara} energy-momentum tensor.

Finding the Sugawara tensor for the $U(1)$ current algebra is easy. Under normalization (\ref{11.5.14}), it is
\begin{equation}
	T_{B}^{\mathrm{s}} =\frac{1}{2}\,:jj: \:, \label{11.5.27}
\end{equation}
which can also be seen by writing it as the free bosonic current $j=\partial H$. 

Tensor $T_{B}^{\mathrm{s}}$ may or may not be equal to the total $T_{B}$ of the CFT. Define
\begin{equation}
	T_{B}^{\prime} = T_{B} - T_{B}^{\mathrm{s}} \:. \label{11.5.28}
\end{equation}
Since $T_{B}$ and $T_{B}^{\mathrm{s}}$ have the same singular terms in their OPEs with $j^{a}$, this gives
\begin{equation}
	T_{B}^{\prime}(z_{1})j^{a}(z_{2})\sim 0 . \label{11.5.29}
\end{equation}
Since $T_{B}^{\mathrm{s}}$ itself is constructed from currents, this implies $T_{B}^{\mathrm{s}}T_{B}^{\prime}\sim 0$. Then
\begin{align}
	T_{B}^{\prime}(z)T_{B}^{\prime}(0) &= T_{B}(z)T_{B}(0)-T_{B}^{\mathrm{s}}(z)T_{B}^{\mathrm{s}}(0)
	-T_{B}^{\prime}(z)T_{B}^{\mathrm{s}}(0) - T_{B}^{\mathrm{s}}(z)T_{B}^{\prime}(0) \nonumber \\
	&\sim \frac{c^{\prime}}{2z^{4}}+\frac{2}{z^{2}}T_{B}^{\prime}(0) +\frac{1}{z}\partial T_{B}^{\prime}(0)\:,
	\label{11.5.30}
\end{align}
which is a standard $TT$ OPE with central charge
\begin{equation}
	c^{\prime}=c-c^{g,k} \label{11.5.31}
\end{equation}

Thus the internal theory is divided into two decoupled CFTs. One possesses the energy-momentum tensor $T_{B}^{\mathrm{s}}$ constructed entirely from currents, while the other energy-momentum tensor $T_{B}^{\prime}$ commutes with the currents. Since the two CFTs are completely independent of each other, we will use the term \emph{current algebra} to refer to the first one. For a unitary CFT, $c^{\prime}$ must be non-negative, so
\begin{equation}
	c^{g,k}\leq c \:, \label{11.5.32}
\end{equation}
if
\begin{equation}
	c^{g,k}=c \:, \label{11.5.33}
\end{equation}
then $T_{B}^{\prime}$ is exactly trivial, in which case $T_{B}=T_{B}^{\mathrm{s}}$.

We now examine several examples. The dual Coxeter number can be written as a summation of roots. For any simply-laced algebra, $h(g)+1=\operatorname{dim}(g)/\operatorname{rank}(g)$, so
\begin{equation}
	c^{g,k} = \frac{k\operatorname{dim}(g)\operatorname{rank}(g)}{\operatorname{dim}(g)+(k-1)\operatorname{rank}(g)}
	\:. \label{11.5.34}
\end{equation}
For any simply-laced algebra at level $k=1$, the central charge is
\begin{equation}
	c^{g,1}=\operatorname{rank}(g) \:. \label{11.5.35}
\end{equation}
For $E_{8}\times E_{8}$ and $SO(32)$ heterotic strings, this is the same as the central charge of the free fermion representation, and for the free boson representation in the next section: they are Sugawara theories. The operator $:jj: $ looks like it is quadratic with respect to the fermions, but by using OPEs and the antisymmetry of fermions, it can be seen that $T_{B}^{\mathrm{s}}$ reduces to the usual $-\tfrac{1}{2}\lambda^{A}\partial\lambda^{A}$.
\begin{tcolorbox}[breakable]
	\begin{remark}
		\textbf{(\ref{11.5.34})--(\ref{11.5.35}).} For simply-laced $g$, $h(g)=\operatorname{dim}(g)/\operatorname{rank}(g)-1$ (box after Table \ref{tab:11.3}). Then
		\begin{equation*}
			c^{g,k}=\frac{k\operatorname{dim}(g)}{k-1+\operatorname{dim}(g)/\operatorname{rank}(g)}
			=\frac{k\operatorname{dim}(g)\operatorname{rank}(g)}{\operatorname{dim}(g)+(k-1)\operatorname{rank}(g)}
			\;\xrightarrow{\ k=1\ }\;\operatorname{rank}(g)\:.
		\end{equation*}
		For $SO(32)$, $c=16$, and for $E_{8}\times E_{8}$, $c=8+8=16$, the same as for 32 free fermions.

		\textbf{$T^{\mathrm{s}}_{B}$ for $SO(n)_{1}$.} Take the orthonormal currents $J^{AB}=\mi\lambda^{A}\lambda^{B}$, $A<B$. For one pair,
		\begin{align*}
			J^{AB}(z_{1})J^{AB}(z_{2})&=-\lambda^{A}_{1}\lambda^{B}_{1}\lambda^{A}_{2}\lambda^{B}_{2}
			\eqs{1}\lambda^{A}_{1}\lambda^{A}_{2}\;\lambda^{B}_{1}\lambda^{B}_{2}
			\eqs{2}\Bigl(\frac{1}{z_{12}}+:\!\lambda^{A}_{1}\lambda^{A}_{2}\!:\Bigr)\Bigl(\frac{1}{z_{12}}+:\!\lambda^{B}_{1}\lambda^{B}_{2}\!:\Bigr)\\
			&\eqs{3}\frac{1}{z_{12}^{2}}-\bigl(\lambda^{A}\partial\lambda^{A}+\lambda^{B}\partial\lambda^{B}\bigr)(z_{1})+O(z_{12})\:.
		\end{align*}
		At (1) $\lambda^{A}_{2}$ was moved past $\lambda^{B}_{1}$. At (2) the two species are independent free fields. At (3) $\lambda^{A}(z_{2})=\lambda^{A}(z_{1})-z_{12}\partial\lambda^{A}(z_{1})+\cdots$ gives $:\!\lambda^{A}_{1}\lambda^{A}_{2}\!:=-z_{12}\lambda^{A}\partial\lambda^{A}+O(z_{12}^{2})$, and the quartic term vanishes at coincidence. Summing over the $\frac{1}{2}n(n-1)$ pairs, each $\lambda^{C}\partial\lambda^{C}$ appears in $n-1$ of them:
		\begin{equation*}
			:\!jj\!:=-(n-1)\sum_{C}\lambda^{C}\partial\lambda^{C}\:,\qquad
			T^{\mathrm{s}}_{B}=\frac{:\!jj\!:}{(k+h)\psi^{2}}\eqs{4}\frac{-(n-1)\sum_{C}\lambda^{C}\partial\lambda^{C}}{2(n-1)}=-\frac{1}{2}\lambda^{C}\partial\lambda^{C}\:.
		\end{equation*}
		At (4) $k=1$, $h=n-2$ and $\psi^{2}=2$. Consistently, $c^{g,1}=\frac{\frac{1}{2}n(n-1)}{n-1}=\frac{n}{2}$, the central charge of $n$ free fermions.
	\end{remark}
\end{tcolorbox}

Another example is $SU(2)=SO(3)$, which has ($\operatorname{dim}(g)=3, \operatorname{rank}(g)=1$)
\begin{equation}
	c^{g,k}=\frac{3k}{2+k} = 1,\:\frac{3}{2},\:\frac{9}{5},\:2,\:\frac{15}{7},\cdots\to3 \:. \label{11.5.36}
\end{equation}
We have already seen the first CFT in this sequence (self-dual point of toroidal compactification) and the second (free fermions). Most levels do not have free-field representations. For any current algebra, the central charge lies in the range
\begin{equation}
	\operatorname{rank}(g)\leq c^{g,k} \leq \operatorname{dim}(g) \:. \label{11.5.37}
\end{equation}
The first equal sign holds only for level 1 simply-laced algebras, while the second holds only for Abelian algebras or in the $k\to\infty$ limit.
\begin{tcolorbox}[breakable]
	\begin{proof}
		The upper bound is immediate: $c^{g,k}=\operatorname{dim}(g)\,\frac{k}{k+h(g)}<\operatorname{dim}(g)$ for $h(g)>0$, and it tends to $\operatorname{dim}(g)$ as $k\to\infty$. A $U(1)$ current contributes exactly $1=\operatorname{dim}$. For the lower bound, use the root sum from the box after Table \ref{tab:11.3} and $\alpha^{2}\leq\psi^{2}$ for every root:
		\begin{equation*}
			h(g)\,\psi^{2}\operatorname{rank}(g)=\sum_{\alpha}\alpha^{2}\;\leq\;\psi^{2}\bigl(\operatorname{dim}(g)-\operatorname{rank}(g)\bigr)\:.
		\end{equation*}
		Then
		\begin{align*}
			c^{g,k}\geq\operatorname{rank}(g)
			\;&\Longleftrightarrow\; k\operatorname{dim}(g)\geq\bigl(k+h(g)\bigr)\operatorname{rank}(g)\\
			\;&\Longleftrightarrow\; k\bigl(\operatorname{dim}(g)-\operatorname{rank}(g)\bigr)\geq h(g)\operatorname{rank}(g)\:,
		\end{align*}
		which holds because $h(g)\operatorname{rank}(g)\leq\operatorname{dim}(g)-\operatorname{rank}(g)$ and $k\geq1$. Equality needs $k=1$ and $\alpha^{2}=\psi^{2}$ for all roots, i.e.\ a simply-laced algebra at level 1. For example, $SU(2)$ at $k=1,2,3,4,5$ gives $\frac{3k}{k+2}=1,\frac{3}{2},\frac{9}{5},2,\frac{15}{7}$, all between $1$ and $3$, which is (\ref{11.5.36}).
	\end{proof}
\end{tcolorbox}

\subsection*{Primary fields}

By repeatedly acting with lowering operators $j_{n}^{a}$ for $n>0$, we can obtain \emph{highest weight} states or \emph{primary} states of the current algebra, which are annihilated by all $j_{n}^{a}$ for $n>0$. Thus it is also annihilated by $L_{n}^{\mathrm{s}}$ in equation (\ref{11.5.26}) for $n>0$, and so is also a highest weight state of the Virasoro algebra. The center-of-mass generator $j_{0}^{a}$ transforms primary states into primary states, so the latter form a representation of algebra $g$,
\begin{equation}
	j_{0}^{a}\lvert r,i\rangle = \lvert r,j\rangle t_{r,ji}^{a} \label{11.5.38}
\end{equation}
where $r$ (without summation) labels the representation. It follows that
\begin{align}
	L_{0}^{\mathrm{s}}\lvert r,i\rangle &= \frac{1}{(k+h(g))\psi^{2}} \lvert r,k \rangle t_{r,kj}^{a}t_{r,ji}^{a} \nonumber\\ 
	&= \frac{Q_{r}}{(k+h(g))\psi^{2}} \lvert r,i \rangle \:, \label{11.5.39}
\end{align}
where $Q_{r}$ is the Casimir (\ref{11.4.7}). The weight of the primary state is thus determined by the algebra, level, and representation
\begin{equation}
	h_{r} = \frac{Q_{r}}{(k+h(g))\psi^{2}} = \frac{Q_{r}}{2\hat{k}+Q_{g}} \:. \label{11.5.40}
\end{equation}
where $Q_{g}$ is the Casimir of the adjoint representation. For $SU(2)$ at level $k$, the weight of the spin $j$ primary field is
\begin{equation}
	h_{j} = \frac{j(j+1)}{k+2} \:. \label{11.5.41}
\end{equation}
\begin{tcolorbox}[breakable]
	\begin{remark}
		\textbf{(\ref{11.5.39}).} On a primary state only the $j^{a}_{0}j^{a}_{0}$ term of (\ref{11.5.26a}) survives, because $j^{a}_{n}\lvert r,i\rangle=0$ for $n\geq1$. Applying (\ref{11.5.38}) twice,
		\begin{equation*}
			j^{a}_{0}j^{a}_{0}\lvert r,i\rangle=j^{a}_{0}\lvert r,j\rangle t^{a}_{r,ji}=\lvert r,k\rangle\,t^{a}_{r,kj}t^{a}_{r,ji}=\lvert r,k\rangle\,(t^{a}_{r}t^{a}_{r})_{ki}\eqs{1}Q_{r}\lvert r,i\rangle\:,
		\end{equation*}
		where (1) is (\ref{11.4.7}) with $d_{ab}=\delta_{ab}$.

		\textbf{Second form of (\ref{11.5.40}).} $(k+h(g))\psi^{2}=k\psi^{2}+h(g)\psi^{2}=2\hat{k}+Q_{g}$, using (\ref{11.5.13}) and $Q_{g}=h(g)\psi^{2}$ (box after Table \ref{tab:11.3}).

		\textbf{$SU(2)$.} With $\psi^{2}=2$ the root is $\alpha=\sqrt{2}$ in the one-dimensional Cartan space, so $J^{3}=\alpha H/\alpha^{2}=H/\sqrt{2}$, with $H$ unit normalized. The invariant inner product treats all three generators alike, so the orthonormal basis is $t^{a}=\sqrt{2}J^{a}$, where the $J^{a}$ are the usual spin matrices. Therefore
		\begin{equation*}
			Q_{j}=\sum_{a}t^{a}t^{a}=2\,\mathbf{J}^{2}=2j(j+1)\:,\qquad (k+h)\psi^{2}=2(k+2)\:,\qquad h_{j}=\frac{2j(j+1)}{2(k+2)}=\frac{j(j+1)}{k+2}\:.
		\end{equation*}
		As a check, $Q_{\text{adj}}=Q_{1}=4=h\psi^{2}$ with $h=2$.
	\end{remark}
\end{tcolorbox}


At any given level, there are at most finitely many representations for the primary state. For any root $\alpha$ of $g$ and any weight $\lambda$ of $r$, the $SU(2)$ algebra (\ref{11.5.12b}) implies that
\begin{align}
	\langle r,\lambda \rvert\, [E_{1}^{\alpha},E_{-1}^{-\alpha}]\,\lvert r,\lambda\rangle 
	&= 2\,\langle r,\lambda \rvert (\alpha \cdot H_{0}+\hat{k})\rvert r,\lambda \rangle/\alpha^{2} \nonumber \\
	&=2(\alpha \cdot \lambda + \hat{k})/\alpha^{2} \:. \label{11.5.42}
\end{align}
The left side is $\lVert E_{-1}^{-\alpha}\lvert r,\lambda\rangle \rVert^{2}\geq 0$, so $\hat{k}\geq -\alpha\cdot \lambda$. The same is true for $-\alpha$, which gives, for any weight $\lambda$ of $r$,
\begin{equation}
	\hat{k} \geq \lvert \alpha \cdot \lambda \rvert \:. \label{11.5.43}
\end{equation}
Taking $\alpha$ to be the long root $\psi$, then the level must satisfy
\begin{equation}
	k \geq \frac{2\lvert \psi \cdot \lambda \rvert}{\psi^{2}} = 2\lvert J^{3}\rvert \:, \label{11.5.44}
\end{equation}
where $J^{3}$ denotes the $SU(2)$ algebra (\ref{11.5.12a}) constructed from charge $j_{0}^{a}$ and root $\psi$. For $g=SU(2)$, the statement is that the spin $j$ of any primary state is at most $\tfrac{1}{2}k$. For example, at $k=1$, only representations $\mathbf{1}$ and $\mathbf{2}$ are possible. For $g=SU(3)$ at $k=1$, only $\mathbf{1}, \mathbf{3}$ and $\bar{\mathbf{3}}$ can appear. For $g=SU(n)$ at level $k$, the number of columns in the Young tableau of the representation that can appear is no more than $k$.
\begin{tcolorbox}[breakable]
	\begin{remark}
		For a simply-laced algebra every root is long, so (\ref{11.5.43}) with $\hat{k}=k$ reads $\lvert\alpha\cdot\lambda\rvert\leq k$ for every root $\alpha$ and every weight $\lambda$ of $r$.
		\begin{itemize}
			\item $SU(2)$: $\alpha=\sqrt{2}$ and $\lambda=\sqrt{2}\,m$ with $m=J^{3}\in\{-j,\dots,j\}$, so $\lvert\alpha\cdot\lambda\rvert=2\lvert m\rvert\leq k$, i.e.\ $j\leq\frac{k}{2}$.
			\item $SU(n)$ in the $U(n)$ basis, $\alpha=e_{i}-e_{j}$. A tableau with row lengths $r_{1}\geq\cdots\geq r_{n}$ has highest weight $\sum_{i}r_{i}e_{i}$. Full columns can be removed, so take $r_{n}=0$; the number of columns is then $r_{1}$. With $\alpha=e_{1}-e_{n}$, $\alpha\cdot\lambda=r_{1}-r_{n}=r_{1}$, so $r_{1}\leq k$. The $U(1)$ part of the weight drops out of $\alpha\cdot\lambda$ because $\alpha$ lies in the traceless hyperplane.
			\item $SU(3)$, $k=1$: the $\mathbf{3}$ has weights $e_{i}$ and $\bar{\mathbf{3}}$ has $-e_{i}$, with $\lvert\alpha\cdot\lambda\rvert\leq1$. The $\mathbf{6}$ has $2e_{1}$ with $(e_{1}-e_{2})\cdot2e_{1}=2$, and the $\mathbf{8}$ contains the root $e_{1}-e_{2}$ itself with $\alpha^{2}=2$. Both are excluded, and so is every larger representation, since it contains such a weight. Only $\mathbf{1},\mathbf{3},\bar{\mathbf{3}}$ remain.
		\end{itemize}
	\end{remark}
\end{tcolorbox}

The expectation value of a primary field is determined entirely by symmetry. We will give the details in Chapter 15.

Finally, we briefly discuss the gauge symmetry of Type I theory in equally abstract language. The matter part of the gauge vertex operator is the boundary
\begin{equation}
	\dot{X}^{\mu}\lambda^{a}\me^{\mi k\cdot X} \label{11.5.45}
\end{equation}
where $\lambda^{a}$ is a field with weight 0. In a unitary CFT, such $\lambda^{a}$ must be constructed using the equation of motion. Then the OPE is
\begin{equation}
	\lambda^{a}(y_{1})\lambda^{b}(y_{2}) = \Bigl[\theta(y_{1}-y_{2})d\indices{^a^b_c}+
	\theta(y_{2}-y_{1})d\indices{^b^a_c}\Bigr] \lambda^{c}(y_{2}) \:, \label{11.5.46}
\end{equation}
so $\lambda^{a}$ form a multiplicative algebra with structure constant $d\indices{^a^b_c}$. The antisymmetric part of $d\indices{^a^b_c}$ is the structure constant of the gauge Lie algebra. This is an abstract description of Chan-Paton factors. The requirement that the $\lambda^{a}$ algebra must be associative will rule out the gauge group $E_{8}\times E_{8}$.


\section{The Bosonic Construction and Toroidal Compactification} \label{sec:11.6}

When constructing the winding state vertex operators in section 8.2, we saw that we could consider left-moving and right-moving scalars separately. Now we try to construct a heterotic string with 26 left-movers and 10 right-moves, which together with $\psi^{\mu}$ gives the correct central charge. The main issue is the spectrum of $k_{L,R}$; as in section 8.4, in most of the discussion, we will use dimensionless momentum
\begin{equation}
	l_{L,R} = (\alpha^{\prime}/2)^{1/2} k_{L,R} \:. \label{11.6.1}
\end{equation}
Ordinary non-compact dimensions correspond to the left-move plus right-move with $l_{L}^{\mu}=l_{R}^{\mu}=l^{\mu}$, and they can take continuous values; let $d\leq 10$ be the number of non-compact dimensions. The remaining momenta,
\begin{equation}
	(l_{L}^{m},l_{R}^{n}) \:, \qquad d\leq m \leq 25 \:, \: d\leq n \leq 9 \:, \label{11.6.2}
\end{equation}
take values on lattice $\Gamma$. From the discussion of Narain compactification in section 8.4, we know that the requirements for a self-consistent CFT are OPE locality plus modular invariance. After taking the GSO projection on the right-movers, the conditions on $\Gamma$ are precisely the same as in the bosonic case. Define the product
\begin{equation}
	l \circ l^{\prime} = l_{L}\cdot l_{L}^{\prime} - l_{R} \cdot l_{R}^{\prime} \:, \label{11.6.3}
\end{equation}
the lattice must be an \emph{even self-dual Lorentzian lattice with signature $(26-d,10-d)$},
\begin{subequations}
	\begin{align}
		l \circ l & \in 2\mathds{Z} \qquad \text{for all } l\in\Gamma \:, \label{11.6.4a} \\
		\Gamma &= \Gamma^{\ast} \:. \label{11.6.4b}
	\end{align} \label{11.6.4}
\end{subequations}

As in the bosonic case, where the signature was $(26-d,26-d)$, all such lattices have been classified. First examine the maximum number of non-compact dimensions, $d=10$; in this case, the $\circ$ product has only a positive sign, so $l_{L}^{m}$ forms a 16-dimensional even self-dual Euclidean lattice. Even self-dual Euclidean lattices exist only when the dimension is a multiple of 8; for 16 dimensions, there are precisely two such lattices, $\Gamma_{16}$ and $\Gamma_{8}\times \Gamma_{8}$. Lattice $\Gamma_{16}$ is the set of all points of the following form
\begin{subequations}
	\begin{align}
		&(n_{1},\cdots,n_{16})\quad \text{or} \quad (n_{1}+\tfrac{1}{2},\cdots ,n_{16}+\tfrac{1}{2}) \:, \label{11.6.5a} \\
		&\sum_{i}n_{i} \in 2\mathds{Z} \:, \label{11.6.5b}
	\end{align}  \label{11.6.5}
\end{subequations}
where $n_{i}$ is any integer. Similarly, lattice $\Gamma_{8}$ is
\begin{subequations}
	\begin{align}
		&(n_{1},\cdots,n_{8})\quad \text{or} \quad (n_{1}+\tfrac{1}{2},\cdots ,n_{8}+\tfrac{1}{2}) \:, \label{11.6.6a} \\
		&\sum_{i}n_{i} \in 2\mathds{Z} \:. \label{11.6.6b}
	\end{align} \label{11.6.6}
\end{subequations}
\begin{tcolorbox}[breakable]
	\begin{remark}
		Treat both lattices at once as $\Gamma_{N}$, $N=8$ or $16$: $D_{N}=\{n\in\mathds{Z}^{N}:\sum n_{i}\in2\mathds{Z}\}$ together with $D_{N}+s$, where $s=(\frac{1}{2})^{N}$. Write points as $n$ or $m+s$ with $n,m\in D_{N}$.

		\textbf{Even.} Using $n_{i}^{2}\equiv n_{i}\pmod 2$,
		\begin{align*}
			n^{2}&=\sum n_{i}^{2}\equiv\sum n_{i}\equiv0\pmod2\:,\\
			(m+s)^{2}&=\sum_{i}\bigl(m_{i}^{2}+m_{i}\bigr)+\frac{N}{4}\eqs{1}\frac{N}{4}\pmod 2\:.
		\end{align*}
		At (1) $m_{i}(m_{i}+1)$ is even. So the half-integer points are even iff $N/4\in2\mathds{Z}$, i.e.\ $N\in8\mathds{Z}$. This is why these lattices exist in dimensions 8 and 16.

		\textbf{Integral.} $n\cdot n'\in\mathds{Z}$. Next, $n\cdot(m+s)=n\cdot m+\frac{1}{2}\sum n_{i}\in\mathds{Z}$ since $\sum n_{i}$ is even. Finally $(m+s)\cdot(m'+s)=m\cdot m'+\frac{1}{2}\bigl(\sum m_{i}+\sum m'_{i}\bigr)+\frac{N}{4}\in\mathds{Z}$. So $\Gamma_{N}\subset\Gamma_{N}^{\ast}$. The set is closed under addition because $2s=(1^{N})\in D_{N}$ for even $N$.

		\textbf{Self-dual.} $D_{N}$ has index 2 in $\mathds{Z}^{N}$, so its unit cell has volume 2. Adding the coset $D_{N}+s$ doubles the density:
		\begin{equation*}
			\operatorname{vol}(\Gamma_{N})=\tfrac{1}{2}\operatorname{vol}(D_{N})=1=\operatorname{vol}(\Gamma_{N}^{\ast})\:.
		\end{equation*}
		Since $\Gamma_{N}\subset\Gamma_{N}^{\ast}$ and the index is $\operatorname{vol}(\Gamma_{N})/\operatorname{vol}(\Gamma_{N}^{\ast})=1$, we get $\Gamma_{N}=\Gamma_{N}^{\ast}$.

		\textbf{Length-$\sqrt{2}$ vectors.} Integer points with $n^{2}=2$ are $(\pm1,\pm1,0^{N-2})$, $4\binom{N}{2}$ of them. Half-integer points have $(m+s)^{2}\geq N/4$.
		\begin{align*}
			\Gamma_{16}&:\quad 4\tbinom{16}{2}=480\ \text{integer vectors, no half-integer ones }(16/4=4>2)\:,\\
			&\quad 480+16=496=\operatorname{dim}SO(32)\:,\\
			\Gamma_{8}&:\quad 4\tbinom{8}{2}+2^{7}=112+128=240\:,\quad 240+8=248=\operatorname{dim}E_{8}\:.
		\end{align*}
		For $\Gamma_{8}$ the half-integer vectors are all $(\pm\frac{1}{2})^{8}$ with an even number of minus signs, as in the box after (\ref{11.4.21}). Both $\Gamma_{16}$ and $\Gamma_{8}\times\Gamma_{8}$ have $480$ roots and give $496$-dimensional gauge groups.
	\end{remark}
\end{tcolorbox}


The left-moving zero-point energy is $-1$ as in the bosonic string, so massless states will have left-moving vertex operators $\partial X^{\mu}, \partial X^{m}$ or $\me^{\mi k_{L}\cdot X(z)}$ satisfying $l_{L}^{2}=2$. Doing the tensor product with the usual right-moving $\mathbf{8}_{v}+\mathbf{8}$, the first gives the usual graviton, dilaton, and antisymmetric tensor. 16 $\partial X^{m}$ currents form the largest commuting set of $m$-momenta. Momenta $l_{L}^{m}$ are the charges under this current, so they are the roots of the gauge group. For $\Gamma_{16}$, points with length $\sqrt{2}$ are the roots of $SO(32)$ (\ref{11.4.16}). For $\Gamma_{8}$, points with length $\sqrt{2}$ are the roots of $E_{8}$ (\ref{11.4.21}). Thus the two possible lattices exactly give the two gauge groups found previously, $SO(32)$ and $E_{8}\times E_{8}$. The commuting current has singularity $1/z_{2}$, so again $k=1$.

It is easy to see that the previous fermi construction and the current bosonic construction are equivalent under bosonization. The integer points on the lattices (\ref{11.6.5}) and (\ref{11.6.6}) map to the NS sector of the current algebra, and half-integer points map to the R sector. The constraint that the total $k_{L}^{m}$ must be even is the GSO projection on the left-moving side in each theory. As we saw in the previous section, the dynamics of the current algebra are determined entirely by its symmetry, so we can give a representation-independent description of the left-moving side, describing it as a level 1 current algebra of $SO(32)$ or $E_{8}\times E_{8}$.

Let's review some general results for Lie algebras and lattices. The set of all integer-coefficient linear combinations of the roots of Lie algebra $g$ is called the \emph{root lattice} $\Gamma_{g}$ of $g$. Now take any representation $r$ and let $\lambda$ be any weight of $r$. For all $v\in\Gamma_{g}$, the set of points $\lambda+v$ is denoted $\Gamma_{r}$. By examining various $SU(2)$ subgroups it can be proved that, for simply-laced Lie algebras with $\sqrt{2}$ root length,
\begin{equation}
	\Gamma_{r} \subset \Gamma_{g}^{\ast} \:. \label{11.6.7}
\end{equation}
The union of all $\Gamma_{r}$ is the weight lattice $\Gamma_{w}$, and
\begin{equation}
	\Gamma_{w}=\Gamma_{g}^{\ast} \:.
\end{equation}
\begin{tcolorbox}[breakable]
	\begin{remark}
		For each root $\alpha$ the generators (\ref{11.5.10}) form an $SU(2)$, so on any weight $\lambda$ of a representation $2J^{3}=2\alpha\cdot\lambda/\alpha^{2}$ is an integer. With $\alpha^{2}=2$,
		\begin{equation*}
			\alpha\cdot\lambda\in\mathds{Z}\quad\text{for every root }\alpha
			\quad\Longrightarrow\quad v\cdot\lambda\in\mathds{Z}\quad\text{for every }v=\textstyle\sum_{\alpha}n_{\alpha}\alpha\in\Gamma_{g}\:,
		\end{equation*}
		which is $\lambda\in\Gamma_{g}^{\ast}$. Adding a root to a weight keeps it in $\Gamma_{g}^{\ast}$ because the roots themselves are integral ($\alpha\cdot\beta\in\mathds{Z}$ by the same argument in the adjoint), so $\Gamma_{r}\subset\Gamma_{g}^{\ast}$. Conversely, every point of $\Gamma_{g}^{\ast}$ is the highest weight of some representation, up to a Weyl reflection, which is the content of $\Gamma_{w}=\Gamma_{g}^{\ast}$. For $SO(2n)$ one checks directly that the four classes (\ref{11.6.9}) exhaust $\Gamma^{\ast}_{D_{n}}$. A vector $x$ has $x\cdot(e_{i}\pm e_{j})\in\mathds{Z}$ for all $i\neq j$ iff all $x_{i}\in\mathds{Z}$ or all $x_{i}\in\mathds{Z}+\frac{1}{2}$. Modulo the root lattice (integer vectors with even sum), the integer vectors fall into the classes $(0)$ and $(v)$ according to the parity of $\sum x_{i}$, and the half-integer ones into $(s)$ and $(c)$ in the same way.
	\end{remark}
\end{tcolorbox}
For example, the weight lattice of $SO(2n)$ has 4 sublattices:
\begin{subequations}
	\begin{align}
		&(0)\::\quad 0+\text{any root}\:; \label{11.6.9a} \\
		&(v)\::\quad (1,0,0,\cdots,0)+\text{any root}\:; \label{11.6.9b} \\
		&(s)\::\quad (\tfrac{1}{2},\tfrac{1}{2},\tfrac{1}{2},\cdots,\tfrac{1}{2})+\text{any root}\:; \label{11.6.9c} \\
		&(c)\::\quad (-\tfrac{1}{2},\tfrac{1}{2},\tfrac{1}{2},\cdots,\tfrac{1}{2})+\text{any root}\:; \label{11.6.9d} 
	\end{align} \label{11.6.9}
\end{subequations}
they are the root lattice, the lattice containing weights of the vector representation, and two lattices containing weights of the $2^{n-1}$-dimensional spinor representations, respectively. Lattice $\Gamma_{8}$ is the root lattice of $E_{8}$, and simultaneously, since it is self-dual, it is the weight lattice. $SO(32)$ root lattices give integer lattices in $\Gamma_{16}$. The entire $\Gamma_{16}$ is the root lattice plus a spinor lattice of $SO(32)$.

For any simply-laced Lie algebra $g$, its level 1 current algebra can be represented by $\operatorname{rank}(g)$ left-moving bosons; as the roots are rescaled to length $\sqrt{2}$, the momentum lattice becomes the root lattice of $g$. The constant $\epsilon(\alpha,\beta)$ appearing in Lie algebra (\ref{11.4.12}) can be determined by vertex operator OPEs; such is the case when an explicit form of the cocycle is needed. By simultaneously taking $\operatorname{rank}(g)$ right-moving bosons, we can obtain a modularly invariant CFT, and the momentum lattice becomes
\begin{equation}
	\Gamma=\sum_{r}\Gamma_{r}\times \tilde{\Gamma}_{r} \:. \label{11.6.10}
\end{equation}
That is, the spectrum takes all sublattices of the weight lattice, and at the same time, left-moving and right-moving momenta are taken within the same sublattice.
\begin{tcolorbox}[breakable]
	\begin{proof}
		Pick a representative $\lambda_{r}$ in each class, so a general point is $l=(\lambda_{r}+v,\lambda_{r}+\tilde{v})$ with $v,\tilde{v}\in\Gamma_{g}$. The root lattice of a simply-laced algebra with $\alpha^{2}=2$ is even, since $\alpha\cdot\beta\in\mathds{Z}$ and $\alpha^{2}=2$. Also $\lambda_{r}\cdot v\in\mathds{Z}$ by (\ref{11.6.7}).

		\textbf{Even.}
		\begin{equation*}
			l\circ l=(\lambda_{r}+v)^{2}-(\lambda_{r}+\tilde{v})^{2}=2\lambda_{r}\cdot(v-\tilde{v})+v^{2}-\tilde{v}^{2}\in2\mathds{Z}\:.
		\end{equation*}
		\textbf{Integral.} For $l'=(\lambda_{r'}+v',\lambda_{r'}+\tilde{v}')$,
		\begin{equation*}
			l\circ l'=\lambda_{r}\cdot\lambda_{r'}-\lambda_{r}\cdot\lambda_{r'}+(\text{products involving at least one root vector})\in\mathds{Z}\:.
		\end{equation*}
		\textbf{Unimodular.} Let $N=\lvert\Gamma_{g}^{\ast}/\Gamma_{g}\rvert$ be the number of classes. Then $\operatorname{vol}(\Gamma_{g})=N\operatorname{vol}(\Gamma_{g}^{\ast})=N/\operatorname{vol}(\Gamma_{g})$, so $\operatorname{vol}(\Gamma_{g})^{2}=N$. The lattice $\Gamma$ is $N$ translates of $\Gamma_{g}\times\Gamma_{g}$, so
		\begin{equation*}
			\operatorname{vol}(\Gamma)=\frac{\operatorname{vol}(\Gamma_{g})^{2}}{N}=1\:.
		\end{equation*}
		An integral lattice of unit volume equals its dual, so $\Gamma=\Gamma^{\ast}$. For example, $SU(2)$ has $N=2$ classes (integer and half-integer spin) and $\Gamma$ is the momentum lattice of a circle at the self-dual radius, where the $SU(2)\times SU(2)$ symmetry appears. For $SO(2n)$, $N=4$ by (\ref{11.6.9}).
	\end{proof}
\end{tcolorbox}

\subsection*{Toroidal compactification}

As in the bosonic case, any even self-dual lattice with signature $(26-d,10-d)$ can be obtained from any single lattice via $O(26-d,10-d,\mathds{R})$. Again, start from any given solution $\Gamma_{0}$; for example, it can be a ten-dimensional theory where all compactified dimensions are orthogonal and reside at $SU(2)\times SU(2)$ radii. Then any lattice
\begin{equation}
	\Gamma=\Lambda\Gamma_{0} \:,\qquad \Lambda\in O(26-d,10-d,\mathds{R}) \label{11.6.11}
\end{equation}
defines a self-consistent heterotic string theory. As in the bosonic case, there exists an equivalence relation
\begin{equation}
	\Lambda_{1}\Lambda\Lambda_{2}\Gamma_{0} \cong \Lambda\Gamma_{0} \label{11.6.12}
\end{equation}
where
\begin{equation}
	\Lambda_{1} \in O(26-d,\mathds{R})\times O(10-d,\mathds{R}) \:,\qquad
	\Lambda_{2} \in O(26-d,10-d,\mathds{Z}) \:. \label{11.6.13}
\end{equation}
Then the moduli space is
\begin{equation}
	\frac{ O(26-d,10-d,\mathds{R})}{ O(26-d,\mathds{R})\times O(10-d,\mathds{R})\times  O(26-d,10-d,\mathds{Z})}\:.
	\label{11.6.14}
\end{equation}
The discrete $T$-duality group $O(26-d,10-d,\mathds{Z})$, which remains invariant on $\Gamma_{0}$, is understood to act on the right.

Now examine the unbroken gauge symmetry. There are $26{-}d$ gauge bosons with vertex operators $\partial X^{m}\tilde{\psi}^{\mu}$ and $10{-}d$ gauge bosons with vertex operators $\partial X^{\mu}\tilde{\psi}^{m}$. They are the original 16 commuting symmetries from the ten-dimensional theory and $10-d$ Kaluza-Klein gauge bosons, plus the compactification of another $10-d$ antisymmetric tensors. For each point on lattice $\Gamma$ satisfying
\begin{equation}
	l_{L}^{2}=2 \:, \qquad l_{R}=0 \:, \label{11.6.15}    
\end{equation}
there are additional gauge bosons $\me^{\mi k_{L}\cdot X_{L}}\tilde{\psi}_{\mu}$. There are no gauge bosons at points where $l_{R}\neq 0$, because the mass of such a state is at least $\frac{1}{2}l_{R}^{2}$. For a general boost $\Lambda$, which gives a general point in moduli space, there are no points with $l_{R}=0$ in $\Gamma$, and thus no extra gauge bosons; the gauge group is $U(1)^{36-2d}$. By compactifying the original ten-dimensional theory on a torus without Wilson lines, one can clearly obtain $SO(32)\times U(1)^{20-2d}$ or $E_{8}\times E_{8}\times U(1)^{20-2d}$, as in field theory. However, as with bosonic string theory, there are points in the moduli space where stringy expanded gauge symmetry exists. For example, the lattice $\Gamma_{26-d,10-d}$ defined following lattices $\Gamma_{8}$ and $\Gamma_{16}$ will give $SO(52-2d)\times U(1)^{10-d}$. Similar to the bosonic case, the low-energy physics near the points where expanded symmetry exists is the Higgs mechanism. The rank of all groups obtained in this way is $36-2d$. This is the maximum value in perturbation theory, but non-perturbative effects lead to larger gauge symmetries (see chapter 19).

The number of moduli comes from the dimension of the $SO$ group, which is
\begin{equation}
	\frac{1}{2}\Bigl[(36-2d)(35-2d)-(26-d)(25-d)-(10-d)(9-d)\Bigr]=(26-d)(10-d) \:. \label{11.6.16}
\end{equation}
\begin{tcolorbox}[breakable]
	\begin{remark}
		The discrete group does not change the dimension, so the dimension of (\ref{11.6.14}) is $\operatorname{dim}O(p,q)-\operatorname{dim}O(p)-\operatorname{dim}O(q)$ with $p=26-d$, $q=10-d$ and $p+q=36-2d$:
		\begin{equation*}
			\tfrac{1}{2}\bigl[(p+q)(p+q-1)-p(p-1)-q(q-1)\bigr]\eqs{1}\tfrac{1}{2}\bigl[(p+q)^{2}-p^{2}-q^{2}\bigr]=pq\:.
		\end{equation*}
		At (1) the linear terms cancel, $-(p+q)+p+q=0$. The count matches the backgrounds: $G_{mn}$ and $B_{mn}$ give $\frac{1}{2}q(q+1)+\frac{1}{2}q(q-1)=q^{2}$, and the Wilson lines $A^{I}_{m}$ give $16q$. In total $q^{2}+16q=q(q+16)=(10-d)(26-d)$.
	\end{remark}
\end{tcolorbox}
As in the bosonic string, they can be interpreted as background fields of a ten-dimensional gauge theory. The compactified components of the metric and antisymmetric tensor altogether give $(10-d)^{2}$, as before. Additionally, there can be Wilson lines, namely constant backgrounds of gauge field $A_{m}$. As discussed in Chapter 8, due to the potential $\operatorname{Tr}([A_{m},A_{n}]^{2})$, fields in different directions commute along flat directions, so they can be chosen to reside in a $U(1)^{16}$ subgroup. Thus, there are $16(10-d)$ parameters for $A_{m}$ in total $(26-d)(10-d)$.

In Chapter 8, we studied the quantization with antisymmetric tensor and open string Wilson line backgrounds. Details are omitted here and the results are given directly. If we do the compactification $x^{m}\cong x^{m}+2\pi R$ under constant backgrounds $G_{mn}$, $B_{mn}$ and $A_{m}^{I}$, canonical quantization gives
\begin{subequations}
	\begin{align}
		k_{Lm}&= \frac{n_{m}}{R}+\frac{w^{n}R}{\alpha^{\prime}}(G_{mn}+B_{mn})-q^{I}A_{m}^{I}-\frac{w^{n}R}{2}A_{n}^{I}A_{m}^{I}\:,\label{11.6.17a}\\
		k_{L}^{I} &= (q^{I}+w^{m}RA_{m}^{I})(2/\alpha^{\prime})^{1/2} \:, \label{11.6.17b} \\
		k_{Rm}&= \frac{n_{m}}{R}+\frac{w^{n}R}{\alpha^{\prime}}(-G_{mn}+B_{mn})-q^{I}A_{m}^{I}-\frac{w^{n}R}{2}A_{n}^{I}A_{m}^{I}\:,\label{11.6.17c}
	\end{align}
\end{subequations}
where $n_{m}$ and $w^{m}$ are integers depending on which type of string is compactified, and $q^{I}$ resides on $\Gamma_{16}$ or $\Gamma_{8}\times\Gamma_{8}$. Note that when the gauge field is set to zero, this result reduces to the bosonic result (\textcolor{blue}{8.4.7}). The linear term of $A^{I}$ in $k_{Lm}$ and $k_{Rm}$ comes from the effect of Wilson lines on periodicity, the same as (\textcolor{blue}{8.6.7}). The source of the linear term in $A^{I}$ in $k_{L}^{I}$ is as follows. For a string winding along a compactified dimension, Wilson lines mean current algebra fermions are no longer periodic. The corresponding vertex algebra (\ref{10.3.25}) indicates that momentum is shifted. Finally, the quadratic term of $A^{I}$ can be easily verified by checking that $\alpha^{\prime}k\circ k/2$ is even.
\begin{tcolorbox}[breakable]
	\begin{remark}
		Write (\ref{11.6.17a}) and (\ref{11.6.17c}) as $k_{Lm}=P_{m}+W_{m}$ and $k_{Rm}=P_{m}-W_{m}$ with
		\begin{equation*}
			P_{m}=\frac{n_{m}}{R}+\frac{w^{n}R}{\alpha'}B_{mn}-q^{I}A^{I}_{m}-\frac{w^{n}R}{2}A^{I}_{n}A^{I}_{m}\:,\qquad W_{m}=\frac{w^{n}R}{\alpha'}G_{mn}\:,
		\end{equation*}
		and define $a^{I}=w^{m}RA^{I}_{m}$. Indices are contracted with $G^{mn}$. Then
		\begin{align*}
			k_{Lm}k_{L}^{m}-k_{Rm}k_{R}^{m}&=4P_{m}W^{m}\eqs{1}\frac{4R}{\alpha'}w^{m}P_{m}\\
			&\eqs{2}\frac{4}{\alpha'}\Bigl[n_{m}w^{m}-q^{I}a^{I}-\tfrac{1}{2}a^{I}a^{I}\Bigr]\:,\\
			k_{L}^{I}k_{L}^{I}&=\frac{2}{\alpha'}\bigl(q^{I}+a^{I}\bigr)^{2}=\frac{2}{\alpha'}\bigl[q^{2}+2q\cdot a+a^{2}\bigr]\:.
		\end{align*}
		At (1) $W^{m}=G^{mn}W_{n}=w^{m}R/\alpha'$. At (2) $w^{m}w^{n}B_{mn}=0$ by antisymmetry. Adding,
		\begin{equation*}
			\frac{\alpha'}{2}\,k\circ k=\frac{\alpha'}{2}\bigl(k_{L}^{2}+k^{I}_{L}k^{I}_{L}-k_{R}^{2}\bigr)
			=2n_{m}w^{m}-2q\cdot a-a^{2}+q^{2}+2q\cdot a+a^{2}=2n_{m}w^{m}+q^{2}\in2\mathds{Z}\:,
		\end{equation*}
		because $q$ lies in the even lattice $\Gamma_{16}$ or $\Gamma_{8}\times\Gamma_{8}$. The linear terms in $A$ cancel between $k_{Lm}$ and $k_{L}^{I}$, and the quadratic term $-\frac{1}{2}w^{n}RA^{I}_{n}A^{I}_{m}$ is exactly what cancels the $a^{2}$ from $(k_{L}^{I})^{2}$. Without it $k\circ k$ would depend continuously on $A$ and could not be even.
	\end{remark}
\end{tcolorbox}

To compare this spectrum with the Narain description, we must go to coordinates where $G_{m^{\prime}n^{\prime}}=\delta_{m^{\prime}n^{\prime}}$, which makes $k_{m^{\prime}}=e_{m^{\prime}}{}^{n}k_{n}$, the frame being defined as $\delta_{p^{\prime}q^{\prime}}=e_{p^{\prime}}{}^{m}e_{q^{\prime}}{}^{n}G_{mn}$. The discrete $T$-duality group is generated by $T$-dualities on each axis, finite large spacetime coordinate transformations, and quantum shifts of antisymmetric tensor background and Wilson lines.

There is one interesting point here. Since coset space (\ref{11.6.14}) is the general solution for compatibility conditions, whether we compactify the $SO(32)$ theory or the $E_{8}\times E_{8}$ theory, we must obtain the same family of theories. From another perspective, note that the coset space is non-compact due to its Lorentzian signature—we can go to the limit of infinite Narain boost. One such limit yields the ten-dimensional $SO(32)$ theory, while another yields the ten-dimensional $E_{8}\times E_{8}$ theory. Clearly, we can consider all different toroidally compactified heterotic strings as different states in the same theory. Then these two ten-dimensional theories are two different limits of this theory.

We now make the connection between these theories more explicit. Compactifying the $SO(32)$ theory on a circle of radius $R$, let this circle have $G_{99}=1$ and a Wilson line
\begin{equation}
	RA_{9}^{I}=\operatorname{diag}\Bigl(\tfrac{1}{2}^{8},0^{8}\Bigr)
	\label{11.6.18}
\end{equation}
Adjoint states where one index takes $1\leq A\leq 16$ and the other index takes $17\leq A\leq 32$ are antiperiodic due to Wilson lines, so the gauge symmetry breaks to $SO(16)\times SO(16)$. Now examine the $E_{8}\times E_{8}$ theory, this time with radius of the circle $R^{\prime}$, and $G_{99}=1$ and a Wilson line
\begin{equation}
	R^{\prime}A_{9}^{I}=\operatorname{diag}\Bigl(1,0^{7},1,0^{7}\Bigr) \:.
	\label{11.6.19}
\end{equation}
In each $E_{8}$, integer-charge states from the $SO(16)$ root lattice remain periodic, but half-integer states from the $SO(16)$ spinor lattice become antiperiodic. As above, gauge symmetry remains $SO(16)\times SO(16)$. To see the connection between these two theories, we only focus on states that are neutral under $SO(16)\times SO(16)$, namely those with $k_{L}^{I}=0$. In these two theories, due to the shift in $k_{L}^{I}$, they only manifest for even $w=2m$. The two neutral spectra are
\begin{equation}
	k_{L,R}=\frac{\tilde{n}}{R} \pm \frac{2mR}{\alpha^{\prime}}\:,\qquad k_{L,R}^{\prime} =\frac{\tilde{n}^{\prime}}{R^{\prime}} \pm \frac{2m^{\prime}R^{\prime}}{\alpha^{\prime}}\:,
\end{equation}
we have not written out the index 9. The prime refers to the $E_{8}\times E_{8}$ theory, and we have $\tilde{n}=n+2m, \tilde{n}^{\prime}=n^{\prime}+2m^{\prime}$. We used the explicit forms of Wilson lines and $k_{L}^{I}=0$ in each case. Under $(\tilde{n},m)\leftrightarrow (m^{\prime},\tilde{n}^{\prime})$ and $(k_{L},k_{R})\leftrightarrow (k_{L}^{\prime},-k_{R}^{\prime})$, if $RR^{\prime}=\alpha^{\prime}/2$, then these two spectra are identical. This symmetry extends to the entire spectrum.
\begin{tcolorbox}[breakable]
	\begin{remark}
		\textbf{Symmetry breaking.} A state of charge $q$ picks up $\exp(2\pi\mi\,q\cdot a)$ around the circle, with $a=RA_{9}$. For $SO(32)$, $a=((\frac{1}{2})^{8},0^{8})$. The roots $\pm e_{i}\pm e_{j}$ with $i,j$ in the same block have $q\cdot a\in\{0,\pm1\}$, while those with one index in each block have $q\cdot a=\pm\frac{1}{2}$. The latter are antiperiodic, which leaves $SO(16)\times SO(16)$. For $E_{8}\times E_{8}$, $a'=(1,0^{7},1,0^{7})$. The $SO(16)$ roots in each $E_{8}$ have $q\cdot a'\in\mathds{Z}$, while the spinor weights $(\pm\frac{1}{2})^{8}$ have $q\cdot a'=\pm\frac{1}{2}$, so each $E_{8}$ is broken to $SO(16)$.

		\textbf{Neutral states, $SO(32)$.} With one compact direction, $G_{99}=1$, $B=0$: $k_{L}^{I}\propto q+wa=0$ requires $q=-wa=((-\frac{w}{2})^{8},0^{8})$. For odd $w$ this vector mixes half-integer and integer entries and is not in $\Gamma_{16}$, so $w=2m$ and $q=((-m)^{8},0^{8})\in\Gamma_{16}$. Then
		\begin{align*}
			k_{L,R}&=\frac{n}{R}\pm\frac{wR}{\alpha'}-q\cdot A-\frac{wR}{2}A\cdot A
			\eqs{1}\frac{n}{R}\pm\frac{2mR}{\alpha'}+\frac{4m}{R}-\frac{2m}{R}=\frac{n+2m}{R}\pm\frac{2mR}{\alpha'}\:.
		\end{align*}
		At (1) $q\cdot A=q\cdot a/R=-8\cdot\frac{m}{2}/R=-4m/R$ and $A\cdot A=a^{2}/R^{2}=2/R^{2}$.

		\textbf{Neutral states, $E_{8}\times E_{8}$.} Now $q=-w'a'=(-w',0^{7},-w',0^{7})$, which lies in $\Gamma_{8}\times\Gamma_{8}$ iff $w'$ is even (each block needs an even coordinate sum), so $w'=2m'$. With $q\cdot a'=-4m'$ and $a'^{2}=2$ the same computation gives $k'_{L,R}=\frac{n'+2m'}{R'}\pm\frac{2m'R'}{\alpha'}$.

		\textbf{Duality.} Set $R'=\alpha'/2R$ and map $(\tilde{n}',m')=(m,\tilde{n})$:
		\begin{equation*}
			k'_{L}=\frac{2R\tilde{n}'}{\alpha'}+\frac{2m'}{\alpha'}\frac{\alpha'}{2R}=\frac{2mR}{\alpha'}+\frac{\tilde{n}}{R}=k_{L}\:,\qquad
			k'_{R}=\frac{2mR}{\alpha'}-\frac{\tilde{n}}{R}=-k_{R}\:.
		\end{equation*}
		Both $(\tilde{n},m)$ and $(\tilde{n}',m')$ range over all of $\mathds{Z}^{2}$, so the map is a bijection of the neutral spectra, and $k_{R}\to-k_{R}$ is the usual T-duality reflection of the right movers.
	\end{remark}
\end{tcolorbox}

Finally, let's see how realistic the theories obtained via compactification to 4 dimensions are. At a general point in moduli space, the massless spectrum is obtained via dimensional reduction, and states are classified by their 4-dimensional symmetries. Analyzing this spectrum using 4-dimensional $SO(2)$ helicity, the $SO(8)$ spin decomposes into
\begin{subequations}
	\begin{align}
		\mathbf{8}_{v} &\to +1,\,0^{6},\,-1\:, \label{11.6.21a}  \\
		\mathbf{8} &\to +\tfrac{1}{2}^{4},\,-\tfrac{1}{2}^{4} \:, \label{11.6.21b}
	\end{align}
\end{subequations}  
thus
\begin{subequations}
	\begin{align}
		\mathbf{8}_{v}\times\mathbf{8}_{v} &\to +2,\,+1^{12},\,0^{38},\,-1^{12},\,-2\:, \label{11.6.22a} \\
		\mathbf{8}\times \mathbf{8}_{v} &\to \tfrac{3}{2}^{4},\,\tfrac{1}{2}^{28},\,-\tfrac{1}{2}^{28},\,-\tfrac{3}{2}^{4} \:. \label{11.6.22b}
	\end{align}
\end{subequations}
\begin{tcolorbox}[breakable]
	\begin{remark}
		Under $SO(8)\supset SO(2)\times SO(6)$ the $SO(2)$ charge is the 4d helicity, i.e.\ the first weight component $s_{1}$. The vector weights are $\pm e_{1}$ (helicity $\pm1$) and $\pm e_{2,3,4}$ (helicity 0, six states), which is (\ref{11.6.21a}). The spinor weights $(\pm\frac{1}{2})^{4}$ with an even number of minus signs have $s_{1}=+\frac{1}{2}$ for four of them and $-\frac{1}{2}$ for the other four, which is (\ref{11.6.21b}). Helicities add in tensor products, so the multiplicities follow by convolution:
		\begin{align*}
			\mathbf{8}_{v}\times\mathbf{8}_{v}:\quad \pm2&:1\cdot1=1\:,\qquad \pm1:\ 1\cdot6+6\cdot1=12\:,\qquad 0:\ 1\cdot1+1\cdot1+6\cdot6=38\:,\\
			\mathbf{8}\times\mathbf{8}_{v}:\quad \pm\tfrac{3}{2}&:4\cdot1=4\:,\qquad \pm\tfrac{1}{2}:\ 4\cdot6+4\cdot1=28\:,
		\end{align*}
		with totals $1+12+38+12+1=64$ and $2(4+28)=64$. The helicity-zero states from $\mathbf{8}_{v}\times\mathbf{8}_{v}$ are the $6\times6=36$ internal components $G_{mn}$ ($21$) and $B_{mn}$ ($15$), plus the dilaton and the axion dual to $B_{\mu\nu}$, giving $36+2=38$. The helicity-$\pm1$ states are the $6+6$ vectors $G_{\mu m}$ and $B_{\mu m}$. The $\pm\frac{3}{2}$ states are four gravitini, one for each of the $4$ spinors of $SO(6)\cong SU(4)$, which is the $N=4$ R-symmetry.
	\end{remark}
\end{tcolorbox}
In this supergravity multiplet, there is one graviton with helicity $\pm2$ and 4 gravitinos with helicity $\pm\frac{3}{2}$. Toroidal compactification does not break any supersymmetry. Since supercharges have 4 components in 4 dimensions, 16 supersymmetries reduce to $d=4, N=4$ supersymmetry. For this compactification, this supergravity multiplet also introduces 12 Kaluza-Klein and antisymmetric tensor gauge bosons, some fermions, and 36 moduli. The last two states with spin zero are the dilaton and the axion. In 4 dimensions, the second-rank tensor $B_{\mu\nu}$ is equivalent to a scalar. This is the axion, the physics of which will be discussed further in Chapter 18.

In ten dimensions, the only fields carrying gauge charges are the gauge fields and gauginos. They reduce to an $N=4$ vector multiplet—a gauge field, 4 Weyl spinors, and 6 scalars, all in the adjoint representation. For expanded gauge symmetries not manifested in ten dimensions, what is still obtained due to supersymmetry is the same $N=4$ vector multiplet. In compactification with $N=4$ supersymmetry, because fermions must be in the adjoint representation of the gauge group, it cannot yield the Standard Model. We will explain in detail in section 16.2 that one gravitino is good, but 4 gravitinos are excessive. We will see in Chapter 16 that a relatively simple orbifold compactification reduces this symmetry to $N=1$ and gives a relatively realistic spectrum.

\subsection*{Supersymmetry and BPS states}

Thinking for a moment shows that the supersymmetry algebra of toroidally compactified theories must take the following form
\begin{equation}
	\{Q_{\alpha},Q^{\dag}_{\beta}\} = 2P_{\mu}(\Gamma^{\mu}\Gamma^{0})_{\alpha\beta} +2P_{Rm}(\Gamma^{m}\Gamma^{0})_{\alpha\beta} \:. \label{11.6.23}
\end{equation}
It differs from the simple dimensional reduction of the ten-dimensional algebra in that we replaced $P_{m}$ with $P_{Rm}$, namely the total right-moving momentum $k_{Rm}$ of all strings in a given state. They are equal only when the total winding number is zero. This algebra must take this form because spacetime supersymmetry only involves the right-moving part of the heterotic string.

We now look for \emph{Bogomolnyi-Prasad-Sommerfield} (BPS) states, namely states annihilated by certain $Q_{\alpha}$. For a single string of mass $M$, take the vacuum expectation value of algebra (\ref{11.6.23}) on any state $\lvert \psi\rangle$ in the rest frame. The left side is a non-negative matrix. The right side is
\begin{equation}
	2(M+k_{Rm}\Gamma^{m}\Gamma^{0})_{\alpha\beta} \:. \label{11.6.24}
\end{equation}
The zero eigenvectors of this matrix are the supersymmetries that annihilate $\lvert\psi\rangle$. Since $(k_{Rm} \Gamma^{m} \Gamma^{0})^{2} = k_{R}^{2}$, the eigenvalues of the matrix are
\begin{equation}
	2(M\pm\lvert k_{R}\rvert) \:, \label{11.6.25}
\end{equation}
with each of the two signs accounting for half. Thus BPS states have $M^{2}=k_{R}^{2}$.
\begin{tcolorbox}[breakable]
	\begin{remark}
		In the rest frame $P_{0}=-M$ and $P_{i}=0$ for the noncompact spatial directions. With $(\Gamma^{0})^{2}=-1$,
		\begin{equation*}
			2P_{\mu}\Gamma^{\mu}\Gamma^{0}=2(-M)(\Gamma^{0})^{2}=2M\:,
		\end{equation*}
		which gives (\ref{11.6.24}). Let $K=k_{Rm}\Gamma^{m}\Gamma^{0}$. Then
		\begin{align*}
			K^{2}&=k_{Rm}k_{Rn}\,\Gamma^{m}\Gamma^{0}\Gamma^{n}\Gamma^{0}\eqs{1}-k_{Rm}k_{Rn}\,\Gamma^{m}\Gamma^{n}(\Gamma^{0})^{2}
			=k_{Rm}k_{Rn}\,\Gamma^{m}\Gamma^{n}\eqs{2}k_{R}^{2}\:,\\
			\operatorname{Tr}K&=k_{Rm}\operatorname{Tr}(\Gamma^{m}\Gamma^{0})\eqs{3}0\:.
		\end{align*}
		At (1) $\Gamma^{0}$ was moved past $\Gamma^{n}$ ($m,n$ spatial). At (2) only the symmetric part $\frac{1}{2}\{\Gamma^{m},\Gamma^{n}\}=\delta^{mn}$ survives the symmetric sum. At (3) $\operatorname{Tr}(\Gamma^{m}\Gamma^{0})=\frac{1}{2}\operatorname{Tr}\{\Gamma^{m},\Gamma^{0}\}=0$. $K$ is also Hermitian, since $(\Gamma^{m}\Gamma^{0})^{\dagger}=\Gamma^{0\dagger}\Gamma^{m}=-\Gamma^{0}\Gamma^{m}=\Gamma^{m}\Gamma^{0}$. So its eigenvalues are $\pm\lvert k_{R}\rvert$, in equal numbers since it is traceless, and those of $2(M+K)$ are $2(M\pm\lvert k_{R}\rvert)$. Positivity of $\{Q,Q^{\dagger}\}$ gives $M\geq\lvert k_{R}\rvert$. A zero eigenvalue, i.e.\ a supercharge annihilating $\lvert\psi\rangle$, requires $M=\lvert k_{R}\rvert$, and then exactly half the supercharges annihilate the state.
	\end{remark}
\end{tcolorbox} Recalling the mass-shell condition for the right-moving side of the heterotic string,
\begin{equation}
	M^{2} = \begin{cases}
		k_{R}^{2}+4(\tilde{N}-\tfrac{1}{2})/\alpha^{\prime} & \quad(\text{NS})\:, \\
		k_{R}^{2}+4\tilde{N}/\alpha^{\prime} &\quad (\text{R})\:,
	\end{cases} \label{11.6.26}
\end{equation}
BPS states are those whose right-moving side is in the R ground state or in an NS state excited once by $\psi_{-1/2}$. The latter is the lowest state remaining under GSO projection, so it makes sense to change terminology and call them ground states here. Then BPS states are those whose right-moving side resides in the $\mathbf{8}_{v}+\mathbf{8}$ ground state, but with arbitrarily large $k_{R}$. They can match various possible left-moving states. The left-moving mass-shell condition is
\begin{equation}
	M^{2}=k_{L}^{2}+4(N-1)/\alpha^{\prime} \label{11.6.27}
\end{equation}
or
\begin{equation}
	N=1+\alpha^{\prime}(k_{R}^{2}-k_{L}^{2})/4=1-n_{m}w^{m}-q^{I}q^{I}/2 \:. \label{11.6.28}
\end{equation}
\begin{tcolorbox}[breakable]
	\begin{remark}
		For a BPS state the right side is in its ground state, $\tilde{N}=0$ (R) or $\tilde{N}=\frac{1}{2}$ (NS), so (\ref{11.6.26}) gives $M^{2}=k_{R}^{2}$. Equating with (\ref{11.6.27}),
		\begin{equation*}
			k_{R}^{2}=k_{L}^{2}+\frac{4(N-1)}{\alpha'}\quad\Longrightarrow\quad N=1+\frac{\alpha'}{4}\bigl(k_{R}^{2}-k_{L}^{2}\bigr)\eqs{1}1-\frac{1}{2}\cdot\frac{\alpha'}{2}\,k\circ k\eqs{2}1-n_{m}w^{m}-\frac{1}{2}q^{I}q^{I}\:.
		\end{equation*}
		At (1) $k_{L}^{2}$ includes the 16 internal components $k_{L}^{I}$, so $k_{L}^{2}-k_{R}^{2}=k\circ k$. At (2) the lattice identity $\frac{\alpha'}{2}k\circ k=2n_{m}w^{m}+q^{2}$ from the box following (\ref{11.6.17a})--(\ref{11.6.17c}) was used. The condition is independent of the moduli $G$, $B$, $A$. Since $N\geq0$, BPS states need $n_{m}w^{m}+\frac{1}{2}q^{2}\leq1$. For $N=0$ these are the states with $n\cdot w+\frac{1}{2}q^{2}=1$.
	\end{remark}
\end{tcolorbox}
Any left-moving oscillator states are possible as long as compactified momentum and winding number satisfy condition (\ref{11.6.28}). For any given left-moving state, the 16 left-moving states $\mathbf{8}_{v}+\mathbf{8}$ form an extremely short multiplet of the supersymmetry algebra, as opposed to the 256 states in a normal massive multiplet.

Now let's look at the ten-dimensional origin of the modified supersymmetry algebra (\ref{11.6.23}). Rewrite this algebra as
\begin{equation}
	\{Q_{\alpha},Q^{\dag}_{\beta}\}=2P_{M}(\Gamma^{M}\Gamma^{0})_{\alpha\beta}-2\frac{\Delta X_{m}}{2\pi\alpha^{\prime}}(\Gamma^{m}\Gamma^{0})_{\alpha\beta}\:,\label{11.6.29}
\end{equation}
where $\Delta X^{m}$ is the total winding number of the string. Examine the limit where the compactification radius becomes macroscopically visible, which makes a wound string also macroscopically visible. The central charge term in the supersymmetry algebra must be proportional to a conserved charge, so we seek a charge proportional to the string length $\Delta X$. In fact, the coupling of the string with the antisymmetric tensor is
\begin{subequations}
	\begin{align}
		\frac{1}{2\pi\alpha^{\prime}}\int_{M} B &=\frac{1}{2}\int \dif^{10}x\:j^{MN}(x)B_{MN}(x) \label{11.6.30a} \\
		j^{MN}(x) &= \frac{1}{2\pi\alpha^{\prime}} \int_{M}\dif^{2}\sigma\: (\partial_{1}X^{M}\partial_{2}X^{N}-\partial_{1}X^{N}\partial_{2}X^{M})\delta^{10}(x-X(\sigma))\:.\label{11.6.30b}
	\end{align}
\end{subequations}
Integrating the current at a fixed time gives the charge
\begin{equation}
	Q^{M} = \int\dif^{9}x\:j^{M0}=\frac{1}{2\pi\alpha^{\prime}}\int\dif X^{M} \:, \label{11.6.31}
\end{equation}
\begin{tcolorbox}[breakable]
	\begin{remark}
		\textbf{(\ref{11.6.30a})--(\ref{11.6.30b}).} The pullback of $B$ to the world-sheet, rewritten with a delta function, is
		\begin{align*}
			\frac{1}{2\pi\alpha'}\int_{M}B&=\frac{1}{2\pi\alpha'}\int\dif^{2}\sigma\,B_{MN}(X)\,\partial_{1}X^{M}\partial_{2}X^{N}\\
			&\eqs{1}\frac{1}{2\pi\alpha'}\int\dif^{2}\sigma\,\frac{1}{2}B_{MN}(X)\bigl(\partial_{1}X^{M}\partial_{2}X^{N}-\partial_{1}X^{N}\partial_{2}X^{M}\bigr)\\
			&\eqs{2}\frac{1}{2}\int\dif^{10}x\,B_{MN}(x)\,j^{MN}(x)\:.
		\end{align*}
		At (1) $B_{MN}$ is antisymmetric. At (2) $1=\int\dif^{10}x\,\delta^{10}(x-X(\sigma))$ was inserted and the $\sigma$ integral moved into $j^{MN}$.

		\textbf{(\ref{11.6.31}).} In static gauge $X^{0}=\sigma^{2}$, so $\partial_{1}X^{0}=0$ and $\partial_{2}X^{0}=1$:
		\begin{align*}
			\int\dif^{9}x\,j^{M0}&=\frac{1}{2\pi\alpha'}\int\dif^{2}\sigma\,\partial_{1}X^{M}\,\partial_{2}X^{0}\int\dif^{9}x\,\delta^{10}(x-X(\sigma))\\
			&\eqs{3}\frac{1}{2\pi\alpha'}\int\dif\sigma^{1}\dif\sigma^{2}\,\partial_{1}X^{M}\,\delta(x^{0}-\sigma^{2})
			=\frac{1}{2\pi\alpha'}\int\dif\sigma^{1}\,\partial_{1}X^{M}\:.
		\end{align*}
		At (3) the nine spatial delta functions were integrated, leaving $\delta(x^{0}-X^{0})$. The result is $\frac{1}{2\pi\alpha'}\int\dif X^{M}$ along the string at time $x^{0}$.

		\textbf{Right-moving combination.} For a string wound $w$ times on a circle of radius $R$, $\int\dif X^{m}=2\pi Rw$, so $Q_{m}=wR/\alpha'$ and, at $B=A=0$, $G=1$,
		\begin{equation*}
			P_{m}-Q_{m}=\frac{n}{R}-\frac{wR}{\alpha'}=k_{Rm}\:,
		\end{equation*}
		by (\ref{11.6.17c}). So (\ref{11.6.32}) reduces to (\ref{11.6.23}) in the compact directions, and $\Delta X_{m}/2\pi\alpha'=Q_{m}$ is the term in (\ref{11.6.29}).
	\end{remark}
\end{tcolorbox}
the integral along the worldline of the string. The entire supersymmetry algebra is
\begin{equation}
	\{Q_{\alpha},Q^{\dag}_{\beta}\}=2(P_{M}-Q_{M})(\Gamma^{M}\Gamma^{0})_{\alpha\beta}\:. \label{11.6.32}
\end{equation}

In ten-dimensional non-compact spacetime, charge (\ref{11.6.31}) is zero for a finite closed string but non-zero for an infinitely long string; for example, an infinitely long string can serve as a limit of a wound string as $R\to\infty$. It is usually useful to consider such a macroscopic string, which clearly has infinite total mass and total charge, but the mean value per unit length is finite. In compactification, the combination $P_{m}-Q_{m}$ is the right-moving gauge charge. Left-moving charges do not appear in the supersymmetry algebra.

A natural question is whether algebra (\ref{11.6.32}) is now complete; in fact, it is not. Examine compactification to four-dimensional spacetime at a general point in moduli space, where the gauge group breaks to $U(1)^{28}$. In grand unified theories where the $U(1)$ of the Standard Model is embedded in a simple group, there are always magnetic monopoles generated from the quantization of topologically non-trivial classical solutions. String theory is not an ordinary grand unified theory, but it also has magnetic monopoles. Compactification of heterotic strings gives three kinds of gauge symmetries: ten-dimensional symmetry, Kaluza-Klein symmetry, and antisymmetric tensor symmetry. For each symmetry, there is a corresponding monopole solution: the 't Hooft-Polyakov monopole solution, the Kaluza-Klein monopole, and the $H$-monopole. Of course, since various charges can be exchanged by $O(22,6,\mathds{Z}) T$-duality, monopoles must be similarly so. Monopole charges appear in the supersymmetry algebra; in the current case, they still appear in the right-moving charges. In low-energy supergravity theory, there is a symmetry that exchanges electric and magnetic charges, so they must appear in the supersymmetry algebra in a symmetric way.